% concatenated from RMS_talk.tex
\documentclass[a4paper,11pt]{amsart}
\usepackage{wrapfig}
\usepackage{ragged2e}
\usepackage{graphicx} % Required for inserting images
\usepackage{amsmath}
\usepackage{amssymb}
\usepackage{amsthm}
\usepackage{hyperref}
\usepackage{enumerate}
\usepackage{enumitem}
\usepackage[normalem]{ulem}
\usepackage{mathdots}
\usepackage{tikz}
\usepackage{tikz-cd}
\usepackage[]{mdframed}
\usepackage[dvipsnames]{xcolor}
\usepackage{eqnarray}
%\usepackage{ntheorem}
\usepackage{xy}
\input xy
\xyoption{all}
\usepackage{multirow}
\usepackage{leftindex}
\usepackage{array}
\usepackage{float}
\usepackage{booktabs}
\usepackage{caption}
\usepackage{geometry}
\geometry{margin=1.2in}
\numberwithin{equation}{section}
\numberwithin{equation}{subsection}
\newtheorem{theorem}{Theorem}[section]
\newtheorem{proposition}[theorem]{Proposition}
\newtheorem{conjecture}[theorem]{Conjecture}
\newtheorem{corollary}[theorem]{Corollary}
\newtheorem{lemma}[theorem]{Lemma}
\newtheorem{claim}[theorem]{Claim}
\newtheorem{defn}[theorem]{Definition}
\newtheorem*{remark*}{Remark}
\newtheorem{definition}[theorem]{Definition}
\newtheorem{remark}[theorem]{Remark}

\newtheorem{example}[theorem]{Example}
\newtheorem{problem}[theorem]{Problem}
\newtheorem{algorithm}[theorem]{Algorithm}
\newtheorem{question}[theorem]{Question}
\newtheorem{solution}[theorem]{Solution}

\newcommand{\Q}{\mathbb{Q}}
\newcommand{\R}{\mathbb{R}}
\newcommand{\C}{\mathbb{C}}
\newcommand{\F}{\mathbb{F}}
\newcommand{\Ca}{\mathbb{C^\alpha}G}
\newcommand{\Cb}{\mathbb{C^\beta}G}
\newcommand{\Z}{\mathbb{Z}}
\newcommand{\N}{\mathbb{N}}
\newcommand{\n}{\mathbb{R}^n}
\newcommand{\m}{\mathbb{R}^m}
\newcommand{\dg}{\delta(G)}
\newcommand{\di}{\delta_{irr}(G)}
\newcommand{\tg}{\tau(G)}
\newcommand{\ti}{\tau_{irr}(G)}
\newcommand{\irr}{\mathrm{Irr}}
\newcommand{\Hom}{\mathrm{Hom}}
\newcommand{\PGL}{\mathrm{PGL}}
\newcommand{\bchi}{\Bar{\chi}}
\newcommand{\Ho}{{\mathrm{H}}}
\newcommand{\x}{\Bar{x}}
\newcommand{\y}{\Bar{y}}
\newcommand{\z}{\Bar{z}}
\newcommand{\tra}{\operatorname{tra}}
\newcommand{\im}{\operatorname{Im}}
\newcommand\myeq{\stackrel{\mathclap{\normalfont\mbox{s}}}{~}}
\newcommand{\GL}{\mathrm{GL}}
\newcommand{\U}{Um_n(A)}
\newcommand{\Nayak}[1]{{\color{red} #1}}
\newcommand{\poonam}[1]{{\color{blue} #1}}
\usepackage[utf8]{inputenc}


\title[On the Faithful Projective Representations of Finite Groups]
{On the Faithful Projective Representations of Finite Groups and their Minimal Dimension}

\author{Sumana Hatui} \author{Poonam Nayak}
\address{School of Mathematical Sciences,  National Institute of Science Education and Research Bhubaneswar, 752050, Odisha, India.}
\address{Homi Bhabha National Institute, Training School Complex, Anushakti Nagar, Mumbai, 400094, India.
}
\email{sumanahatui@niser.ac.in}
\email{poonam.nayak@niser.ac.in, rashmitanayak2021@gmail.com}

\keywords{Faithful projective representations, Projective embedding degree, Representation groups, Second cohomology groups}

\subjclass[2020]{20C25, 20D15, 20J06}

\begin{document}
	
	\begin{abstract}
		The first part of this article is devoted to characterizing the cocycles $\alpha$ of a finite group $G$ that give rise to faithful projective representations of $G$. 
		We prove that a $p$-group $G$ admits a faithful irreducible projective representation if and only if the cohomology class $[\alpha]$ does not lie in the image of the inflation map
		$\inf: \mathrm{H}^2\!\left(G / N, \mathbb{C}^{\times}\right) \longrightarrow \mathrm{H}^2\!\left(G, \mathbb{C}^{\times}\right)$
		for any non-trivial central subgroup $N$ of $G$.
		In the case where $[\alpha] \in \operatorname{Im}(\operatorname{inf})$, we determine a criterion such that a direct sum of irreducible $\alpha$-representations is faithful. We conclude this part by describing the behaviour of cocycles $\alpha$ that yield faithful irreducible representations for direct products of groups.
		
		In the second part, we introduce the notion of the projective embedding degree of a finite group $G$, defined as the smallest integer $n$ such that $G$ embeds into $\PGL_n(\C)$; equivalently, it is the smallest $n$ such that $G$ has a faithful complex projective representation of degree $n$. We also define the analogous notion of the irreducible projective embedding degree of $G$. These invariants have been investigated for several classes of groups, including direct products of groups, finite abelian groups, extra-special $p$-groups, Heisenberg groups, and groups of order $p^3$, $p^4$ (for primes $p$), and $p^5$ (for $p \geq 5$). 
		
		% we introduce the notion of the \emph{projective embedding degree} of a finite group $G$, defined as the smallest integer $n$ such that $G$ embeds into $\PGL_n(\C)$; equivalently, it is the smallest $n$ such that $G$ has a faithful complex projective representation of degree $n$. We also define the \emph{irreducible projective embedding degree} of $G$ as the smallest $n$ for which $G$ has a faithful irreducible complex projective representation of degree $n$. These invariants have been investigated for several classes of groups, including direct products of groups, finite abelian groups, extra-special $p$-groups, Heisenberg groups, and groups of order $p^3$, $p^4$ (for primes $p$), and $p^5$ (for $p\geq 5$). 
		
		%Next, we characterize those cocycles $\alpha$ that yield faithful irreducible $\alpha$-representations for direct products of groups. We conclude this part by determining a criterion when a direct sum of irreducible $\alpha$-representations gives rise to a faithful $\ alpha$-representation.
		
	\end{abstract}
	
	\maketitle
	
	%\tableofcontents
	
	\section{Introduction}
	%Projective representations arise naturally in several areas of mathematics, including group cohomology, number theory, and mathematical physics. 
	A complex projective representation of a finite group $G$ is a homomorphism $\rho \colon G \to \PGL(V)$, where $V$ is a finite-dimensional complex vector space. Equivalently, we say a map $\rho: G \to GL(V)$ is a projective representation of $G$ if there is a map $\alpha \colon G \times G \to \mathbb{C}^\times$ such that $$\rho(g)\rho(h)= \alpha(g,h)\rho(gh), g, h \in G \text{ and } \rho(1)=1_V.$$
	Then $\alpha$ becomes a cocycle of $G$, and we refer $\rho$ as an $\alpha$-representation.
	Projective representations are a generalization of ordinary representations. The equivalence classes of projective representations of $G$ are classified by the elements of the second
	cohomology group $\mathrm{H}^2(G,\mathbb{C}^{\times})$, known as the Schur multiplier of $G$. It is also denoted by $M(G)$.
	
	A fundamental and long-standing problem in the representation theory of finite groups $G$ is to determine the smallest positive integer $n$ for which $G$ embeds into $\mathrm{GL}_n(\mathbb{C})$.
	%This problem naturally motivates the study of faithful ordinary representations of finite groups $G$. 
	To address this problem, the embedding degree $\delta(G)$ and irreducible embedding degree $\delta_{irr}(G)$ of a group $G$ were introduced, which are defined as follows. The embedding degree (or the representation dimension), $\delta(G)$, is the minimal dimension of a faithful complex ordinary representation of $G$, and the irreducible embedding degree $\delta_{irr}(G)$ is the minimal dimension of a faithful irreducible complex ordinary representation of $G$. 
	%, i.e., $G$ is isomorphic to a complex linear group of degree $n$. 
	%The \textit{embedding degree}, also called the representation dimension of $G$, denoted by $\delta(G)$, is the minimum dimension of a faithful complex ordinary representation of $G$. 
	The representation dimension has been studied extensively for finite $p$-groups (see \cite{bardestani, cernele, kaur}) as well as for finite groups of Lie type (\cite{Lusztig, Lübeck}).
	While questions of faithfulness, irreducibility, and minimal dimension for ordinary representations of finite groups have been extensively studied in \cite{ Martino, moreto, Prajapati, SZECHTMAN}, analogous questions for projective representations have received comparatively little attention.
	%—such as faithfulness, irreducibility, and minimal dimension—
	Existing work on projective representations has
	primarily focused on faithfulness and irreducibility \cite{Bekka, frucht, HIGGS, Morris}, with further results obtained for specific families of groups, such as metabelian groups \cite{NG}, and finite classical groups in \cite{LANDAZURI, Zalesskii}. %REMOVED YAMAZAKI from citations after irreducibility.
	
	These earlier studies naturally lead to the problem of determining the smallest positive integer $n$ for which a finite group $G$ embeds into $\mathrm{PGL}_n(\mathbb{C})$; equivalently, of deciding when $G$ admits a faithful (possibly irreducible) projective representation, and what the minimal possible dimension of such a representation is. 
	%In this paper, we investigate faithful complex projective representations of finite groups, placing special emphasis on $p$-groups and determining their minimal dimension. \Nayak{Check it}.
	%as the corresponding questions for finite nilpotent groups reduce to the case of $p$-groups. \poonam{Only in case of irreducible} 
	Let $Z^2(G,\mathbb C^\times)$ denote the set of cocycles of $G$ over $\mathbb C^\times$.
	Recall that if $N \lhd G$, the inflation map is defined as follows. For a given $\alpha \in Z^2(G/N, \mathbb{C}^{\times})$, define $\alpha' : G \times G \to \mathbb{C}^{\times}$ by $\alpha'(x, y) = \alpha(xN, yN)$. Then the assignment $[\alpha] \mapsto [\alpha']$ induces a homomorphism $\text{inf} : \text{H}^2(G/N, \mathbb{C}^{\times}) \to \text{H}^2(G, \mathbb{C}^{\times})$, called the inflation map. 
	
	Our first result appears in  Theorem \ref{cohomological characterization}, where we show that a finite $p$-group $G$ admits a faithful irreducible $\alpha$-representation if and only if the cohomology class
	$[\alpha]$ does not lie in the image of any inflation homomorphism
	\[
	\mathrm{inf} \colon
	\mathrm{H}^2\!\left(G/N, \mathbb{C}^{\times}\right)
	\longrightarrow
	\mathrm{H}^2\!\left(G, \mathbb{C}^{\times}\right),
	\]
	where $N$ ranges over all non-trivial central subgroups of $G$. As a corollary, we obtain corresponding results for non-capable groups (see Corollary \ref{non-capable groups}). Next, in Theorem \ref{direct_sum}, we determine a criterion when a direct sum of irreducible $\alpha$-representations gives rise to a faithful $\alpha$-representation for $[\alpha] \in \operatorname{Im}(\operatorname{inf})$. Furthermore, for direct products of groups $G$ with $Z(G) \subseteq [G, G]$, we provide an explicit description of the cocycles $\alpha$ that give rise to faithful irreducible $\alpha$-representations in Theorem \ref{directproduct}. 
	
	Since every finite nilpotent group is isomorphic to the direct product of its Sylow $p$-subgroups, it follows from Theorem~\ref{schur multiplier of direct product} and Theorem~\ref{tensor product is faithful} that the study of faithful irreducible projective representations of finite nilpotent groups reduces to the corresponding problem for $p$-groups. Therefore, we restrict our attention to $p$-groups.
	
	
	For an $\alpha$-representation  $\rho : G \to \mathrm{GL}(V)$ of $G$, we define the kernel of $\rho$ by $\ker \rho=\{ g \in G \mid \rho(g) = \lambda I_V$ for some $\lambda\in \C^{\times}$\}. If $\ker \rho = 1$, then we say that $\rho$ is a faithful projective representation. Note that every finite group admits a faithful complex projective representation, which follows
	from the fact that the $\alpha$-regular representation of $G$ is always faithful (see \cite[Chapter 7, Section 1]{karpilovsky}).
	%It follows that every finite group admits a faithful complex projective representation, and hence embeds into $\PGL_n(\mathbb{C})$ for some $n \geq 1$.
	Building on this observation, we introduce numerical invariants that capture the smallest possible dimensions of faithful projective representations. We define the \textit{projective embedding degree} $\tau(G)$ of \(G\) to be the smallest positive integer \(n\) for which \(G\) embeds into \(\PGL_n(\mathbb{C})\).
	Equivalently,
	\[
	\tau(G) := \min\{\deg(\rho) : \rho \text{ is a faithful complex  projective representation of } G\},
	\]
	where $\deg(\rho)$ denotes the degree of  $\rho$.
	The \textit{irreducible projective embedding degree} $\ti$ of  $G$ is defined as follows:
	\[
	\ti := \min\{\deg(\rho) : \rho \text{ is a faithful irreducible complex projective representation of } G\}.
	\] 
	We note that \(\tau_{irr}(G)\) need not exist for every group \(G\). However, whenever it exists, one has $\tau(G) \leq \tau_{irr}(G).$
	In Section~\ref{emdedding degree}, we establish sufficient conditions under which two groups \(G_1\) and \(G_2\) satisfy
	$\tau(G_1)= \tau(G_2)$
	(see Theorem~\ref{tau(G_1) = tau(G_2)}). We also study the invariants \(\tau(G)\) and \(\tau_{irr}(G)\) for several important classes of finite groups, including direct products of groups, finite abelian groups, extra-special \(p\)-groups, and Heisenberg groups.
	In Section~\ref{additional}, we present several examples of groups \(G\) and compute \(\tau(G)\) and \(\tau_{irr}(G)\). We also carry out these computations for non-abelian groups of orders \(p^3\) and \(p^4\), where \(p\) is a prime, as well as for groups of order \(p^5\) with \(p \geq 5\); see Section~\ref{groupsoforderp5}. The determination of these invariants is more intricate and requires a number of non-trivial arguments. The proofs make essential use of auxiliary lemmas and techniques developed in Sections~\ref{sec:Preliminaries}, \ref{faithful_projective}, \ref{emdedding degree}, and \ref{groupsoforderp5}.
	
	
	% We note that $\ti$ does not exist for all groups $G$. However, if it exists, then $\tau(G) \leq \ti$.In Section \ref{emdedding degree}, we determine sufficient conditions for two groups $G_1$ and $G_2$ to satisfy $\tau(G_1) = \tau(G_2)$ (see Theorem \ref{tau(G_1) = tau(G_2)}), and we study $\tau(G)$ and $\tau_{irr}(G)$ for several significant families of finite groups, including direct product of groups, finite abelian groups, extra-special \(p\)-groups and Heisenberg groups. 
	
	%In Section \ref{additional}, we present several examples of groups $G$ and compute $\tau(G)$ and $\tau_{irr}(G)$. Next, we compute these invariants for non-abelian groups $G$ of order \(p^3\), \(p^4\) for primes $p$, and for groups of order \(p^5\) for primes \(p \geq 5\) in Section \ref{groupsoforderp5}. The determination of these invariants is not straightforward and requires several non-trivial arguments. The proofs make use of a number of auxiliary lemmas and techniques established in Sections \ref{sec:Preliminaries}, \ref{faithful_projective}, \ref{emdedding degree}, and \ref{groupsoforderp5}. %The proofs concerning the invariants for groups of order $p^n$ are not straightforward and require several novel techniques. For clarity and independent interest, these methods are presented separately as lemmas in Section \ref{groupsoforderp5}. \poonam{You had said that in the introduction, we should add that we employ several techniques for computing $\tau(G)$ and $\tau_{irr}(G)$ for groups of order $p^5$, so I have added this at the end of the paragraph.}
	
	%Further, in Lemma  \ref{direct sum of k irreducible ordinary representations}, we derive the minimum number of elements in $\mathrm{Irr}(G)$ such that their direct sum yields a faithful projective representation of $G$, where 
	
	\begin{remark}
		One might expect that there is not much significant difference between $\tau(G)$ and $\delta(G)$; however, this is not the case. Our computations show that there is a large class of groups $G$ for which $\tau(G)$ and $\delta(G)$ differ substantially, for example, extra-special $p$-groups, Heisenberg groups, and groups of order $p^5$ in $\Phi_7$. (See Section \ref{EH}, Theorem \ref{p5} and \cite[Theorem 3.1]{kaur}.) 
	\end{remark} 
	%Let $G$ be a finite group. Recall that the \textit{embedding degree}, or the representation dimension, of a finite group $G$, denoted by $\delta(G)$, is the minimal dimension of a faithful complex ordinary representation of $G$. Likewise, $\delta_{irr}(G)$ is the minimal dimension of a faithful irreducible complex ordinary representation of $G$. \\
	
	%The proof of this result is given in Section \ref{abelian}.
	
	%From the above result, the question of determining $\tau_{irr}(G)$ boils down to computing  
	%A proof of this result is given in Section.
	
	%In this article, our goal is to compute the projective embedding degree and the projective irreducible embedding degree for
	
	%Finally, we compute $\tau(G)$ and $\tau_{irr}(G)$ explicitly for non-abelian groups $G$ of order $p^4$ and $p^5$ for primes $p \ge 5$ (see Section \ref{groupsoforderp5}). For isoclinic groups, the equality $\delta(G) = \delta_{irr}(G)$ holds (\cite{kaur}, Theorem 3.1). However, this is not true in the present case, as is evident from the computation of $\tau(G)$ and $\tau_{irr}(G)$.
	
	Let us fix some notation here. For $x,y\in G$, the commutator $x^{-1}y^{-1}xy$ is denoted by $[x,y]$. The subgroups $G'$ and $Z(G)$ refer to the commutator subgroup and the centre of $G$, respectively. For $x \in G$, $C_G(x)$ denotes the centralizer of $x$ in $G$. The one-dimensional trivial ordinary representation of $G$ is denoted by $1_G$. We write $cd(G)= \{\deg(\rho) \mid \rho \in \irr(G) \}$, where $\mathrm{Irr}(G)$ denotes the set of all linearly inequivalent irreducible ordinary representations of $G$. If $\psi$ is an ordinary representation of $G$, then we write $K_{\psi}= \{x \in G \mid \psi(x)=I\}$. For $\alpha \in Z^2(G,\mathbb C^\times)$, $\mathrm{Irr}^\alpha(G)$ denotes the set of all irreducible $\alpha$-representations of $G$ up to linear equivalence. Furthermore, we say that a cocycle $\alpha \in Z^{2}(G,\mathbb{C}^{\times})$ is non-trivial if $[\alpha] \neq [1]$.  
	
	Throughout this paper, we restrict our attention to finite groups and complex representations. 
	By the tensor product $G_1 \otimes G_2$ of two groups $G_1$ and $G_2$, we always mean the abelian tensor product, i.e., $G_1/G_1' \otimes_{\mathbb Z} G_2/G_2'$. \\
	
	\begin{mdframed}
		\textbf{Unless stated otherwise, we say $\rho: G \to GL(V)$ is faithful on $G$ if it is faithful as a projective representation of $G$.}
	\end{mdframed}
	
	\section{Preliminaries} \label{sec:Preliminaries}
	In this section, we present several results that will be used in the proofs of subsequent sections. For basic definitions and results on projective representations, we refer to \cite[Chapter 3]{karpilovsky}.
	\begin{definition}[$\alpha$-regular element]
		For a cocycle $\alpha$ of $G$, an element $x \in G$ is called an $\alpha$-regular element of $G$ if $\alpha(g,x) =\alpha(x,g)$ for all $g\in C_G(x)$. Further, we say that $x \in G$ is an $\alpha$-regular element of $Z(G)$ if $x \in Z(G)$ and $x$ is an $\alpha$-regular element of $G$.    
	\end{definition}
	
	\iffalse
	\begin{lemma}(\cite{Serre}, Chapter 8, Theorem 16)
		Let G be a $p$-group. Then each irreducible ordinary representation of G is induced by an ordinary representation of degree 1 of a subgroup of G (i.e., is monomial). 
	\end{lemma}
	\fi
	
	\begin{lemma}[\cite{karpilovsky}, Chapter 3, Lemma 1.1] \label{projectively equivalent to an ordinary representation}
		Suppose $\alpha_i \in Z^2(G,\mathbb C^\times)$, $i=1,2$.
		Let $\rho_1$ be an $\alpha_1$-representation. If $[\alpha_1]= [\alpha_2]$, then there exists an $\alpha_2$-representation $\rho_2$ which is projectively equivalent to $\rho_1$. In particular, if $\alpha_1$ is a coboundary, then $\rho_1$ is projectively equivalent to an ordinary representation.
		
	\end{lemma}
	
	The next result follows immediately.
	\begin{lemma} \label{projectively equivalent} 
		Let $\rho_{1}$ and $\rho_{2}$ be projective representations that are projectively equivalent. If $\rho_{1}$ is faithful, then $\rho_{2}$ is also faithful.
	\end{lemma}
	
	%%Let $\rho_i: G \rightarrow G L\left(V_i\right), i=1,2$, be two projective representations of $G$, where $V_i$ is a complex vector space. Since $\rho_{1}$ and $\rho_{2}$ are projectively equivalent, therefore, there exists a map $\mu:G \rightarrow \C^{\times}$ with $\mu(1)=1$ and a vector space isomorphism $T: V_1 \rightarrow V_2$ such that $\rho_2(g) =\mu(g) T \rho_1(g) T^{-1}$ $ \forall g \in G$. \\
	
	%If $\rho_2(g)=\lambda I_{V_2}$  for some $g\in G$, $\lambda \in \mathbb{C}^{\times}$, then,  $\rho_1(g)=c I_{V_1}$ where $c=(\mu(g))^{-1} \lambda$. Hence, the result follows.
	%\end{proof}
	
	In view of the above results, throughout this paper, we consider projective representations associated with a fixed cocycle representative of each cohomology class. In particular, if $\alpha$ is a coboundary, then we restrict our attention to ordinary representations ($\alpha=1$).
	
	\begin{lemma} \label{one dim direct sum in non-abelian group}
		\begin{enumerate}
			\item If $G$ is a non-abelian group, then the direct sum of one-dimensional ordinary representations can never be a faithful projective representation.
			
			\item A faithful irreducible ordinary representation of $G$ is a faithful irreducible projective representation if and only if $Z(G)=1$.
			
			\item If $Z(G) = 1$, then $\tg \leq \dg $.
			Furthermore, if $Z(G) = 1$ and $M(G)=1$, then $\dg=\tg$  \text{ and }  $\ti = \delta_{irr}(G)$.
		\end{enumerate}
		
	\end{lemma}
	
	\begin{proof}
		(1) Suppose $\rho= \oplus_{k=1}^{n}\rho_{i}$, where $\rho_{i}s$ are one-dimensional ordinary representations of $G$.
		For $x \in G^\prime$, $\rho_{i}(x)=1$ and so
		$\rho(x)= I_{n \times n}$. Hence, $\rho$ cannot be faithful.
		
		(2) Let $\psi$ be a faithful irreducible ordinary representation of $G$. Then, by \cite[Chapter 2, Lemma 2.27, (f)]{isaacs}, $\{g$$|$$\psi(g)=\lambda I$ for some $\lambda\in \mathbb{C}^{\times} \} =Z(G)$.
		
		(3) Let $\psi$ be a faithful ordinary representation of $G$. Suppose $\psi(a)= \lambda I$ for some $\lambda$ $\in$ $\mathbb{C}^{\times}$ and $a \in G$. Then, $\psi(a)\psi(g)= \psi(g) \psi(a)$ $\forall g \in G$, which implies $\psi(a^{-1}g^{-1}ag)=I$. Since $\psi$ is a faithful ordinary representation, $[a,g]=1$ for all $g\in G$. 
		Hence, $a$ $\in$ $Z(G)$ and $\psi$ is a faithful projective representation. Thus, $\tg \leq \dg $. Furthermore, if $Z(G) = 1$ and $M(G)=1$, the result is obtained using $(2)$ and the argument above. 
		
	\end{proof}
	
	\begin{remark} \label{some results on p-groups} Let $G$ be a non-abelian $p$-group.
		The following facts are straightforward using the above result.
		\begin{enumerate}
			\item[(i)] $\tau(G) \geq p$.
			
			\item[(ii)] $\tau(G)=p$ if and only if $\tau_{irr}(G) =p$. In this case, the corresponding faithful representation must be a true projective representation.
			
			\item[(iii)] If $\tau_{irr}(G)$ does not exist or $\tau_{irr}(G)>p$, then $\tau(G) > p$.
		\end{enumerate}
	\end{remark}
	
	%This implies that $\phi($  a$\in$ Z(G). Then,$\{g$$|$$\phi(g)=\lambda I$ for some $\lambda$ $\in$ $\mathbb{C}^{\times} \}= Z(G)$.
	
	%Since  every one-dimensional projective representation of $G$ is projectively equivalent to an  ordinary representation, by Theorem 1.11, WLOG, we can assume that $\rho_{i}$, $i=1,\ldots,n$ are ordinary representations.\\
	
	%Before stating our next result, we define some notation. 
	
	Let $N$ be a normal subgroup of $G$ and  $\rho \in \operatorname{Irr}(N)$. The induced representation of $\rho$ from $N$ to $G$ is denoted by $\mathrm{Ind}_N^G(\rho)$. We say $\rho \in \mathrm{Irr}(N)$  extends to $G$ if there exists $\tilde{\rho} \in \mathrm{Irr}(G)$ such that $\tilde{\rho}|_N = \rho$.
	Let $\chi_\rho$ be the character affording $\rho$. For $g \in G$, define conjugate character $\chi_\rho^{(g)}: N \to \C$ by $\chi_\rho^{(g)}(x) = \chi_\rho(gxg^{-1})$ $\forall$ $x \in N$. Then, 
	\[
	I_G(\chi_\rho) = \{ g \in G \mid \chi_\rho^{(g)} = \chi_\rho \}
	\]
	is called the inertia group of $\rho$ in $G$. For any character $\chi$ of $N$, we define  $$\operatorname{Irr}(G \mid \chi)= \{\eta \in \irr(G) \mid \langle \chi_{\eta\mid_N}, \chi \rangle \neq 0\}$$ denote the set of inequivalent irreducible ordinary representations of $G$ lying above $\chi$. 
	Next, we recall a few results of Clifford's theory.
	\begin{theorem}[\cite{isaacs}, Theorem 6.11, Corollary 11.22] Let $N \triangleleft G$ and $\rho \in \mathrm{Irr}(N)$. If $\chi$ is the character affording $\rho$, then the following hold.
		
		\begin{enumerate} \label{Clifford's theory}
			
			%\item  Let $\nu$ be an irreducible constituent of $\chi|_N$. Suppose $\nu=\nu_1$, $\nu_2$, ..., $\nu_t$ are the distinct conjugates of $\nu$ in $G$. Then, there exists an integer $e \geq 1$ such that $\chi|_N = e(\sum_{i=1}^t \nu_i)$. \poonam{Am I using it anywhere?}
			\item The map 
			\[
			\theta \mapsto \operatorname{Ind}^{G}_{I_G(\chi)}(\theta)
			\] is a bijection of $\operatorname{Irr}(I_G(\chi) \mid \chi)$ onto $\operatorname{Irr}(G \mid \chi)$.
			
			\item Suppose $G/N$ is cyclic and $I_G(\chi) = G$. Then $\rho$ extends  to $G$.
			
		\end{enumerate}    
	\end{theorem}
	
	Next, we recall the following definition from \cite{gonzalo}. %\poonam{They have given separate definitions 2.6 and 2.11.}
	\begin{definition} \label{covering group}
		Let $G$ be a group. A group $A.G$ is called a covering of $G$ if there exists a central extension
		\[
		1 \to A \to A.G \to G \to 1
		\]
		such that $A \subseteq M(G)$ and $A \subseteq (A.G)'$. 
		Furthermore, if $A \cong M(G)$, then $A.G$ is called a representation group of $G$, and we denote it by  $G^*$.
	\end{definition} 
	%\poonam{We denote a representation group of $G$ by $G^*$, and this notation will be used throughout without further reference. Can I omit it?} 
	Let $1 \to A \to H \xrightarrow{f} G \to 1$ be a central extension and $\mu$ be a section of $f$, i.e., $\mu: G \to H$ is a map such that $f \circ \mu =id_G$ and $\mu(1)= 1$. 
	For a $\chi \in \mathrm{Hom}(A, \mathbb{C}^{\times})$, define 
	$\alpha_\chi : G \times G \to \mathbb{C}^{\times}$ by $\alpha_\chi(x, y) = \chi(\mu(x)\mu(y)\mu(xy)^{-1})$, $x,y \in G$. Then the assignment $\chi \mapsto [\alpha_\chi]$ induces a homomorphism $\mathrm{tra}: \mathrm{Hom}(A, \mathbb{C}^{\times}) \to \mathrm{H}^2(G, \mathbb{C}^{\times})$, called the transgression map associated with the given central extension. For simplicity, we shall denote the representative cocycle $\alpha_\chi$ by $\mathrm{tra}(\chi)$. 
	
	Now, the result below follows from \cite[Chapter 2, Theorem 2.5]{karpilovsky} and \cite[Proposition 2.1.7]{schurmultiplier}.
	\begin{lemma} \label{exact sequence}
		Let $A \subseteq G' \cap Z(G)$ and $1 \to A \to G \to G/A \to 1$ be the natural exact sequence. 
		Then the induced sequence
		
		$$ 1  \xrightarrow{}\mathrm{Hom}(A,\mathbb{C}^{\times}) 
		\xrightarrow{\mathrm{tra}} 
		\mathrm{H}^2(G/A,\mathbb{C}^{\times}) 
		\xrightarrow{\mathrm{inf}} \mathrm{H}^2(G,\mathbb{C}^{\times})$$
		%\end{aligned}
		is exact.
		
	\end{lemma}
	
	The next result establishes a relationship between the projective representations of a group $G$ and those of a suitable quotient group of $G$,
	which follows from the proof of \cite[Theorem 3.2]{hatui}.
	\begin{theorem} \label{bijective correspondence via inflation}
		Let $A$ be a subgroup of a finitely generated group $G$ such that $A \subseteq G' \cap Z(G)$ and $[\alpha] \in \mathrm { Im } \big(\mathrm{inf}: \mathrm{H}^2(G/A, \C^{\times}) \to \mathrm{H}^2(G, \C^{\times})\big)$. Suppose $\rho : G \to \mathrm{GL}(V) \in \mathrm{Irr}^\alpha(G)$.
		Define $\tilde{\rho} : G/A \to \mathrm{GL}(V)$ by 
		$\tilde{\rho}(gA) = \rho\big(\mu(gA)\big)$, where $\mu : G/A \to G$ is a section of $G \twoheadrightarrow G/A$.
		Then the map
		\[
		\phi : \mathrm{Irr}^\alpha(G) \longrightarrow 
		\bigcup_{\{[\beta] \in \mathrm{H}^2(G/A, \mathbb{C}^{\times}) \mid \mathrm{inf}([\beta]) = [\alpha]\}} 
		\mathrm{Irr}^\beta(G/A)
		\]
		given by $\phi(\rho) = \tilde{\rho}$ is a dimension-preserving bijective map. 
	\end{theorem} 
	
	\begin{remark} \label{remark for bijective correspondence via inflation}
		A closer look at the map $\phi$  gives the following.
		\begin{enumerate}
			\item  If \(\beta \in Z^{2}(G/A,\C^{\times})\) and $\alpha(g,h) = \beta(gA,hA), g,h \in G$, then $\phi$ gives a bijective correspondence between \(\mathrm{Irr}^{\alpha}(G)\) and
			$\bigcup_{\chi \in \mathrm{Hom}(A,\C^{\times})}
			\mathrm{Irr}^{\,\beta\mathrm{tra}(\chi)}(G/A)$.
			
			\item If $\tilde{\rho} \in \operatorname{Irr}^{\beta \mathrm{tra}(\chi)}(G/A)$, then $\rho : G \to \mathrm{GL}_n(\C)$ defined by
			\[
			\rho(a \mu(gA)) = \chi(a)\tilde{\rho}(gA) \quad \text{for all } a \in A, g \in G  
			\] is such that $\phi(\rho) = \tilde{\rho}$ and $\rho |_A=\chi$.
		\end{enumerate}
	\end{remark}
	We now recall the fact that every projective representation of $G$ lifts to an ordinary representation of $G^{*}$, where $G^*$ is a representation group of $G$. Let 
	$$1 \longrightarrow A \longrightarrow G^* \xrightarrow{f} G \longrightarrow 1$$
	be the corresponding central extension and $\mu$ be a section of $f$. Then the associated transgression map
	$
	\operatorname{tra} : \operatorname{Hom}(A, \mathbb{C}^{\times}) \longrightarrow 
	\mathrm{H}^{2}(G, \mathbb{C}^{\times})
	$
	is an isomorphism (see \cite [Theorem 2.1.4]{schurmultiplier}).
	Suppose $\chi \in \mathrm{Hom}(A, \mathbb{C}^{\times})$ and  $\alpha= \mathrm{tra}(\chi)$. For a given $\Gamma \in \operatorname{Irr}(G^* \mid \chi)$, let $\rho : G \to \mathrm{GL}_n(\C)$ be defined by
	\[
	\rho(g) = \Gamma(\mu(g)) \quad \text{for all }g \in G.
	\]
	It is easy to see that $\rho \in \operatorname{Irr}^{\alpha}(G)$; in this case, we say $\rho$ lifts to $\Gamma$.
	The next result follows from the proof of \cite[Chapter 4, Lemma 1.2]{karpilovsky3}. 
	
	\begin{theorem} \label{bijective correspondence}
		With the above notations, for $\alpha = \operatorname{tra}(\chi)$, the sets $\operatorname{Irr}(G^{*} \mid \chi)$ and $\operatorname{Irr}^{\alpha}(G)$ are in dimension-preserving bijective correspondence, given by the map
		\[
		\Gamma \longmapsto \rho.
		\]
		In particular, the sets
		\[
		\operatorname{Irr}(G^{*})
		\quad \text{and} \quad \bigcup_{[\alpha] \in \mathrm{H}^{2}(G, \mathbb{C}^{\times})}
		\operatorname{Irr}^{\alpha}(G)
		\] are in dimension-preserving bijective correspondence. 
	\end{theorem}
	
	\iffalse
	\begin{corollary} \label{bijective correspondence} Let $\alpha = \operatorname{tra}(\chi)$. As $G\cong G^*/A$, the map
		(as defined in Theorem \ref{bijective correspondence via inflation}) \[
		\phi : \mathrm{Irr}(G^*) \longrightarrow 
		\bigcup_{\{[\beta] \in \mathrm{H}^2(G, \mathbb{C}^{\times}) \mid \mathrm{inf}([\beta]) = [\alpha]\}} 
		\mathrm{Irr}^\beta(G/A)
		\]
		define a bijection between the sets $\operatorname{Irr}(G^{*} \mid \chi)$ and $\operatorname{Irr}^{\alpha}(G)$.
		In particular, the sets
		\[
		\bigcup_{[\alpha] \in \mathrm{H}^{2}(G, \mathbb{C}^{\times})}
		\operatorname{Irr}^{\alpha}(G)
		\quad \text{and} \quad
		\operatorname{Irr}(G^{*})
		\]
		are in bijective correspondence.
	\end{corollary}
	
	%\poonam{I have added brackets around $\alpha$.}
	
	\begin{remark} \label{bijective correspondence for $G^*$}
		It is easy to see that if $\tilde{\rho} \in \operatorname{Irr}^{\alpha}(G)$, then $\rho : G^* \to \mathrm{GL}_n(\C)$ defined by
		\[
		\rho(a \mu(g)) = \chi(a)\tilde{\rho}(g) \quad \text{for all } a \in A, g \in G  
		\] is such that $\phi(\rho) = \tilde{\rho}$.
	\end{remark}
	\fi
	
	We now summarize the following results from  \cite[Chapter 10, Lemma 4.3, Lemma 4.12]{karpilovsky}.
	
	\begin{lemma} \label{faithful projective representation}
		
		%\item If $G$ possesses a faithful irreducible projective representation, then $Z(G^{*}) \cong \mathrm{H}^2(G,\mathbb{C}^{\times})$.
		
		Let $\alpha \in Z^2(G,\mathbb{C}^{\times})$.
		\begin{enumerate}
			\item If $\rho \in \operatorname{Irr}^{\alpha}(G)$ is faithful, then $1$ is the only $\alpha$-regular element of $Z(G)$.
			
			\item Suppose that $G$ is nilpotent. If $1$ is the only $\alpha$-regular element of $Z(G)$, then every $\alpha$-representation of $G$ is faithful.
			
			%\item $G$ has a faithful irreducible projective representation if and only if $G \cong H/Z(H)$ for some finite group $H$ with cyclic centre $Z(H)$.
			%\end{enumerate}
			%\end{enumerate}
		\end{enumerate}
	\end{lemma}
	
	\begin{lemma} \label{restriction}
		Let $G$ be a nilpotent group. Suppose that $\alpha \in Z^2(G,\mathbb{C}^{\times})$ and $\rho \in \mathrm{Irr}^\alpha(G)$. Then, the following hold:
		\begin{enumerate}
			\item If $\rho$ is faithful, then any element of $\mathrm{Irr}^\alpha(G)$ is faithful.
			\item $\rho$ is faithful if and only if $\rho|_{Z(G)}$ is faithful. 
		\end{enumerate}
		
	\end{lemma}
	
	\begin{proof}
		(1) follows  from Lemma \ref{faithful projective representation}. 
		
		(2) If ker $\rho \neq 1$, then the result follows by the fact that ker $\rho$ intersects $Z(G)$ non-trivially (see \cite[Chapter 10, Lemma 1.7]{karpilovsky2}). 
		
	\end{proof}
	
	%\begin{corollary}
	%Let $G$ be a $p$-group and $\tau_{irr}(G)=p^k$. If $\rho \in \mathrm{Irr}^\alpha(G)$ is faithful and $\operatorname{deg}(\rho)=p^k$, then the degree of every irreducible $\alpha$-representation is atleast $p^k$.    
	%\end{corollary} 
	
	\begin{corollary} \label{restriction on A is faithful}
		Let $G$ be a $p$-group such that $Z(G)$ is cyclic. Suppose that $\alpha \in Z^2(G,\mathbb{C}^{\times})$ and $\rho \in \mathrm{Irr}^\alpha(G)$. Then, $\rho$ is faithful if and only if $\rho|_A$ is faithful for any non-trivial central subgroup $A$ of $G$.   
	\end{corollary}
	
	\begin{proof}
		If $\rho$ is faithful, then clearly $\rho|_A$ is faithful. Now, let $Z(G)= \langle x \rangle$ and $A= \langle x^{p^m} \rangle$. Assume that $\rho|_A$ is faithful. If $\rho|_{Z(G)}$ is not faithful, then $\exists$ $g \in Z(G)$ such that $\rho(g)= \lambda I$ for some $\lambda \in \C^\times$. Let $g= x^{kp^n}$ for some $1 \leq k \leq p-1$. If $m \leq n$, then $g \in A$, which leads to a contradiction. Suppose $m> n$. As $Z(G)=\langle x \rangle= \langle x^k \rangle$, there is a $r>0$ such that  $\rho(x^{p^m})= \rho(x^{krp^m})= \lambda_1 I$ for some $\lambda_1 \in \C^\times$, a contradiction.  
		Hence, $\rho|_{Z(G)}$ is faithful. The result now follows from Lemma \ref{restriction}.
	\end{proof}
	
	\iffalse
	Now the following cases occur.  
	
	\textbf{Case (a)} If $p^m|p^n$, then $g \in A$ which leads to a contradiction. 
	
	\textbf{Case (b)} Suppose $p^n|p^m$. As $Z(G)=\langle x \rangle=\langle x^k \rangle$, there is a $r>0$ such that  $\rho(x^{p^m})=\rho(x^{krp^m})=\lambda_1 I$ for some $\lambda_1 \in \C^\times$, a contradiction.  
	Hence, $\rho|_{Z(G)}$ is faithful. The result now follows from Lemma \ref{restriction}.
	\fi
	
	\begin{lemma}[\cite{karpilovsky3}, Chapter 8, Theorem 2.21] \label{properties of abelian groups}
		Let $G$ be a finite abelian group. Suppose $\alpha \in Z^2(G,\mathbb C^\times)$ and $G_0$ is the subgroup of $G$ consisting of all $\alpha$-regular elements of $G$. Then $\alpha|_{G_0 \times G_0}$ is a coboundary and $\operatorname{deg}(\rho) = \sqrt{|G/G_0|}$ for any $\rho \in \mathrm{Irr}^\alpha(G)$.
	\end{lemma}  
	
	Frucht \cite{frucht} showed that a finite abelian group $G$ has a faithful irreducible projective representation if and only if $G$ is of symmetric type, i.e., $G$ can be written as a direct product of two isomorphic subgroups. 
	
	\begin{theorem} \label{projective irreducible embedding degree of finite abelian groups}
		Let $G$ be a finite abelian group of symmetric type. Then $\tau_{irr}(G) = \sqrt{|G|}$.    
	\end{theorem}
	
	\begin{proof}
		This is a consequence of \cite[Chapter 10, Lemma 4.5]{karpilovsky}.     
	\end{proof}
	
	We now state some results on capable groups. A group $\Gamma$ is called capable if there exists a group $\triangle$ with $\Gamma \cong \triangle/Z(\triangle)$. The epicentre of $G$ is denoted by $Z^{*}(G)$, which is the smallest central subgroup of $G$ such that $G/Z^*(G)$ is capable.
	
	\begin{theorem}[\cite{schurmultiplier}, Theorem 2.5.10] \label{epicenter} Let $Z$ be a central subgroup of $G$. Then the following conditions are equivalent:
		\begin{enumerate}
			\item[(a)] $Z \subseteq Z^*(G)$
			\item[(b)] $\mathrm{inf} : \mathrm{H}^2(G/Z,\mathbb{C}^{\times})\rightarrow\mathrm{H}^2(G,\mathbb{C}^{\times})$ is surjective.   
		\end{enumerate}
	\end{theorem}
	
	Recall that a group $G$ is said to be a central product of its normal subgroups $H$ and $K$ if $G=HK$, $[H, K]= 1$, and $H \cap K \neq 1$.
	
	\begin{lemma} \label{central product}
		Let $G$ be a central product of its subgroups $H$ and $K$ with $H^{\prime} \cap K^{\prime} \neq 1$. Then $G$ is not capable.      
	\end{lemma}
	
	\begin{proof}
		Let $Z = H^{\prime}$ $\cap$ $K^{\prime}$. By \cite[Theorem 3.1]{central}, inf : $\mathrm{H}^2\left(G / Z, \mathbb{C}^{\times}\right) \rightarrow \mathrm{H}^2\left(G, \mathbb{C}^{\times}\right)$ is a surjective map. Thus, the proof follows from Theorem \ref{epicenter}. 
		
	\end{proof}
	The next result follows from the proof of \cite[Chapter 4, Theorem 3.1]{karpilovsky2}.%\cite[Theorem 1]{Bekka}. 
	\begin{theorem} \label{capable groups}
		Let $G$ be a group. The following properties are equivalent:
		\begin{enumerate}
			\item[(i)] $G$ affords an irreducible faithful projective representation;
			\item[(ii)] there exists a group $H$ admitting an irreducible faithful ordinary representation such that $H/Z(H) \cong G$ .   
		\end{enumerate}
		Further, if $\delta_{irr}(H)$ exists, then $\tau_{irr}(G)$ also exists and satisfies $\tau_{irr}(G) \le \delta_{irr}(H)$. %$H$ can be chosen so that $Z(H) \subseteq H'$ and $Z(H) \subseteq M(G)$. 
		
	\end{theorem}
	
	\begin{lemma}[\cite{Beyl}, Corollary 8.2] \label{capability of extra special $p$-group}
		An extra-special $p$-group is capable if and only if it is either a dihedral group of order $8$ or of order $p^3$ having exponent $p > 2$.    
	\end{lemma}   
	
	%\begin{proof} These assertions follow easily from Lemma \ref{one dim direct sum in non-abelian group}.  \end{proof}
	
	%Next, we record some general properties concerning projective representations of abelian groups.  %We now record some general properties concerning irreducible $\alpha$-representations of $G$. 
	
	\section{Results on Faithful Projective Representations of $G$} \label{faithful_projective}
	Our first main result provides a cohomological characterization of the existence of faithful irreducible projective representations.
	
	\begin{theorem} \label{cohomological characterization}
		A p-group $G$ has a faithful irreducible $\alpha$-representation if and only if $[\alpha] \notin \mathrm {Im}\big(\operatorname{inf} : \mathrm{H}^2\left(G / N, \mathbb{C}^{\times}\right) \rightarrow \mathrm{H}^2\left(G, \mathbb{C}^{\times}\right)\big)$ for any non-trivial central subgroup $N$ of $G$.    
	\end{theorem}
	
	\begin{proof}
		Let $\rho \in \mathrm{Irr}^\alpha(G)$ be not faithful. Then, for $N= \operatorname{ker}(\rho) \cap Z(G)$, we have the exact sequence
		$$
		\mathrm{H}^2(G / N, \mathbb{C}^{\times}) \xrightarrow{\text {inf}} \mathrm{H}^2(G, \mathbb{C}^{\times}) \xrightarrow{\chi} \mathrm{H}^2(N, \mathbb{C}^{\times}) \oplus \mathrm{Hom}(G \otimes N, \mathbb{C}^{\times})
		$$
		where $\chi=(\operatorname{res}, \psi)$ as defined by Iwahori and Matsumoto \cite{matsumoto}. To be more precise, $\operatorname{res}: \mathrm{H}^2(G, \mathbb{C}^{\times}) \rightarrow \mathrm{H}^2(N, \mathbb{C}^{\times})$ is the restriction homomorphism and $\psi: \mathrm{H}^2(G, \mathbb{C}^{\times}) \rightarrow \mathrm{Hom}(G \otimes N, \mathbb{C}^{\times})$ is defined as $\psi(\xi)(\overline{x}, n)=f(x, n)f(n, x)^{-1}$ for all $\overline{x}=x G^{\prime} \in G / G^{\prime}$ and $n \in N$, where $\xi\in \mathrm{H}^2(G, \mathbb{C}^{\times})$ and $f$ is a cocycle representative of $\xi$. 
		Since $\alpha|_{\operatorname{ker}(\rho) \times \operatorname{ker}(\rho)}$ is a coboundary,  $\operatorname{res}([\alpha])$ is trivial. We also have $\alpha(x, n)= \alpha(n, x)$ $\forall$ $n$ $\in$ $N$, $x\in G$. Therefore, $[\alpha] \in \operatorname{ker}(\psi)$. Thus, $[\alpha] \in \mathrm{ker}(\operatorname{\chi}) = \mathrm{Im} (\operatorname{inf})$. 
		
		Conversely, let $N$ be a non-trivial central subgroup of $G$, $[\alpha] = \operatorname{inf}([\beta])$ for some $[\beta] \in \operatorname{H}^2(G/N, \mathbb{C}^{\times})$ and $\rho\in \mathrm{ Irr}^\alpha(G)$. 
		%Fix a $[\beta]$ $\in$ $\mathrm{H}^2\left(G / N, \mathbb{C}^{\times}\right)$ such that $\operatorname{inf}([\beta]) = [\alpha].$
		Then $\alpha(g,h)= \beta(gN,hN)$ $\forall$ $g,h \in G$. If $a \in N$, then $\alpha(a,g) = \alpha(g,a)= 1$ $\forall$ $g \in G$ and so $a$ is an $\alpha$-regular element of $Z(G)$. Hence, $\rho$ cannot be faithful, due to Lemma \ref{faithful projective representation}$(1)$.
		
	\end{proof}
	
	% As a corollary, we obtain the following result for non-capable groups. The next result for the central product of groups also follows immediately, as a corollary of the above result.  \Nayak{Result (cite the reference)}
	
	\begin{corollary} \cite[Chapter 4, Theorem 3.1]{karpilovsky2} \label{non-capable groups}
		Let $G$ be a non-capable group. Then, $G$ does not have a faithful irreducible projective representation.    
	\end{corollary}
	
	\begin{proof}
		By Theorem \ref{epicenter}, inf : $\mathrm{H}^2(G/Z^*(G),\mathbb{C}^{\times})\rightarrow \mathrm{H}^2(G,\mathbb{C}^{\times})$ is surjective. Therefore, the result follows by the same argument used in the second paragraph of the proof of Theorem \ref{cohomological characterization}. 
		
	\end{proof}
	
	%As an application, we derive the following result for the central product of groups. We first state a theorem that will be used in proving our next theorem.
	
	Let $G$ be a $p$-group with a subgroup $A$ such that $A \subseteq Z(G) \cap G^\prime$. If $[\alpha] \in \mathrm{Im} (\mathrm{inf})$, then $\exists$ $\beta \in \mathrm{Z}^2(G/A, \C^{\times})$ such that $\alpha(g,h) = \beta(gA,hA)$ $\forall g,h \in G$. Recall the map $\phi$ from Theorem \ref{bijective correspondence via inflation} which gives a bijection between $\mathrm{Irr}^\alpha(G)$ and $\bigcup_{\{[\beta] \in \mathrm{H}^2(G/A, \mathbb{C}^{\times}) \mid \mathrm{inf}([\beta]) = [\alpha]\}} 
	\mathrm{Irr}^\beta(G/A)$. Under these assumptions and notation, we have the following result.
	\begin{theorem} \label{direct_sum} 
		Suppose $\rho_i \in \mathrm{Irr}^\alpha(G)$ and  $\phi(\rho_i)= \tilde\rho_i$ for $i=1,...,k$. Then $\tilde\rho_i$ is a $\beta\mathrm{tra}(\chi_i)$-representation of $G/A$ for some $\chi_i\in \mathrm{Hom}(A,\mathbb{C}^{\times})$.
		Furthermore, if one of the following conditions is satisfied, 
		\begin{enumerate}
			\item $Z(G)$ is cyclic,
			\item $\big(\cap_{i=1}^k \operatorname{ker}(\tilde\rho_i)\big)\cap \big(Z(G)/A \big)= 1$,
		\end{enumerate}
		then $\oplus_{i=1}^k \rho_i$ is faithful on $G$ if and only if $\oplus_{i=1}^k \chi_i$ is faithful on $A$.
		
	\end{theorem}
	
	%\item If $\mathrm{inf}: \mathrm{H}^2(G/Z(G), \C^{\times}) \to \mathrm{H}^2(G, \C^{\times})$ is surjective, then $\tau(G) = \min\{\deg(\rho_1 \oplus \rho_2) \mid \chi_1 \oplus \chi_2 \text{ is a faithful projective representation of } Z(G)\}.$
	
	\begin{proof}
		By Remark \ref{remark for bijective correspondence via inflation}$(1)$, $\tilde\rho_i$ is a $\beta\mathrm{tra}(\chi_i)$ representation of $G/A$ for some $\chi_i \in \mathrm{Hom}(A,\mathbb{C}^{\times})$.
		\item 
		\begin{enumerate}
			%\item It follows directly from Lemma \ref{one dim direct sum in non-abelian group}. 
			\item It is easy to see that $\chi_i= \rho_i|_A$. The result now follows from Corollary \ref{restriction on A is faithful}. 
			%assume that $\oplus_{i=1}^k \rho_i$ is not faithful. Let $Z(G) = \left< g \right>$. Let $A=\left< g^i \right>$. Then, $\exists$ $a \in Z(G)$ such that $\oplus_{i=1}^k \rho_i(g^j)=\lambda I$. Take $i=3$, $j=5$.
			\item 
			Consider the central extension
			
			\[
			1 \to A \to Z(G) \to Z(G)/A \to 1
			\]
			Let $\mu: Z(G)/A \to Z(G)$ be a section of $Z(G) \twoheadrightarrow Z(G)/A$.
			Let $x \in Z(G)$ such that $(\oplus_{i=1}^k \rho_i)(x)=\lambda I$ for some $\lambda \in \mathbb{C}^{\times}$. Then $x= a\mu(gA)$
			for some $gA \in Z(G)/A$, $a\in A$ and by Remark \ref{remark for bijective correspondence via inflation}(2), $\chi_i(a) \tilde\rho_i(gA) = \lambda I$ $\forall$ $1 \leq i \leq k$. From the given hypothesis, it follows that $g \in A$ and so $x \in A$. Thus, $\oplus_{i=1}^k \rho_i$ is faithful on $Z(G)$ if and only if $\oplus_{i=1}^k \chi_i$ is a faithful on $A$. Hence, the result follows from Lemma \ref{restriction}$(2)$. 
			
		\end{enumerate}
	\end{proof}
	
	\subsection{Direct products} Our next aim is to describe the faithful irreducible projective representations of the direct product of groups. 
	For direct products of groups $G$ satisfying $Z(G) \subseteq [G, G]$, we characterize cocycles $\alpha$ for which the corresponding irreducible $\alpha$-representations are faithful.   
	\begin{theorem}[\cite{karpilovsky}, Chapter 2, Theorem 3.13] \label{schur multiplier of direct product}
		If $G = H \times K$, then 
		$$
		\mathrm{H}^2(G, \mathbb{C}^{\times}) \cong\mathrm{H}^2(H, \mathbb{C}^{\times}) \times \mathrm{H}^2(K, \mathbb{C}^{\times}) \times \mathrm{Hom}(H \otimes K, \mathbb{C}^{\times}).$$ 
	\end{theorem}
	% Let $\beta\in Z^{2}(H,\mathbb{C}^{\times})$, $\gamma  \in Z^{2}(K,\mathbb{C}^{\times})$ and $f \in \mathrm{Hom}(H/H' \otimes_\mathbb{Z} K/K', \mathbb{C}^{\times})$.
	The above isomorphism $\mathrm{H}^2(H, \mathbb{C}^{\times}) \times \mathrm{H}^2(K, \mathbb{C}^{\times}) \times \mathrm{Hom}(H \otimes K, \mathbb{C}^{\times}) \to \mathrm{H}^2(H \times K, \mathbb{C}^{\times})$ is given 
	by the map
	\[
	([\beta],[\gamma],f) \longmapsto [\alpha], 
	%~~\beta\in Z^{2}(H,\mathbb{C}^{\times}), \gamma  \in Z^{2}(K,\mathbb{C}^{\times}), f \in \mathrm{Hom}(H \otimes K, \mathbb{C}^{\times})
	\]
	such that $\alpha(h_1k_1,h_2k_2)= \beta(h_1,h_2)\,\gamma(k_1,k_2)\,f(\overline{h_1}\otimes\overline{k_2}), $ for $ h_1,h_2 \in H, $  $k_1,k_2 \in K, \overline{h_1}= h_1H', \overline{k_2}= k_2K'$.
	Via this isomorphism, we write $\alpha=(\beta,\gamma,f)$. \\
	
	%, where $\beta= \alpha|_{H \times H} \in Z^{2}(H,\mathbb{C}^{\times})$, $\gamma = \alpha|_{K \times K} \in Z^{2}(K,\mathbb{C}^{\times})$and $f \colon H^{2}(G, \mathbb{C}^{\times}) \to \operatorname{Hom}(H \otimes K, \mathbb{C}^{\times})$ is a homomorphism given by \[\nu([\alpha])(\tilde{h} \otimes \tilde{k}) = \alpha(h, k) \alpha(k, h)^{-1}, \] for $\tilde{h} = h H'$ and $\tilde{k} = k K'$. 
	
	
	Recall that, for two matrices $A=(a_{ij})_{m\times m}$ and $B=(b_{ij})_{n\times n}$, their tensor product is defined by
	\[
	A \otimes B
	=
	\begin{pmatrix}
		a_{11}B & \cdots & a_{1m}B \\
		\vdots  & \ddots & \vdots \\
		a_{m1}B & \cdots & a_{mm}B
	\end{pmatrix}.
	\]
	The next result is immediate.
	\begin{lemma} \label{tensor product of matrices}
		Let $A \in GL_m(\mathbb{C})$ and $B \in GL_n(\mathbb{C})$. Then $A \otimes B = \lambda I_{mn}$ for some $\lambda \in \mathbb{C}^\times$ if and only if there exist scalars $\lambda_1, \lambda_2 \in \mathbb{C}^\times$ such that $A = \lambda_1 I_m \text{ and } B = \lambda_2 I_n.$
	\end{lemma}
	
	Let $\rho_1: H \to GL(V_1)$ be a $\beta$-representation of $H$ and $\rho_2: K \to GL(V_2)$ be a $\gamma$-representation of $K$. The tensor product $\rho_1 \otimes \rho_2: H \times K \to GL(V_1\otimes V_2)$ is defined by $$(\rho_1 \otimes \rho_2) (h,k)= \rho_1(h) \otimes \rho_2(k), h \in H, k \in K.$$ It is easy to check that $\rho_1 \otimes \rho_2$ is a $(\beta, \gamma, 1)$-representation of $G$
	(\cite[Chapter 5, Section 1]{karpilovsky}).
	
	%\begin{remark} \label{tensor product definition}Let $\alpha= (\beta,\gamma,1)$ and $\rho \in \mathrm{Irr}^\alpha(G)$. Then, by \cite[Chapter 5, Corollary 1.3]{karpilovsky}, $\rho=\rho_1 \otimes \rho_2$ for some $\rho_1\in \mathrm{Irr}^\beta(H)$ and $\rho_2\in \mathrm{Irr}^\gamma(K)$.     \end{remark}
	
	\begin{theorem} \label{tensor product is faithful} 
		Let $G = H \times K$ and $\alpha = (\beta,\gamma,1)$ $\in$ $Z^{2}(G,\mathbb{C}^{\times})$. Then, up to linear equivalence, an irreducible $\alpha$-representation $\rho$ is of the form $\rho_1 \otimes \rho_2$ for some $\rho_1\in \mathrm{Irr}^\beta(H), \rho_2\in \mathrm{Irr}^\gamma(K)$. Furthermore,
		$\rho$ is faithful on $G$ if and only $\rho_1$ and $\rho_2$ are faithful on $H$ and $K$ respectively.
		%If $\rho_1$ and $\rho_2$ are projective representations of $H$ and $K$ respectively, then $\rho_1 \otimes \rho_2$ is a faithful projective representation of $H \times K$ if and only if $\rho_1$ and $\rho_2$ are faithful.
	\end{theorem}
	
	\begin{proof}
		By \cite[Chapter 5, Corollary 1.3]{karpilovsky}, it follows that, up to linear equivalence, $\rho= \rho_1 \otimes \rho_2$ for some $\rho_1 \in \mathrm{Irr}^\beta(H)$ and $\rho_2 \in \mathrm{Irr}^\gamma(K)$. 
		
		%We show that if $\rho_1$ and $\rho_2$ are projective representations of $H$ and $K$ respectively, then $\rho_1 \otimes \rho_2$ is a faithful projective representation of $H \times K$ if and only if $\rho_1$ and $\rho_2$ are faithful. 
		
		Suppose $\rho= \rho_1 \otimes \rho_2$ is faithful and $\rho_2$ is not faithful.
		Then, $\exists$ $k^{\prime} \in K$, $k^{\prime} \neq 1$ such that $\rho_{2}(k')= d I_{n \times n}$ for some $d \in \mathbb{C}^{\times}$. Then, we have $$(\rho_1 \otimes \rho_{2})\left(1,k^{\prime}\right)= \rho_{1}\left(1\right) \otimes \rho_{2}(k') = I_{m\times m}  \otimes d I_{n \times n}= d I_{m n \times m n}.$$ 
		This contradicts our assumption that $\rho_{1} \otimes \rho_{2}$ is faithful. So, $\rho_{2}$ must be faithful. Using similar arguments, one can verify that $\rho_{1}$ is also faithful.
		Conversely, let $\rho_1$ and $\rho_2$ be faithful.
		If $(\rho_{1} \otimes \rho_{2})(h,k)=$ $\lambda I$ for some $\lambda \in \mathbb{C}^{\times}$, then $\rho_{1}(h) \otimes \rho_{2}(k) = \lambda I$. Hence, the result follows from Lemma \ref{tensor product of matrices}. 
	\end{proof}
	
	\begin{theorem} \label{directproduct}
		Let $G = H \times K$, where $H$ and $K$ are non-abelian $p$-groups. Let $\alpha$ $\in$ $Z^{2}(G,\mathbb{C}^{\times})$. 
		\begin{enumerate}
			\item If $\alpha=(1,1,f)$, then $\rho \in \mathrm{Irr}^\alpha(G)$ is not faithful.
			
			\item If $Z(G) \subseteq G^{\prime}$ and $\alpha = (\beta,\gamma,f)$, then an irreducible $\alpha$-representation of $G$ is faithful if and only if $1$ is the only $\beta$-regular element of $Z(H)$ and $1$ is the only $\gamma$-regular element of $Z(K)$. Furthermore, $G$ admits a faithful irreducible $\alpha$-representation if and only if $H$ and $K$ admit faithful irreducible $\beta$ and $\gamma$-representation respectively.     
		\end{enumerate}
	\end{theorem}
	%In this section, we discuss the existence of faithful irreducible projective representations of the direct product of groups.
	
	\begin{proof}
		(1) Let $N = \big(H^{\prime}$ $\cap$ $Z(H)\big) \times \big(K^{\prime}$ $\cap$ $Z(K)\big)$. Since $\alpha|_{N \times G} = \alpha|_{G \times N}= 1$, there exists a $\beta\in \mathrm{Z}^2(G/N, \C^{\times})$ such that $\operatorname{inf}([\beta])=[\alpha]$.  
		The result now follows from Theorem \ref{cohomological characterization}. 
		\\  
		
		(2) By Lemma \ref{faithful projective representation}, it suffices to show that $1$ is the only $\alpha$-regular element of $Z(G)$ if and only if $1$ is the only $\beta$-regular element of $Z(H)$ and $1$ is the only $\gamma$-regular element of $Z(K)$.    
		Let $(h,k)$ be an $\alpha$-regular element of $Z(G)$. Then, $\alpha \big((h,k),(h_1,k_1)\big)= \alpha\big((h_1,k_1),(h,k)\big)$ $\forall$ $(h_1,k_1) \in G$. Putting $k_1 = 1$, we deduce that $\beta(h,h_1) = \beta(h_1,h)$ $\forall$ $h_1 \in H$. Hence, $h$ is a $\beta$-regular element of $Z(H)$. Similarly, it can be shown that $k$ is a $\gamma$-regular element of $Z(K)$.
		Conversely, let $h$ be a $\beta$-regular element of $Z(H)$ and $k$ be a $\gamma$-regular element of $Z(K)$. It is easy to check that $\alpha\big((h,k),(h_1,k_1)\big)= \alpha\big((h_1,k_1),(h,k)\big)$ $\forall$ $(h_1,k_1)\in G$. Therefore, $(h,k)$ is an $\alpha$-regular element of $Z(G)$. This completes the proof of the result. 
	\end{proof}
	
	\section{$\tau(G)$ and $\tau_{irr}(G)$ for finite groups $G$} \label{emdedding degree} We first establish a relationship between $\delta(G)$ and $\tau(G)$. 
	
	%\subsection{Relation between embedding degree and projective embedding degree}
	
	\begin{lemma} \label{relation}
		For any finite group $G$, $\tau(G) \le \delta(G)+ 1$.
	\end{lemma}
	
	\begin{proof} Let $\delta(G)= n$ and $\psi$ be a faithful ordinary representation of $G$ of degree $n$. 
		Suppose $\rho$: $G\,\to\,GL_{n+1}(\mathbb{C})$ is defined by \[\rho(g)=
		\begin{bmatrix}
			\psi(g) & 0 \\
			0 &  1
		\end{bmatrix}.
		\]
		Observe that $\rho$ is a faithful projective representation of $G$. Hence, the result follows.
	\end{proof}
	Note that the above bound is optimal, as is evident from the next theorem. 
	%The above result shows that the inequality in Lemma \ref{relation} provides the best possible bound between $\tau(G)$ and $\delta(G)$, and it cannot be improved further. \\
	\begin{lemma} \label{group with cyclic centre of order p}
		Let $G$ be a $p$-group such that $M(G)= 1$ and $Z(G)$ is of order $p$. Then $\tau_{irr}(G)$ does not exist and $\tau(G)= \delta(G)+1$.
	\end{lemma}
	
	\begin{proof}
		
		Since $M(G)=1$, by Lemma \ref{projectively equivalent to an ordinary representation} and Lemma \ref{projectively equivalent}, it suffices to consider only ordinary representations of $G$.
		Since faithful projective representations are also faithful ordinary representations, by Lemma \ref{relation}, we have $\dg \leq \tau(G) \leq \dg+1$. Let $\tau(G)= \dg = m$ and $\rho$ be an ordinary representation of degree $m$ which is faithful as a projective representation of $G$. We write $\rho = \oplus_{i=1}^k \rho_i$, where $\rho_i \in \mathrm{Irr}(G)$. By Lemma \ref{one dim direct sum in non-abelian group}$(2)$, it follows that $\tau_{irr}(G)$ does not exist. Hence, $k > 1$. 
		Since $\rho_i$ is not a faithful ordinary representation, by the fact used in the proof of Lemma \ref{restriction}, $\rho_i|_{Z(G)}$ cannot be a faithful ordinary representation. Thus, $\rho_i|_{Z(G)}= n_i 1_{Z(G)}$ as $Z(G)$ is of order $p$. Therefore, $\rho|_{Z(G)}= n 1_{Z(G)}$ for some $n \in \mathbb N$. Since $\rho|_{Z(G)}$ is not faithful, by Lemma \ref{restriction}, $\rho$ cannot be a faithful on $G$. Consequently, $\tau(G)= \delta(G)+1$. 
		
	\end{proof}
	In the next result, we describe $\tau_{irr}(G)$ and $\tg$ for $G= H \times K$. %We also analyze how $\tau(H \times K)$ relates to $\tau(H)$ and $\tau(K)$.
	
	%\subsection{Direct product of groups}
	
	% In this section, we discuss $\tau_{irr}(G)$ for the direct product of groups. 
	
	\begin{theorem} \label{projective irreducible embedding degree for direct product} Let $G = H \times K$. \\
		(i) If $gcd(|H|,|K|)=1$, then \textbf{$\tau_{irr}(G) = \tau_{irr}(H) \tau_{irr}(K)$}.   \\
		(ii) If $M(H)= 1$ and $M(K)= 1$, then \textbf{$\tau(G) \leq \tau(H) + \tau(K)$}. \\
		(iii) \textbf{$\tau(G) \leq \tau(H) \tau(K)$}.    
	\end{theorem}
	
	\begin{proof} 
		(i) Let $\tau_{{irr}}(H)=m$,  $\tau_{irr}(K)=n$ and $\tau_{irr}(G)=d$. 
		Since $\gcd(|H|, |K|) = 1$, by Theorem~\ref{schur multiplier of direct product}, any cocycle $\alpha \in Z^2(G,\mathbb{C}^\times)$ is of the form $\alpha = (\beta, \gamma, 1).$ 
		Hence, by Theorem \ref{tensor product is faithful}, any $\rho \in \mathrm{Irr}^\alpha(G)$ is of the form $\rho = \rho_{1} \otimes \rho_{2}$, where $\rho_{1} \in \operatorname{Irr}^{\beta}(H), \rho_{2} \in \operatorname{Irr}^{\gamma}(K)$, and
		$\rho$ is faithful if and only if $\rho_{1}, \rho_{2}$ are also faithful. Since 
		%$\operatorname{deg} (\rho_{1})\geq m,\operatorname{deg} (\rho_{2}) \geq n$, 
		deg $(\rho_{1} \otimes \rho_{2}) =(\operatorname{deg} \rho_{1})(\operatorname{deg} \rho_{2}), ~d = mn$ and the result follows.
		
		%Let $\phi_1\in \operatorname{Irr}^{\beta}(H)$ and $\phi_2\in \operatorname{Irr}^{\gamma}(K)$ be faithful. Let $\operatorname{deg} (\phi_{1})=m$ and $\operatorname{deg} (\phi_{2})=n$. By Theorem \ref{tensor product is faithful}, $\phi_1\otimes \phi_2$ is a faithful irreducible $\alpha$-representation of $G$ of degree $mn$ for $\alpha =(\beta, \gamma, 1)$. Therefore, $d = mn$.
		
		(ii) Let $\tau(H)= m$ and $\tau(K)= n$. Let $\eta$ and $\psi$ be faithful $\beta$ and $\gamma$-representation of $H$ and $K$ of degree $m$ and $n$ respectively. Since $M(H)=1$ and $M(K)=1$, we can assume that $\beta= \gamma=1$. Define $\sigma: G \to GL_{m+n}(\mathbb{C})$ by:
		\[
		\sigma(h,k)= \begin{bmatrix}
			\eta(h) &  \\
			& \psi(k)
		\end{bmatrix}
		\]
		It is easy to verify that $\sigma$ is a faithful $\alpha$-representation of $G$ for $\alpha =(1,1,1)$. Thus, $\tau(G) \leq \tau(H) + \tau(K)$.
		
		(iii) Let $\eta$ and $\psi$ be faithful $\beta$ and $\gamma$-representation of $H$ and $K$ respectively. Using similar arguments as in the proof of Theorem \ref{tensor product is faithful}, $\eta \otimes \psi$ is faithful. Consequently, the result follows.   
	\end{proof}
	
	\iffalse
	\begin{proposition} Let $N \trianglelefteq G$. Then the following holds. \\
		(i) $\tau(G/N) \leq \tau(G)$. \\
		(ii) If $\tau_{irr}(G)$ exists, then $\tau_{irr}(G/N)$ exists and satisfies $\tau_{irr}(G) \leq \tau_{irr}(G/N)$.  
	\end{proposition}
	
	\begin{proof}
		(i) Let $\rho$ be a faithful projective representation of $G$ of degree $\tau(G)= m$. Suppose $L$ is a transversal for $N$ in $G$ containing $1$ and let $\rho': G/N \to \GL_m(\C)$ be defined by $\rho'(gA)= \rho(g)$ ($g \in L)$.   
	\end{proof}
	\fi
	\subsection{Non-capable groups $G_1$ and $G_2$ for which $\tau(G_1)= \tau(G_2)$}
	\iffalse
	%Criteria for groups $G_1$ and $G_2$ to have $\tau(G_1)= \tau(G_2)$
	Let $G_1,G_2$ be two groups and $Z_j$ be central subgroups of $G_j$
	such that for $j=1,2$, the sequences
	$$
	1 \to \operatorname{Hom}(Z_j,\mathbb{C}^{\times})
	\xrightarrow{\operatorname{tra}_j}
	H^2(G_j/Z_j,\mathbb{C}^{\times})
	\xrightarrow{\inf_j}
	H^2(G_j,\mathbb{C}^{\times})
	\to 1
	$$
	are exact.
	\fi
	In the next result, we determine sufficient conditions under which non-capable groups have the same projective embedding degree.
	%In this section, we determine sufficient conditions for two groups to have the same projective embedding degree. \\
	We say that an isomorphism $\psi : G_2 \to G_1$ induces an isomorphism $\bar{\psi} : \Ho^2(G_1,\mathbb{C}^{\times}) \to
	\Ho^2(G_2,\mathbb{C}^{\times})$
	if $\bar{\psi}$ is defined by $\bar{\psi}([\alpha]) = [\beta]$ such that $\beta(g,h)= \alpha(\psi(g),\psi(h))$
	for $g,h \in G_2$.
	\begin{theorem} \label{tau(G_1) = tau(G_2)}
		Let $G_1$ and $G_2$ be two non-capable groups and $A_j \subseteq Z^*(G_j) \cap G_j'$ for $j= 1,2$. Suppose there exist isomorphisms $\eta : A_2 \to A_1$ and $\gamma : G_2/A_2 \to G_1/A_1$ such that Diagram~1 is commutative for the induced isomorphisms
		$\bar{\eta}$ and $\bar{\gamma}$. Then, $\tau(G_1) = \tau(G_2)$. 
		\begin{figure}[H]
			\[
			\xymatrix{
				1 \ar[r] & \mathrm{Hom}(A_{1}, \mathbb C^\times) \ar[d]^{\bar{\eta}} \ar[r] ^{\mathrm{tra}} & \Ho^2(G_1/A_{1}, \mathbb C^\times)  \ar[d]^{\bar{\gamma}} \\
				1 \ar[r] & \mathrm{Hom}(A_{2}, \mathbb C^\times) \ar[r] ^{\mathrm{tra}} & \Ho^2(G_2/A_{2}, \mathbb C^\times)\\
			}
			\] \caption{Diagram 1}
		\end{figure}
	\end{theorem}
	
	\begin{proof}
		By Theorem \ref{epicenter}, inf : $\mathrm{H}^2(G_j/A_j, \mathbb{C}^{\times}) \rightarrow\mathrm{H}^2(G_j,\mathbb{C}^{\times})$ is surjective for $j=1, 2$. Let $\rho_i \in \irr^{\alpha}(G_1)$ such that $\rho= \oplus_{i=1}^l \rho_i$ be a faithful $\alpha$-representation of $G_1$ of degree $\tau(G_1)$. By Theorem \ref{direct_sum}, there exists $\tilde{\rho_i} \in \irr^{\beta\tra(\chi_i)}(G_1/A_1)$ such that $\phi(\rho_i)= \tilde{\rho_i} $. Let $\psi_i:= \bar{\eta} (\chi_i)= \chi_i \circ \eta \in \Hom(A_2, \C^\times)$.  
		Consider $\beta'\in Z^2( G_2/A_2, \C^\times)$ such that $\beta'(gA_2,hA_2)= \beta(\gamma(gA_2), \gamma(hA_2))$ for $g,h \in G_2$. Hence, $[\beta']= \bar{\gamma}([\beta])$. Let $\alpha' \in Z^2(G_2,\C^\times)$ defined by $\alpha'(g,h)= \beta'(gA_2,hA_2)$ $\forall$ $g$,$h$ $\in$ $G_2$. Then, by the commutativity of the diagram, we have $\tilde{\rho_i}':= \tilde{\rho_i} \circ \gamma \in \irr^{\beta'\tra(\psi_i)}(G_2/A_2)$, and by Theorem \ref{direct_sum}, there exists $\rho'_i \in \irr^{\alpha'}(G_2)$ such that $\phi(\rho'_i)= \tilde{\rho_i}'$. %\poonam{Theorem \ref{direct_sum} used twice}
		
		We show that $\oplus_{i=1}^l \rho'_i$ is a faithful $\alpha'$-representation of $G_2$.
		Let $\mu: G_1/A_1 \to G_1$ and $\nu: G_2/A_2 \to G_2$ be sections of $G_1 \twoheadrightarrow G_1/A_1$ and $G_2 \twoheadrightarrow G_2/A_2$, respectively.
		Let $x \in G_2$ such that $(\oplus_{i=1}^l \rho_i')(x)= \lambda I$ for some $\lambda \in \mathbb{C}^{\times}$. Since $x= a\nu(gA_2)$
		for some $gA_2 \in G_2/A_2$, $a\in A_2$, by Remark \ref{remark for bijective correspondence via inflation}(2), $\psi_i(a) \tilde\rho_i'(gA_2)= \lambda I$ $\forall$ $1 \leq i \leq k$. This implies that $(\chi_i \circ \eta(a))(\tilde{\rho_i} \circ \gamma(gA_2))= \lambda I$ and so $\rho_i(\eta(a)\mu(\gamma(gA_2))) = \lambda I$ $\forall i$. Since $\oplus_{i=1}^l \rho_i$ is faithful on $G_1$, $\eta(a) \mu(\gamma(gA_2))=1$. As $\eta(a) \in A_1$, $\mu(\gamma(gA_2)) \in A_1$. Thus, $\gamma(gA_2) =1$. Since $\gamma$ is an isomorphism, $g \in A_2$. Now, $\eta$ is an isomorphism, forcing $a= 1$. Hence, $\oplus_{i=1}^l \rho'_i$ is faithful of degree $\tau(G_1)$. Thus, $\tau(G_2) \leq \tau(G_1)$. By similar arguments, $\tau(G_1) \leq \tau(G_2)$. This completes the proof. %$\mu(\gamma(gA_2))= 1$ and so
	\end{proof}
	
	\begin{remark} Let $A_j \subseteq Z(G_j) \cap G_j'$ such that the Diagram 1 given in Theorem \ref{tau(G_1) = tau(G_2)} is commutative for the isomorphisms $\eta$ and $\gamma$. Then, by similar arguments used in the proof of Theorem \ref{tau(G_1) = tau(G_2)}, it can be shown that $\delta(G_1) = \delta(G_2)$.   
	\end{remark}
	%\poonam{Do we need to mention that $G_i$ need not be non-capable?} %Even Theorem \ref{direct_sum} holds true for ordinary case.} 

Note that Theorem \ref{tau(G_1) = tau(G_2)} also provides a method for constructing many additional non-capable groups with the same projective embedding degree, as shown in the following result.

When $G_1$ and $G_2$ satisfy the hypotheses of Theorem \ref{tau(G_1) = tau(G_2)},
we say that the pair $(G_1, G_2)$ satisfies the hypotheses.
\begin{corollary} If two pairs $(G_1, G_2)$ and $(H_1, H_2)$ satisfy the hypotheses of Theorem \ref{tau(G_1) = tau(G_2)}. Then
	
	(i) $\tau(G_1 \times H_1) = \tau(G_2 \times H_2)$.
	
	(ii) For any finite abelian group $A$, $\tau(G_1 \times A) = \tau(G_2 \times A)$. %\poonam{Is $G_1\times A$ non-capable?}
\end{corollary} 

\begin{proof}
	(i) By \cite[Proof of Corollary 1.5]{sabnam}, $(G_1 \times H_1, G_2 \times H_2)$ satisfies the hypotheses of Theorem \ref{tau(G_1) = tau(G_2)}, the result follows. \\ 
	(ii) Since $(G_1 \times A, G_2 \times A)$ satisfies the hypotheses of Theorem \ref{tau(G_1) = tau(G_2)}, the result follows.  %is a consequence of Theorem \ref{tau(G_1) = tau(G_2)}. \poonam{Theorem \ref{tau(G_1) = tau(G_2)} used twice}
\end{proof}

\begin{example} The groups $\Phi_5(311)$, $\Phi_5(2211)a$, $\Phi_5(2211)b, \Phi_5(21^4)c$, listed from \cite{james} have same projective embedding degree, follows from \cite[Proof of Theorem 3.4]{sabnam} and Theorem \ref{tau(G_1) = tau(G_2)}. 
\end{example} 


%\Nayak{Rewritten above, check that} Let $G_1$ and $G_2$ be two groups of order $p^6$ belonging to the set \Nayak{check the list, Is it correct?} $\{\Phi_5(311)$, $\Phi_5(2211)a$, $\Phi_5(2211)b, \Phi_5(21^4)c\}$. Since $(G_1, G_2)$ satisfies the hypotheses of Theorem \ref{tau(G_1) = tau(G_2)}, $\tau(G_1) = \tau(G_2)$, follows from \cite[Proof of Theorem 3.4]{sabnam}.    


%We say that $G \sim_\C H$ via the map $\phi$ if there is an isomorphism $\phi : {\Ho}^2(G,\C^\times)\to {\Ho}^2(H,\C^\times)$ such that $\C^\alpha G \cong \C^{\phi(\alpha)}H$ for each $[\alpha]\in H^2(G,\C^\times)$.

\subsection{Finite abelian groups} \label{abelian}
In this section, we determine $\tau(G)$ for finite abelian groups. It is easy to check that $\tau(G) = \tau_{irr}(G) = 1$ if and only if $G= 1$. 
%The values of $\delta(G)$ and $\delta_{irr}(G)$ given in \cite{kaur} will be used to compute $\tau(G)$ and $\tau_{irr}(G)$ in this section, without further citation.  
Recall that the rank of a group $G$ is defined as the minimal number of generators of $G$. Note that if $G$ is a finite abelian group of rank $n$, then $\delta(G)= n$, follows from \cite[Lemma 3.4]{moreto}.    
%The following lemmas provide tools for computing $\tau(G)$ for finite abelian groups.

\begin{lemma} \label{1dim direct sum in abelian group}
	Let $G$ be a finite abelian $p$-group of rank $n$ with $|G|>1$. Suppose $\rho= \oplus_{i=1}^r \rho_i$, where $\rho_i \in \mathrm{Irr}(G)$. If $1 \leq r \leq n$, then $\rho$ is not a faithful projective representation of $G$.   
\end{lemma}

\begin{proof}
	The proof is straightforward for $n=1$. So, we assume $n>1$. First, we consider $G=(\mathbb Z/p\mathbb Z)^n$. Let $e_{k}= \underbrace{(0,0, \ldots, 1, \ldots, 0)}_\text{$1$ in the $k{\text {th}}$ position} \in G$ for $1 \leq k \leq n$.
	
	We assume that $r=n$.  
	Suppose $\rho$ is a faithful projective representation of $G$. Then, for each $k$, there exists $j_{l k} \in \{0,1, \ldots, p-1\}$ ($1\leq l \leq n $) such that all $j_{l k}$'s are not equal and $$\rho(e_{k})=
	\left[\begin{array}{ccc}e^{\frac{2 \pi ij_{1k}}{p}} &  & \\
		& \ddots & \\
		& & e^{\frac{2 \pi ij_{nk}}{p}}\end{array}\right]$$
	is a diagonal matrix. Hence, $$\rho\left(m_{1}, m_{2}, \ldots, m_{n}\right)
	=\left[\begin{array}{ccc}e^{\frac{2 \pi i \sum\limits_{k=1}^{n} m_{k} j_{1 k}}{p}} &  & \\
		& \ddots & \\
		& & e^{\frac{2 \pi i \sum \limits_{k=1}^{n} m_{k} j_{n k}}{p}}\end{array}\right].$$                                            
	Since $\rho$ is also a faithful ordinary representation, $\rho\left(m_{1}, m_{2} \ldots, m_{n}\right)= I$ implies that $\left(m_{1}, \ldots, m_{n}\right)= (0,0, \ldots, 0)$, which says
	$$
	\left[\begin{array}{ccc}
		j_{11} & \cdots & j_{1 n} \\
		\vdots & &\vdots \\
		j_{n 1} & \cdots & j_{n n}
	\end{array}\right]\left[\begin{array}{c}
		x_{1} \\
		\vdots \\
		{x}_{n}
	\end{array}\right]= \left[\begin{array}{c}
		0 \\
		\vdots \\
		0
	\end{array}\right]
	$$
	has only a trivial solution in $\Z/p\Z$. Then, the matrix $A=(j_{lk})_{n\times n}$ is invertible,
	and for $b=[1,1,\cdots, 1]^{\top}$, $A x=b$ has a unique solution $c=A^{-1} b,$ where $c \neq 0$.
	Let $c^{\top}= \left[c_{1}, c_{2}, \ldots ., c_{n}\right]$, 
	$c_{i}$ $\in \{0,1, \ldots, p-1\}$.
	Then, $\sum\limits_{k=1}^{n} j_{i k} c_{k}= 1$ for $1\leq i \leq n$. Therefore, 
	\[\rho\left(c_{1}, c_{2}, \ldots, c_{n}\right)=
	\left[\begin{array}{ccc}
		e\frac{2 \pi i \sum\limits_{k=1}^n c_k j_{1k}}{p} & & \\
		& \ddots & \\
		&  & e \frac{2 \pi i \sum\limits_{k=1}^n c_k j_{nk}}{p} 
	\end{array}\right]
	= e^{\frac{2 \pi i}{p}} I_{n \times n},
	\] which contradicts our assumption. Hence, $\rho$ cannot be faithful. 
	
	If $r<n$, let $\psi= \rho \oplus (\oplus_{i=1}^{n-r} 1_G$). If $\rho$ is faithful as projective, then $\psi$ is also faithful as projective, which gives a contradiction to Case 1. Hence, the result holds if $G$ is elementary abelian. 
	Now, let $G$ be any abelian group of rank $n$ and $H= (\Z/p\Z)^n \subseteq G$. If $\rho$ is faithful on $G$, then it is also faithful on $H$, which is not possible. This completes the proof.
	
\end{proof}

\begin{lemma} \label{abelian group of small rank}
	Let $G$ be a finite abelian $p$-group of rank $n$ such that $n \le p-1$ and $|G|>1$. Then, $\tau(G)= n+1$.    
\end{lemma} 

\begin{proof}
	Since $\delta(G)=n$, by Lemma \ref{relation}, we obtain that $\tau(G) \le n+1$. 
	Let $\rho$ be a faithful $\alpha$-representation of $G$ of degree $k$, where $k =\tau(G)$. 
	If $k< n+1$, then $\rho = \oplus_{i=1}^k\rho_i$, where each $\rho_i \in \irr^\alpha(G)$ has degree $1$. Since $\rho_is$ are one-dimensional, $\alpha$ is a coboundary. Therefore, we can assume that $\rho$ is an ordinary representation. %Consequently, the result follows from Lemma 
	However, by Lemma \ref{1dim direct sum in abelian group}, $k> n$. Hence, the result follows.
\end{proof}
%\poonam{Is 1 used correctly for groups?}
\begin{theorem} \label{elementary abelian p-group}
	Let $G \cong (\mathbb Z/p\Z)^{n}$. Then, the following holds. \\
	(i) $\tau(G) = n+1$ if $G \ncong 1, (\Z/2\Z)^2$ and $(\Z/2\Z)^4$.  \\
	(ii) $\tau(G)= n$ if $G \cong 1, (\Z/2\Z)^2$ or $ (\Z/2\Z)^4$.
\end{theorem}

\begin{proof}
	By Lemma \ref{abelian group of small rank}, it is enough to consider $n \geq p$. Let $\rho$ be a faithful $\alpha$-representation of $G$ of degree $\tau(G)$ such that $\rho = \oplus_{i=1}^m\rho_i$, where $\rho_i \in \mathrm{Irr}^\alpha(G)$. Then, there is an $s \geq 0$ such that $\operatorname{deg}(\rho_i) = p^s$ for all $i$, follows from Lemma \ref{properties of abelian groups}. Since $\delta(G)= n$, by Lemma \ref{relation}, $\tau(G) \leq n+1$. First, assume that $\tau(G) < n+1$. \\
	
	$(i)$ Let $s \geq 1$. Then $n = kp^s+r$ for some $ k \geq 0$ and $0 \leq r < p^s$. If $k=0$, then $n=r <p^s \leq \operatorname{deg}(\rho)= \tau(G)$. So, we assume $k \ge 1$. If $m \geq (k+1)$, then $\operatorname{deg}(\rho) > n$. Hence, we consider $m \leq k$. 
	%Let $\rho$ be an irreducible $\alpha$-representation of degree $p^s$. By [Proposition 1.9.14,\cite{karpilovsky3}], $\exists$ a subgroup H of G and an irreducible $\alpha$-representation $\lambda$ of H of degree 1 such that $\rho= \lambda^G$.   \\
	By Lemma \ref{properties of abelian groups}, $\alpha|_{G_0\times G_0}$ is a coboundary, so $\rho|_{G_0}$ is projectively equivalent to an ordinary representation $\rho'$ of $G_0$ and we may assume that $\rho' = \oplus_{i=1}^m \deg(\rho_i)\rho_i'$, where $\rho_i' \in \mathrm{Hom}(G_0, \mathbb C^\times)$. 
	%$o(G_0)=p^{n-2s}$ and $G_0 \subseteq H$ and $[H:G_0]=[G:H]$. Since, $\alpha|_{H \times H}$ is a coboundary, hence.
	%If $\chi_i$ is an irreducible constituent of $\rho_i'$, then by Clifford's theory, $\rho_i'= p^s$ $\chi_i$. This implies that $\rho'= p^s$ $\oplus_{i=1}^m\chi_i$ . \\
	Since $\rho|_{G_0}$ is faithful, $\rho'$ is also faithful by Lemma \ref{projectively equivalent}. Therefore, $\oplus_{i=1}^m \rho_i'$ is faithful on $G_0$. Since $G_0 \cong (\mathbb Z/p\Z)^{n-2s}$ (by Lemma \ref{properties of abelian groups}), it follows from Lemma \ref{1dim direct sum in abelian group} that rank$(G_0)=n-2s < m \leq k$ which is true only for the following cases: 
	\begin{enumerate}
		%\item $p=3$, $s=1$, $k=1$, $r=0$.  $\big(G = (\Z/3\Z)^3\big)$ \poonam{remove}
		\item $p=2$, $s=2$, $k=1$, $r=0$, i.e., $G \cong (\Z/2\Z)^4$. %\poonam{remove $(\Z/2\Z)^5$}
		%\item $p=2$, $s=1$, $k=2$, $r=0$ $\big(G \cong (\Z/2\Z)^4 \big)$ \poonam{remove}
		\item $p=2$, $s=1$, $k=1$, $r=0$, i.e., $G \cong (\Z/2\Z)^2$.  %\poonam{remove $(\Z/2\Z)^3$}     
	\end{enumerate}
	Thus, we conclude that for $G \ncong (\Z/2\Z)^2$ and $(\Z/2\Z)^4$, $\rho$ cannot be faithful when $s \geq 1$ and $\operatorname{deg}(\rho) < n+1$. \\
	For $s=0$, $\rho$ is a direct sum of one-dimensional representations, so by Lemma \ref{1dim direct sum in abelian group}, $\operatorname{deg}(\rho) \geq$ $n+1$. Hence, $\tg = n + 1$ when $G \ncong 1, (\Z/2\Z)^2, (\Z/2\Z)^4$. \\
	%\poonam{Lemma \ref{1dim direct sum in abelian group} and Lemma \ref{properties of abelian groups} is used twice} \\
	
	$(ii)$ If $G \cong (\Z/2\Z)^2$, by Theorem \ref{projective irreducible embedding degree of finite abelian groups}, $\tg=2$. 
	
	If $G \cong (\Z/2\Z)^4$, taking a subgroup $H =(\Z/2\Z)^3$, we have $\tau(H) = 4$ $\leq $ $\tg$. Now, by Theorem \ref{projective irreducible embedding degree of finite abelian groups}, $\tg = 4$.
	
\end{proof}

\begin{theorem}\label{projective embedding degree of finite abelian groups}
	Let $G$ be a finite abelian group of rank $n$ which is not elementary abelian. Then $\tau(G) = n + 1$.
	
\end{theorem}

\begin{proof}
	Let $G$ $\cong$ $\mathbb{Z}/{m_1}\Z \times \mathbb{Z}/{m_2}\Z \times \cdots \times \mathbb{Z}/{m_n}\Z$, where $m_{i+1} \mid m_i$ for $1 \leq i \leq n-1$. Let $p$ be the smallest prime dividing $m_n$. 
	
	Suppose $p = n= 2$. Then $o(G) > 4$ and consider $H_1 = \Z/2q_1\Z \times \Z/2\Z \subseteq G$ for prime $q_1$ such that $2q_1|m_1$. Since $\delta(H_1)= 2$, by Lemma \ref{relation}, $\tau(H_1) \leq 3$. Let $\psi$ be a faithful projective representation of $H_1$ of degree $\tau(H_1)$. Since $H_1$ is not of symmetric type, $\psi$ is not irreducible. Hence, using Lemma \ref{1dim direct sum in abelian group}, we have $\tau(H_1) >2$, and thus $3 =\tau(H_1)$ $\leq$ $\tau(G)$. Since $\delta(G)= 2$, Lemma \ref{relation} yields that $\tg = 3= n+1$. %\poonam{ Lemma \ref{relation} used many times}
	%\poonam{We have proved later that if $\delta(G)=2$, $\tau(G) = 3$.} 
	
	If $p =2, n= 4$, then $o(G) > 16$ and consider $H = \Z/2q_2\Z \times$ $(\Z/2\Z)^3$, where $q_2$ is a prime such that $2q_2|m_1$. Since $\delta(H)= 4$, by Lemma \ref{relation},  $\tau(H) \leq 5$. Let $\psi$ be a faithful $\alpha$-representation of $H$ of degree $\tau(H)$. 
	Since $H$ is not of symmetric type, $\psi$ is not irreducible. Using Lemma \ref{1dim direct sum in abelian group}, if deg$(\psi) <5$, we can assume $\psi= \oplus_{i=1}^2\psi_i$ such that $\operatorname{deg}(\psi_i)=2$. Let $H_0$ be the set of $\alpha$-regular elements of $H$. Then, $H_0= \Z/2q_2\Z \times$ $(\Z/2\Z)^2$ or $H_0= \Z/q_2\Z \times$ $(\Z/2\Z)^3$. Since rank$(H_0)>2$, proceeding along the same lines as proof of Theorem \ref{elementary abelian p-group}$(i)$, we conclude that $\psi$ cannot be faithful. Hence, $\tau(H) = 5 \leq$ $\tau(G)$. Since $\delta(G)= 4$, $\tg = 5= n+1$ in view of Lemma \ref{relation}. 
	%By similar arguments as above, $\psi$ is not irreducible and $\operatorname{deg}(\psi)>3$. If $\psi= \oplus_{i=1}^2\psi_i$, $\operatorname{deg}(\psi_i)=2$, 
	
	In all other cases, taking $H = (\Z/p\Z)^n$ $\subseteq$ $G$, by Theorem \ref{elementary abelian p-group}, we have $\tau(G)$ $\geq \tau(H) = n + 1$. Since $\delta(G)= n$, Lemma \ref{relation} yields that $\tg = n + 1$. %\poonam{Theorem \ref{elementary abelian p-group} used twice}
	
\end{proof}

\subsection{Extra-special $p$-groups and Heisenberg groups} \label{EH} 
In this section, we study $\tau(G)$ and $\tau_{irr}(G)$ for extra-special $p$-groups and Heisenberg groups. We observe that Theorem \ref{direct_sum} plays a crucial role in these computations.

%Note that, using Theorem \ref{direct_sum}, one can describe $\tau(G)$ for $p$-groups such that $Z(G) \subseteq G^\prime$  \Nayak{Why?} \poonam{We are also using that inf is surjective} . As an application of this result, we study $\tau(G)$ and $\tau_{irr}(G)$ for extra-special $p$-groups and Heisenberg groups. 

We first recall the following part from \cite[Page 772]{HIGGS}, which will be used in the proof of the results presented in this section. 
Let $\alpha$ be a cocycle of an abelian group $G$. Define a map $\alpha' \colon G \times G \to \mathbb{C}^{\times}$ by $\alpha'(x,y) ={\alpha(x,y)}{\alpha(y,x)^{-1}}$ $\forall$ $x,y$ $\in G$. Then, $\alpha'$ is independent of the choice of cocycle representative in $[\alpha]$. Suppose that
$$G \cong \langle x_1 \rangle \times \cdots \times \langle x_n \rangle \cong (\Z/p^k\Z)^n,$$ where each $x_i$ has order $p^k$. Let $\omega$ be a primitive $p^k$th root of unity. Then, for all $i,j \in \{1,\ldots,n\}$, we can write
\[
\alpha'(x_i,x_j)=\omega^{c(i,j)}.
\]
We define the alternating $n \times n$ matrix $C(\alpha)= [c(i,j)]$ over $\Z/p^k\Z$ representing $[\alpha]$. Then, the map $[\alpha] \mapsto C(\alpha)$ is an isomorphism from $\mathrm{H^2}(G, \mathbb{C}^{\times})$ onto the group $R(p^k,n)$ of all alternating $n \times n$ matrices over $\Z/p^k\Z$, under matrix addition.
Let $x= x_1^{b_1}\dots x_n^{b_n}\in G$. Then, $x$ is an $\alpha$-regular element if and only if $C(\alpha)b= 0$, where $b^{\top}= [b_1,...,b_n]$.
Furthermore, if $k=1$, then $\operatorname{rank}\big(C(\alpha)\big) = 2r$,
where $p^r$ is the common degree of all irreducible $\alpha$-representations of $G$. 

\begin{remark} \label{ker(C(alpha))} Note that $|\operatorname{ker}\big(C(\alpha)\big)| = |G_0|$.   
\end{remark}

In the following result, we determine $\tau_{irr}(G)$ and $\tau(G)$ for extra-special $p$-groups of order $p^{2n+1}$, with $n \geq 1$.
Since groups of order $p^3$ will be discussed in Section \ref{groupsoforderp5}, we consider $n> 1$ here.

\begin{theorem} \label{extraspecial $p$-group}
	Let $G$ be an extra-special $p$-group of order $p^{2n+1}$, where $n>1$. Then, $\tau_{irr}(G)$ does not exist and  
	\begin{equation}\nonumber
		\tau(G) =
		\begin{cases}
			2p^{n/2}, & \text{if n is even }  \\
			p^{(n+1)/2}+p^{(n-1)/2},  & \text{if n is odd } 
		\end{cases}
	\end{equation}
\end{theorem}

\begin{proof}
	By \cite[Chapter 4, Theorem 7.1]{karpilovsky}, extra-special $p$-groups are central products of their subgroups of order $p^3$. Hence, by Lemma \ref{central product} and Corollary \ref{non-capable groups}, $\tau_{irr}(G)$ does not exist. %\poonam{Do I need to explain why $H^{\prime} \cap K^{\prime} \neq 1$?}
	
	From the proof of Lemma \ref{central product}, it follows that $\mathrm{inf}: \mathrm{H}^2(G/Z(G), \C^{\times}) \to \mathrm{H}^2(G, \C^{\times})$ is surjective.  %\poonam{Lemma \ref{central product} used twice}
	Let $Z(G) = \langle z \rangle \cong \Z/p\Z$ and $G/Z(G) = \prod_{k=1}^{2n}\langle x_i \rangle \cong (\Z/p\Z)^{2n}$. 
	Recall the transgression map 
	$$\mathrm{tra} : \mathrm{Hom}(Z(G), \mathbb{C}^{\times}) \to \mathrm{H}^2(G/Z(G), \mathbb{C}^{\times})$$ from \cite[Proof of Theorem 1.7]{sabnam}, which is
	defined as follows. For $X = \prod_{k=1}^{2n} x_k^{i_k}$,
	$Y = \prod_{k=1}^{2n} x_k^{j_k} $ in $G/Z(G)$, $$\mathrm{tra}(\chi)(X, Y) = \chi(z)^{-\sum_{k=1}^n j_k i_{n+k}}$$ for  $\chi \in \mathrm{Hom}(Z(G), \mathbb{C}^{\times})$. Let $\omega$ be a primitive $p$th root of unity. Then $\chi(z)= \omega^r$ for some $r\in \{0,1,\dots,p-1\}$.
	Hence, $C(\mathrm{tra}(\chi))= 
	\begin{bmatrix}
		0 & rI_{n \times n} \\
		-rI_{n \times n} & 0
	\end{bmatrix}$.
	It is easy to see that $\text{rank}\big(C(\mathrm{tra}(\chi)\big) = 2n$ for $r\neq 0$. 
	
	Let $\tau(G)= m$ and $\rho_i \in \mathrm{Irr}^\alpha(G)$ such that $\rho = \oplus_{i=1}^l\rho_i$ be a faithful $\alpha$-representation of $G$ of degree $m$. By Theorem \ref{direct_sum}, taking $A= Z(G)$, we have $\tilde\rho_i \in \mathrm{Irr}^{\beta\mathrm{tra}(\chi_i)}\big(G/Z(G)\big)$, where $\chi_i \in  \Hom(Z(G), \mathbb{C}^{\times})$, such that $\oplus_{i=1}^l\chi_i$ is faithful on $Z(G)$. 
	Furthermore, we can assume that $\chi_i \neq \chi_j$ for $i \neq j$.
	Since $\chi_1 \oplus \chi_2$ is faithful on $Z(G)$, we have $l=2$. If $\beta_1= \beta\mathrm{tra}(\chi_1)$, then $\beta\mathrm{tra}(\chi_2)= \beta_1\mathrm{tra}(\chi_3)$, where $\chi_3= \chi_1^{-1}\chi_2 \in \mathrm{Hom}(Z(G), \mathbb{C}^{\times})$. Hence, to compute $\tau(G)$, it is enough to find 
	\[\operatorname{min}\Biggl\{
	\begin{array}{cl} \mathrm{deg}(\tilde{\rho}) \mid & \tilde{\rho}= \tilde{\rho}_1 \oplus \tilde{\rho}_2,\ \tilde{\rho}_1 \in \irr^{\beta_1}(G/Z(G)) \text{ and } \tilde{\rho}_2 \in \mathrm{Irr}^{\beta_1\tra(\chi)}(G/Z(G)),\\
		& \chi\in\operatorname{Hom}(Z(G),\mathbb{C}^{\times}), \chi \neq 1,[\beta_1] \in \mathrm{H}^2\big(G/Z(G), \C^{\times}\big)
	\end{array}
	\Bigg\}\]
	which is equivalent to consider
	\[\operatorname{min}\Biggl\{
	\begin{array}{cl}  \text{rank}\big(C(\beta_1)\big)+\text{rank}\Big(C\big(\beta_1\mathrm{tra}(\chi)\big)\Big)\mid & [\beta_1] \in \mathrm{H}^2(G/Z(G), \C^{\times}),\\
		& \chi \in \operatorname{Hom}(Z(G),\mathbb{C}^{\times}), \chi \neq 1 \end{array}
	\Bigg\}.\] 
	Let $\operatorname{deg}(\rho_1)=\operatorname{deg}(\tilde{\rho}_1)= p^{k_1}$ and $\operatorname{deg}(\rho_2)=\operatorname{deg}(\tilde{\rho}_2)=p^{k_2}$.
	Now, $\text{rank} \big(C(\beta_1)\big) +\text{rank} \Big(C\big(\beta_1\mathrm{tra}(\chi)\big)\Big)=$
	$\text{rank}\big(C(\beta_1)\big)+\text{rank}\Big(C\big(\beta_1\big)+C\big(\mathrm{tra}(\chi)\big)\Big)$
	$\geq \text{rank}\Big(C\big(\mathrm{tra}(\chi)\big)\Big)=2n$. Therefore, $k_1+k_2 \geq n$. It is easy to check that 
	$$\operatorname{min}\{p^{k_1}+p^{k_2}\mid k_1+k_2 \geq n\} = p^{\lfloor{n/2}\rfloor} +p^{\lceil{n/2}\rceil}.$$  
	Taking $C(\beta_1)=
	\begin{bmatrix}
		0 & 
		\begin{bmatrix}
			-L & 0 \\[4pt]
			0 & 0
		\end{bmatrix}
		\\[10pt]
		\begin{bmatrix}
			L & 0 \\[4pt]
			0 & 0
		\end{bmatrix}
		& 0
	\end{bmatrix},
	\quad \text{where} \quad L = I_{\left\lfloor \frac{n}{2} \right\rfloor \times \left\lfloor \frac{n}{2} \right\rfloor}$ and $\chi \in  \operatorname{Hom}(Z(G),\mathbb{C}^{\times})$ given by $\chi(z^i)=\omega^{i}$, we get $\operatorname{deg}(\rho) = p^{\frac{\mathrm{rank}(C(\beta_1))}{2}}+p^{\frac{\mathrm{rank}(C(\beta_1 tra(\chi)))}{2}}
	=p^{\lfloor{n/2}\rfloor} + p^{\lceil{n/2}\rceil}$. Hence, the result follows.
	
\end{proof}
Note that while writing a presentation of a group, we omit relations of the form $[x, y] = 1$ for generators $x, y \in G$.  
Our next aim is to describe $\tau(G)$ and $\tau_{irr}(G)$ for the Heisenberg groups $G = H_{2n+1}(\Z/p^k\Z)$, which admit the presentation: 
$$H_{2n+1}(\Z/p^k\Z)= \langle x_i, y_i, z, 1\leq i \leq n \mid [x_i, y_i] = z, x_i^{p^k} = y_i^{p^k} = z^{p^k} = 1 \rangle.$$
\begin{theorem} \label{Heisenberg group}
	Let $p$ be an odd prime. 
	\begin{enumerate}
		\item Suppose $G$ is a $p$-group such that $G' \subseteq Z(G)$ and $Z(G)$ is cyclic. If $\tau_{irr}(G)$ exists, then $G \cong H_{3}(\Z/p^k\Z)$ and in that case, $\tau_{irr}(G)= \tau(G)= p^k$.  
		
		\item If $G= H_{2n+1}(\Z/p^k\Z)$ with $n > 1$, then $\tau_{irr}(G)$ does not exist and $\tau(G) \leq p^{k\lfloor{n/2}\rfloor} + p^{k\lceil{n/2}\rceil}$.
	\end{enumerate}
	
\end{theorem}

\begin{proof}(1)  %Suppose $G = \langle x, y, z \mid [x, y] = z, x^{p^k} = y^{p^k} = z^{p^k} = 1 \rangle$. 
	From the proof of \cite[Proposition 7.6]{NG}, it follows that if $G' \subseteq Z(G)$, $Z(G)$ is cyclic and $\tau_{irr}(G)$ exists, then $G \cong H_{3}(\Z/p^k\Z)$.
	Now, by \cite[Theorem 1.2]{hatui}, consider the representation group of $G= H_{3}(\Z/p^k\Z)$ given by
	$$G^* = \langle \x, \y, \z, z_1, z_2 \mid [\x, \y] = \z, [\z, \x] = z_1, [\z, \y] = z_2, 
	\x^{p^k} = \y^{p^k} = \z^{p^k} = 1 \rangle. $$
	Then, $G^*/\langle z_1, z_2 \rangle \cong G$. Let $A= \langle z_1, z_2 \rangle$.
	By \cite[Lemma 2.2]{hatui}, any cocycle of $G$ is cohomologous to a cocycle $\alpha$ of the form 
	\[
	\alpha\big(z^{m_1}y^{n_1}x^{q_1}, z^{m_2}y^{n_2}x^{q_2}\big) = \lambda^{(m_2q_1 + n_2 \frac{q_1(q_1-1)}{2})} \mu^{(n_1m_2+q_1 \frac{n_2(n_2-1)}{2}+q_1n_1n_2)} , 
	\]
	for some $\lambda, \mu \in \mathbb{C}^{\times}$ such that $\lambda^{p^k} = \mu^{p^k} = 1$.  
	It is easy to check that $1$ is the only $\alpha$-regular element of $Z(G)$ if and only if either $\lambda$ or $\mu$ is a primitive $p^k$th root of unity. 
	Thus, by Lemma \ref{faithful projective representation}, $G$ admits a faithful irreducible $\alpha$-representation only for those $\alpha$ having $\lambda$ or $\mu$ a primitive $p^k$th root of unity. Let $S$ be the set consisting of all such $[\alpha]$. 
	Therefore, by Theorem \ref{bijective correspondence} we have
	\begin{eqnarray*}
		\tau_{irr}(G)&=&\operatorname{min}\{\operatorname{deg}(\operatorname{\rho})\mid \rho \in \mathrm{Irr}^\alpha(G) \text{ for } [\alpha] \in S\}  \\
		&=&\operatorname{min}\{\operatorname{deg}(\operatorname{\Gamma}) \mid \Gamma \in \mathrm{Irr}(G^*\mid \chi), \chi\in \Hom(A, \C^\times) \text{ such that } \operatorname{tra}(\chi) \in S\}.
	\end{eqnarray*}
	Define a set $T=\{\chi \in \Hom(A,\C^\times) \mid \chi(z_1) \text{ or } \chi(z_2) \text{ is a primitive } p^k\text{th root of unity} \}$.
	First, we show that $\operatorname{tra}(\chi) \in S$ if and only if $\chi \in T$. If $\chi^{p^m}=1$ for some $m<k$, then $\operatorname{tra}(\chi)^{p^m}=1$.   
	%It is easy to observe that, if $\chi(\gamma_1)=\lambda$ and $\chi(\gamma_2) =\delta$, then \[ \operatorname{tra}(\chi) \big((\gamma^{m_1} \alpha^{n_1} \alpha_1^{r_1}), (\gamma^{m_2} \alpha^{n_2} \alpha_1^{r_2})\big) = \delta^{(m_2+p n_2r_1)n_1+ \frac{r_1pn_2(n_2-1)}{2}}\lambda^{r_1m_2},\]
	Hence, if $\operatorname{tra}(\chi) \in S$, then $\chi \in T$. Conversely, let $\chi \in T$ and $\nu: G \to G^*$ be a section defined by $\nu(z^{m}y^{n}x^{q}) = \z^{m}\y^{n}\x^{q}$. Then, 
	$$\operatorname{tra}(\chi)(x, z) = \chi(\nu(x)\nu(z)\nu(xz)^{-1})= \chi(\x\z( \z\x)^{-1})= \chi(\z_1^{-1}) \text { and }$$ 
	$$\operatorname{tra}(\chi)(y, z) = \chi(\nu(y)\nu(z)\nu(yz)^{-1})= \chi(\y\z( \z\y)^{-1})= \chi(\z_2^{-1}).$$ 
	Therefore, $\operatorname{tra}(\chi) \in S$.
	Hence, $$\tau_{irr}(G)= \operatorname{min}\{\operatorname{deg}(\operatorname{\Gamma}) \mid \Gamma \in \mathrm{Irr}(G^*\mid \chi), \chi \in T\}.$$
	Consider the abelian normal subgroup $N = \langle \z, z_1, z_2 \rangle$ of $G^*$ of order $p^{3k}$. Let $\chi \in T$. WLOG, we assume that $\chi(z_1) = \chi(z_2)^a$, where $\chi(z_2)$ is a primitive $p^k$th root of unity and $0 \leq a \leq (p^{k}-1)$. 
	Let $\chi_1 \in \text{Irr}(N)$ such that $\chi_1|_{\langle z_1, z_2 \rangle} = \chi$. It is easy to check that $I_{G^*}(\chi_1)= \langle N,\x^{l}\y^{-al}\rangle$ and $G^*/I_{G^*}(\chi_1)$ is cyclic of order $p^k$. It also follows that, for $\Gamma \in \operatorname{Irr}(G^* \mid \chi)$, deg$(\Gamma) = p^k$ by Theorem \ref{Clifford's theory}. Hence, $\tau_{irr}(G)= p^k$. 
	%Following the idea used in \cite[Section 6.3]{Hatuip2}, $\chi_1$ extends to $I_{G^*}(\chi_1)$ if and only if $a$ is a multiple of $p$, and in this case, Ind$_{I_{G^*}(\chi_1)}^{G^*}(\chi_1)$ has dimension $p^k$. 
	
	%\textbf{Case (b)}
	%Let $\chi'_2 \in \text{Irr}(\langle z_1, z_2 \rangle)$ such that $\chi'_2(z_1) = \xi_2^b$, $\chi'_2(z_2) = \xi_2$ where $\xi_2$ is a primitive $p^k$th root of unity and $0 \leq b  \leq (p^{k}-1)$. Proceeding along similar lines to Case (a), one can see that $\operatorname{Irr}(G^* \mid \chi'_2)$ are of dimension $p^k$.  \\
	
	%Observe that $\operatorname{tra}(\chi)=[\alpha]$, $\alpha\in S$ if and only if $\chi= \chi_i'$, $i=1,2$.
	
	Now we claim that $\tau(G) = p^k$. 
	Since $\tau_{irr}(G)= p^k$, $\tau(G) \leq p^k$. 
	WLOG, we assume that $\chi'(z_1) = \chi'(z_2)^a$, where $\chi'(z_2)$ is a primitive $p^r$th root of unity, $0 <r <k$ and $0 \leq a \leq (p^{r}-1)$.  Let $\chi \in \text{Irr}(N)$ such that $\chi|_{\langle z_1, z_2 \rangle} = \chi'$.
	It is easy to check that $I_{G^*}(\chi) = \langle N,\x^{l}\y^{-al},\y^{p^r} \rangle$. %By \cite[Theorem 3.1]{karpilovsky2} 
	Let $N_1= \langle N,\x^{l}\y^{-al} \rangle$. As $N_1/N$ is cyclic, $\chi$ extends to some $\bar{\chi}\in \Hom(N_1, \C^\times)$ by Theorem \ref{Clifford's theory}. Note that $I_{G^*}(\chi)/N_1 = \langle \y^{p^r}N_1 \rangle$. %\poonam{Theorem \ref{Clifford's theory} used twice}
	
	%We show that if $\alpha= \tra(\chi)$ and $\rho$ is an $\alpha$-representation of $G$ such that $\deg(\rho)< p^k$, then $\rho$ cannot be faithful. 
	
	First, we study the dimensions and behaviour of $\Gamma \in \operatorname{Irr}(G^* \mid \chi)$ as $\Gamma$ is a lift of $\rho \in \irr^\alpha(G)$ for $\alpha= \tra(\chi')$. Let $n$ be the smallest positive integer such that $r \leq n\leq k$ and $\bchi^{(\y^{p^n})}(\x\y^{-a}) = \bchi(\x\y^{-a})$. 
	%\poonam{Is it enough to consider subgroups of this form?} 
	Then, we can check that $\bchi^{(\y^{p^n})}(\x^l\y^{-al}) = \bchi(\x^l\y^{-al})$. Since $\bchi^{(\y^{p^n})}(\z) = \bchi(\z z_2^{-p^n})= \bchi(\z)$, $\bchi$ extends to $N_2= \langle N,\x^l\y^{-al}, \y^{p^n} \rangle =\langle N,\x\y^{-a}, \y^{p^n} \rangle$. Let $\sigma: N_2 \to \C^\times$ be such that $\sigma|_{N_1} = \bchi$. 
	Note that $I_{G^*}(\bar{\chi}) =  N_2$. Hence, by Theorem \ref{Clifford's theory}, every $\psi \in \operatorname{Irr}(I_{G^*}(\chi) \mid \bar{\chi})$ is of the form $\psi= \operatorname{Ind}^{I_{G^*}(\chi)}_{N_2}(\sigma)$ and $\Gamma = \operatorname{Ind}^{G^*}_{{I_{G^*}(\chi)}}(\psi)$. Hence, $\operatorname{deg}(\Gamma)= p^n$. Note that $\{\y^i$, $i=1,2,..,p^r\}$ and $\{\y^{ip^r}$, $i=1,2,..,p^{n-r}\}$ are sets of left coset representatives of $I_{{G}^*}(\chi)$ in $G^*$ and of $N_2$ in $I_{{G}^*}(\chi)$, respectively. Then, 
	$$\psi(\z)= \left[\begin{array}{cccc}
		\chi(\z)\chi(z_2)^{-p^r} & & & \text{\huge0} \\
		& \chi(\z)\chi(z_2)^{-2p^r} & & \\
		& & \ddots & \\
		\text{\huge0} & & & \chi(\z)\chi(z_2)^{-p^n} 
	\end{array}\right]_{p^{n-r} \times p^{n-r}}= \chi(\z)I_{p^{n-r} \times p^{n-r}}$$ and  $$\Gamma(\z)= \left[\begin{array}{ccc}
		\psi(\z)\psi(z_2)& & \text{\huge0} \\
		& \ddots & \\
		\text{\huge0} &  & \psi(\z)\psi(z_2^{p^r})
	\end{array}\right]_{p^n \times p^n} = \chi(\z)\left[\begin{array}{ccc}
		\psi(z_2)& & \text{\huge0} \\
		& \ddots & \\
		\text{\huge0} &  & \psi(z_2^{p^r})
	\end{array}\right]_{p^n \times p^n}.$$ 
	
	\bigskip
	
	Now, we show that if $\alpha= \tra(\chi')$ and $\rho$ is an $\alpha$-representation of $G$ with $\deg(\rho)< p^k$, then $\rho$ cannot be faithful.
	Let $\rho = \oplus_{i=1}^m \rho_i$, where $\rho_i \in \mathrm{Irr}^\alpha(G)$ and $\deg(\rho_i)<p^k$. Let $\Gamma_i \in \irr(G^* \mid \chi')$ correspond to $\rho_i$, where the correspondence is given by Theorem \ref{bijective correspondence}. Then, $\Gamma_i = \operatorname{Ind}^{G^*}_{{I_{G^*}(\chi_i)}}(\psi_i)$, where $\chi_i \in \text{Irr}(N)$ 
	such that $\chi_i |_A= \chi'$, and $\psi_i=$ $\operatorname{Ind}^{I_{G^*}(\chi_i)}_{N_2}(\sigma_i)$ for $\sigma_i \in \Hom(N_2, \C^\times)$ such that $\sigma_i|_{N_1} = \chi_i$. Let $r \leq n_i < k$ be the smallest positive integer such that $\chi_i^{(\y^{p^{n_i}})}(\x\y^{-a}) = \chi_i(\x\y^{-a})$. By viewing the above matrices, $\rho_i(z^{p^{k-1}})= \Gamma_i(\z^{p^{k-1}})= \chi_i(\z^{p^{k-1}})$. If none of $\chi_i(z)$ is a primitive $p^k$th root of unity, then $\rho(z^{p^{k-1}})= I$ and $\rho$ is not faithful.
	Hence, for $\rho|_{Z(G)}$ to be faithful, at least one of the $\chi_i(\z)$ has to be a primitive $p^k$th root of unity.
	Now, for any $n \geq r$, we have $\chi_i^{(\y^{p^{n}})}(\x\y^{-a}) = \chi_i(\x\y^{-a}\z^{-p^n}z_1^{-p^n\binom{1}{2}}z_2^{ap^n+\binom{p^k-p^n}{2}})=\chi_i(\x\y^{-a} \z^{-p^n})$.
	Then, $\chi_i^{(\y^{p^{n_i}})}(\x\y^{-a}) \neq \chi_i(\x\y^{-a})$, which is a contradiction. Hence, $\rho|_{Z(G)}$ is not faithful. In view of Theorem \ref{restriction}, $\rho$ cannot be faithful. 
	
	%Next, let $\chi' \in \text{Irr}(\langle z_1, z_2 \rangle)$ such that $\chi'(z_1) = \xi_4^b$, $\chi'(z_2) = \xi_4$ where $\xi_4$ is a primitive $p^r$th root of unity, $0 < r< k$ and $0 \leq b \leq (p^{r}-1)$. Let $\alpha=\tra(\chi')$ and $\rho \in \irr^\alpha(G)$ such that $\deg(\rho) < p^k$. Then, by similar arguments as above, $\rho$ cannot be faithful. 
	
	Now, let $\chi' \in \text{Irr}(\langle z_1, z_2 \rangle)$ such that $\chi'(z_1) = \chi'(z_2) = 1$. Then, $\tra(\chi')= 1$. Let $\rho\in \irr(G)$ be faithful. Since $\dg= p^k$ (\cite[Theorem 3.1, (8)]{kaur}), $\deg(\rho) \geq p^k$. This completes the proof. \\
	
	%The following two cases will occur.\\
	
	%We use Clifford's theory here. Observe that for any $n \geq r$, $\chi^{(\y^{p^n})}(\x\y^{-a}) = \chi(\x\y^{-a}\z^{-p^n}z_1^{-p^n\binom{1}{2}}z_2^{ap^n+\binom{-p^n}{2}})=\chi(\x\y^{-a} \z^{-p^n})$. Let $\Gamma \in \operatorname{Irr}(G^* \mid \chi)$. Now the following cases occur. \\
	
	% By \cite[Section 6.3]{Hatuip2}, $\chi$ extends to $N_1$ if and only if $a$ is multiple of $p$, let $\chi$ also denote the extension to $N_1$. 
	
	%\textbf{Case (a)} If $\bar{\chi}^{(g)}(\x\y^{-a})\neq \bchi(\x\y^{-a})$ $\forall$ $g \in I_{G^*}(\chi)\backslash N_1$, then $I_{G^*}(\bar{\chi}) =  N_1$. Hence, by Theorem \ref{Clifford's theory}, $\Gamma = \operatorname{Ind}^{G^*}_{{I_{G^*}(\chi)}}(\psi_1)$ and $\operatorname{deg}(\Gamma)= p^k$. \poonam{Should I include case a in case b, i.e, n=k?}
	%by Mackey's irreducibility criterion (\cite[Theorem 8.3.6]{steinberg}), $\psi_1=$ $\operatorname{Ind}^{I_{G^*}(\chi)}_{N_1}(\chi)$ is irreducible.  \\
	
	%\textbf{Case (b)} 
	
	(b) For $n > 1$, the group $H_{2n+1}(\Z/p^k\Z)$ is a central product of $H_{2n-1}(\Z/p^k\Z)$ and $H_{3}(\Z/p^k\Z)$. Therefore, by Lemma \ref{central product} and Corollary \ref{non-capable groups}, $\tau_{irr}(G)$ does not exist, and the proof of the Lemma \ref{central product} also shows that $\mathrm{inf}: \mathrm{H}^2(G/Z(G), \C^{\times}) \to \mathrm{H}^2(G, \C^{\times})$ is surjective. Since $Z(G) = \langle z \rangle \cong \Z/p^k\Z$,  let $G/Z(G) = \prod_{k=1}^{2n}\langle a_i \rangle \cong (\Z/p^k\Z)^{2n}$. Using arguments similar to those in \cite[Proof of Theorem 1.7]{sabnam}, one can check that the transgression map $\mathrm{tra} : \mathrm{Hom}(Z(G), \mathbb{C}^{\times}) \to \mathrm{H}^2(G/Z(G), \mathbb{C}^{\times})$ is defined as follows. For $X = \prod_{k=1}^{2n} a_k^{i_k}$, $Y = \prod_{k=1}^{2n} a_k^{j_k} $ in $G/Z(G)$, $$\mathrm{tra}(\chi)(X, Y) = \chi(z)^{-\sum_{k=1}^n j_k i_{n+k}}$$ for $\chi \in \mathrm{Hom}(Z(G), \mathbb{C}^{\times})$. Let $\omega$ be a primitive $p^k$th root of unity. Then, $\chi(z)= \omega^r$ for some $r\in \{0,1,\dots,p^k-1\}$.
	Hence, $C\big(\mathrm{tra}(\chi)\big)= 
	\begin{bmatrix}
		0 & rI_{n \times n} \\
		-rI_{n \times n} & 0
	\end{bmatrix}$.
	
	Let $\tau(G)= m$ and $\rho$ be a faithful $\alpha$-representation of $G$ of degree $m$. Then, $\rho= \oplus_{i=1}^l\rho_i$, where $\rho_i \in \irr^{\alpha}(G)$. By Theorem \ref{direct_sum} (taking $A= Z(G)$), $\tilde\rho_i \in \mathrm{Irr}^{\beta\mathrm{tra}(\chi_i)}\big(G/Z(G)\big)$, where $\chi_i \in  \Hom(Z(G), \mathbb{C}^{\times})$ such that $\oplus_{i=1}^l\chi_i$ is faithful on $Z(G)$. 
	Furthermore, we can assume that $\chi_i \neq \chi_j$ for $i \neq j$. If $\chi_i(z)$ is not a primitive $p^k$th root of unity for any $i$, then $\oplus_{i=1}^l\chi_i$ cannot be faithful on $Z(G)$. Hence, WLOG, let $\chi_1(z)$ be a primitive $p^k$th root of unity. Since $\chi_1 \oplus \chi_2$ is faithful on $Z(G)$, we have $l=2$. 
	Hence, it is enough to find 
	
	\[\operatorname{min}\Biggl\{
	\begin{array}{cl}  \mathrm{deg}(\tilde{\rho}) \mid & \tilde{\rho}=\tilde{\rho}_1 \oplus \tilde{\rho}_2,\ \tilde{\rho}_1 \in \irr^{\beta\tra(\chi_1)}(G/Z(G)) \text{ and }\tilde{\rho}_2 \in \mathrm{Irr}^{\beta \tra(\chi_2)}(G/Z(G)), \\
		& \chi_i \in  \operatorname{Hom}\big(Z(G),\mathbb{C}^{\times}\big), i=1,2 \text{ and } \chi_1(z) \text{ is a primitive } p^k\text{th root of unity} 
	\end{array}
	\Bigg\}\] 
	\iffalse
	By Remark \ref{ker(C(alpha))} and Lemma \ref{properties of abelian groups}, this is equivalent to consider
	\[\operatorname{max}\Biggl\{
	\begin{array}{cl}  & \Big|\text{ker}\big(C(\beta\tra(\chi_1)\big)\Big|+\Big|\text{ker}\Big(C\big(\beta\mathrm{tra}(\chi_2)\big)\Big)\Big|,  \beta \in Z^2(G/Z(G), \C^{\times}),   \\
		& \chi_i \in  \operatorname{Hom}(Z(G),\mathbb{C}^{\times}), i=1,2 \text{ and } \chi_1(z) \text{ is a primitive } p^k\text{th root of unity}  \end{array}
	\Bigg\}\] 
	\fi
	
	Let $C(\beta)=
	\begin{bmatrix}
		0 & 
		\begin{bmatrix}
			-L & 0 \\[4pt]
			0 & 0
		\end{bmatrix}
		\\[10pt]
		\begin{bmatrix}
			L & 0 \\[4pt]
			0 & 0
		\end{bmatrix}
		& 0
	\end{bmatrix}
	\quad \text{, where} \quad L = I_{\left\lfloor \frac{n}{2} \right\rfloor \times \left\lfloor \frac{n}{2} \right\rfloor}$. Let $\chi \in  \operatorname{Hom}(Z(G),\mathbb{C}^{\times})$ be defined by $\chi(z^i)= \omega^{i}$. Let $\tilde{\rho_1}$ and $\tilde{\rho_2}$ be irreducible $\beta$ and $\beta\tra(\chi)$-representations of $G/Z(G)$. Then, by Remark \ref{ker(C(alpha))} and Lemma \ref{properties of abelian groups}, $\deg(\tilde{\rho_1}) = p^{k{\lfloor{n/2}\rfloor}}$ and $\deg(\tilde{\rho_2})= p^{k{\lceil{n/2}\rceil}}$. 
	Let $\alpha(g,h) = \beta(gZ(G), hZ(G))$. Then, $\alpha \in Z^2(G, \C^\times)$. Let $\rho_i \in \irr^\alpha(G)$ correspond to $\tilde{\rho_i}$, $i=1,2$, where the correspondence is given by Remark \ref{remark for bijective correspondence via inflation}. Since $1|_{Z(G)} \oplus \chi$ is faithful on $Z(G)$, Theorem \ref{direct_sum} implies that $\rho_1 \oplus \rho_2$ is faithful on $G$. Hence, $\tg \leq p^{k{\lfloor{n/2}\rfloor}} + p^{k{\lceil{n/2}\rceil}}$. 
\end{proof}
\section{Examples} \label{additional} 
%${Computation of $\tau(G)$ and $\tau_{irr}(G)$ for further groups}  

%We first compute $\tau(G)$ and $\tau_{irr}(G)$ for $p$-groups with trivial Schur multiplier and cyclic centre of order $p$.

In this section, we discuss several examples illustrating $\tau(G)$ and $\tau_{irr}(G)$ for certain classes of groups $G$, particularly metacyclic groups, finite simple groups and symmetric groups. 
We begin by recalling the following facts, which will be useful in discussing the examples.
\begin{enumerate}
	\item  If $G= D_{2n}$ or $Q_{4n}$, then $\delta(G)= 2$ by \cite[Example 2.2]{kaur}, where $D_{2n}$ and $Q_{4n}$ denote the Dihedral group of order $2n$ and the Quaternion group of order $4n$, respectively.
	\item If $G= S_n, n \geq 5$ or $A_n, n \geq 6$, then $\delta(G) = \delta_{irr}(G)= n-1$ (\cite[Example 2.4, 2.5, 2.6]{kaur}).
	\item  For $n \geq 4$, $M(S_n) \cong \Z/2\Z$ by \cite[Chapter 4, Theorem 8.3(i)]{karpilovsky}. 
	\item For $[\alpha] \neq 1$, every $\alpha$-representation has even degree by \cite[Chapter 5, Theorem 5.3]{karpilovsky3}. 
	
	\item A metacyclic $p$-group $G$, $p \ge 3$, has a faithful irreducible projective representation if and only if $G$ is of symmetric type, i.e., $G$ can be written as a semidirect product of two isomorphic cyclic subgroups  (see \cite [Theorem 5.2]{NG}).
\end{enumerate}

\textbf{Example 5.1.} \label{tau(G) = 2}
By \cite[Section 1.1]{Dolgachev}, the finite subgroups of $\mathrm{PGL}_2(\C)$ are isomorphic to $\Z/n\Z$, $D_{2n}$, $A_4$, $A_5$ or $S_4$. Hence, $\tg = 2$ if and only if $G$ is one of these groups. 
%Therefore     $\tg= = \ti=2$ if and only if 
From Lemma \ref{one dim direct sum in non-abelian group}(1), it follows that if $G$ is non-abelian and $\tg = 2$, then $\ti= 2$. Also, if $\ti= 2$, then $\tg= 2$ and so $G$ is one of the groups $\Z/n\Z$, $D_{2n}$, $A_4$, $A_5$ or $S_4$. Since $\Z/n\Z$ is not of symmetric type, $\tau_{irr}(\Z/n\Z)$ does not exist. Hence, $\ti =2$ if and only if $G$ is one of the groups $D_{2n}$, $A_4$, $A_5$ or $S_4$. \\

\textbf{Example 5.2.} \label{tau(G) = 3}
Let $G$ be a group with $\dg= 2$ such that $G \ncong D_{2n}$. By Lemma \ref{relation}, $\tg \leq 3$. From the character table of $S_4$(\cite[18.1]{liebeck}) and by the proof of Lemma \ref{one dim direct sum in non-abelian group}(1), it follows that $\delta(S_4)=3$. Now,  it follows from Example 5.1 and \cite[Examples 2.2, 2.3 and 2.5]{kaur} that $\tg= 3$.  \\

\textbf{Example 5.3.} 
Let $G$ be a metacyclic $p$-group ($p \ge 3$) of symmetric type. Then, $\tau_{irr}(G)= \sqrt{|G|}$ in view of \cite[Chapter 10, Proof of Theorem 4.16]{karpilovsky}.  \\

\textbf{Example 5.4.}  \label{Q4n}
By Example 5.2, $\tau(Q_{4n})=3$. Moreover, \cite[Chapter 10, Theorem 4.14]{karpilovsky} implies that $\tau_{irr}(Q_{4n})$ does not exist.     \\

\textbf{Example 5.5.} \label{simple group}
Let $G$ be a non-abelian finite simple group.
If $\rho$ is a projective representation of $G$ which is not a direct sum of trivial $1$-dimensional representations, then $\rho$ is faithful. Hence, $\tau(G) = \tau_{irr}(G)$ = minimal degree of an irreducible projective representation greater than $1$.   \\

\textbf{Example 5.6.} \label{tau(G) leq delta(G)}
For $G=S_n$ ($n \geq 3$) or $A_n$ ($n\geq 4$), it follows from Lemma \ref{one dim direct sum in non-abelian group}$(3)$ that $\tau(G) \leq \delta(G)$. \\

\textbf{Example 5.7.} \label{S5}
Consider $G = S_5$. Let $\tau(G)= m$ and $\rho$ be a faithful projective representation of degree $m$. If $\rho$ is an ordinary representation, then $\deg(\rho) \geq \delta(G)= 4$. From the spin character table of $G$ (see \cite[Chapter 5, Theorem 6.8]{karpilovsky3}), it follows that $\deg(\rho) \geq 4$. Hence, $\tau(G) = 4$, in view of Example 5.6. By Lemma \ref{one dim direct sum in non-abelian group}$(2)$, $\tau_{irr}(G) \leq 4$. Again, using the spin character table of $G$, we conclude that $\tau_{irr}(G)= 4$.  \\

\textbf{Example 5.8.} 
Consider $G= S_6$. By Lemma \ref{one dim direct sum in non-abelian group}$(2)$, $\tau_{irr}(G) \leq 5$. Let $\psi$ be a faithful irreducible $\alpha$-representation of $G$ of degree $\tau_{irr}(G)$. If $\alpha= 1$, since $\di=5$, $\deg(\psi) \geq 5$. Let $\rho_6$ be the basic spin representation of $S_6^*$ (See \cite[5.5.B.]{karpilovsky3} for definition). It follows from \cite[Chapter 5, Corollary 5.8, Theorem 5.12]{karpilovsky3} that as a projective representation $\rho_6$ is irreducible and faithful on $S_6$ with $\deg(\rho_6)= 4$. Since any $\alpha$-representation has even degree for $[\alpha] \neq 1$, Example 5.1 yields that $\tau_{irr}(G)= 4$. Now, we have $\tau(G) \leq \tau_{irr}(G)= 4$. If $\rho$ is an ordinary representation of degree $m$ which is faithful as a projective representation, then $m \geq \delta(G)= 5$. Therefore, Example 5.1 implies that $\tau(G)= 4$. \\ % \poonam{Example 5.1 used twice} \\

\textbf{Example 5.9} 
$\tau(A_6)= \tau_{irr}(A_6)= 3$ follows from \cite[Table 2]{gonzalo} using Example 5.1 and Example 5.5.     \\ 

\textbf{Example 5.10.} 
Consider $G = A_7$. Since $S_5 \subseteq A_7$, $\tau(G) \geq 4$ in view of Example 5.7. Thus, $\tau(G) = \tau_{irr}(G) =4$ by virtue of \cite[Table 3]{gonzalo} and Example 5.5.    \\

\textbf{Example 5.11.} 
Consider $G = A_8 \cong \mathrm{PSL}_4(\F_2)$. By \cite[Lemma 5.2]{LANDAZURI}, $\tau(G) \geq 7$. Thus, $\tau(G) = \tau_{irr}(G) = 7$ due to Example 5.6 and Example 5.5.    \\

\textbf{Example 5.12.} 
Consider $G = A_9$. Since $A_8 \subset A_9$, $\tau(G) \geq 7$. Using Example 5.6, we obtain that $\tau(G) \leq 8$. Let $\rho$ be a faithful projective representation of degree $\tau(G)=m$. If $\deg(\rho)=7$, then $\rho$ has to be an ordinary representation. But $\delta(G) = 8$ and so $\tau(G) > 7$. Hence, it follows from Example 5.5 that $\tau(G) = \tau_{irr}(G)= 8$.   \\

% The results for $A_4$ and $A_5$ are derived using (2) (see \cite[5.6]{karpilovsky3} for their projective character tables).
%Let $G =S_4$. From the spin character table of $S_4$ (see \cite[5.6]{karpilovsky3}), we obtain that $\tau(G) = \tau_{irr}(G)= 2$. \\

%By similar arguments, from Table 3 of \cite{gonzalo}, $2< \tau(G) = \tau_{irr}(G) \leq 4$. Let $\rho$ be a faithful projective representation of $G$ of degree $\tau(G)$. If $\rho$ is an ordinary representation, then $\deg(\rho) \geq \delta(G)=6$. Hence, $\tau(G)$ cannot be 3, and the result follows.
%It is easy to see that if $\tau_{irr}(G)=2$, then $\tg = 2$. 

\iffalse
$(b)$ This  follows from the proof of \cite[Chapter 10, Theorem 4.14]{karpilovsky}. We present an alternate proof here. 

If $n$ is odd, it follows from \cite[Example 2.2]{kaur} and Lemma \ref{one dim direct sum in non-abelian group}$(1)$ that $\delta(G)= \delta_{irr}(G)= 2$. From Lemma \ref{one dim direct sum in non-abelian group}$(2)$, it follows that $\tau_{irr}(G)= 2$. 

If $n$ is even, let consider the irreducible ordinary representation $\psi: D_{4n} \to GL_2(\C)$ which is defined by
$\psi(a)=
\left[\begin{array}{ll}cos(\frac{\pi}{n}) &  -sin(\frac{\pi}{n})\\ sin(\frac{\pi}{n})& cos(\frac{\pi}{n})\end{array}\right]$, 
$\psi(b)= \left[\begin{array}{ll}1 &  0 \\ 0 & -1 \end{array}\right]$. 
Then, $D_{4n}/K_{\psi} \cong D_{2n}$. Let $T$ be a transversal for $K_{\psi}$ in $D_{4n}$ containing $1$, and consider $\psi^*:D_{2n} \to GL_2(\C)$  defined by $\psi^*\big(t K_{\psi})= \psi(t)$. It follows from \cite[Page 303]{karpilovsky} that $\psi^*$ is faithful on $D_{2n}$. Hence, the result follows. \\
\fi
%, $Q_{n}$ does not have a faithful irreducible projective representation. Hence, $\tau(Q_8)>2$, due to Remark \ref{some results on p-groups}. Since $\delta(Q_n)=2$ \cite[Example 2.2]{kaur}, therefore, by Lemma \ref{relation}, we have, $\tau(Q_8) \leq 3$. Hence, the result follows.\\
The values of $\ti$ and $\tg$ obtained in the above examples are summarized in the following table. If $\tau_{irr}(G)$ does not exist for a group $G$, then we denote this by writing $\tau_{irr}(G)=$ -.  

\begin{table}[H]
	\centering
	\small
	\begin{tabular}{|c|c|c|}
		\hline
		Group & $\tau(G)$ & $\tau_{irr}(G)$ \\
		\hline
		$D_{2n}$ & $2$ & $2$ \\
		$Q_{4n}$ & $3$ & - \\
		$A_4$ & $2$ & $2$ \\
		$A_5$ & $2$ & $2$ \\
		$A_6$ & $3$ & $3$ \\
		$A_7$ & $4$ & $4$ \\
		$A_8$ & $7$ & $7$ \\
		$A_9$ & $8$ & $8$ \\
		$S_4$ & $2$ & $2$ \\
		$S_5$ & $4$ & $4$ \\
		$S_6$ & $4$ & $4$ \\
		\hline
	\end{tabular}
	\caption{Some examples} %for which $\tau_{irr}(G)$ does not exist}
	\end{table}
	
	\section{$\tau(G)$ and $\tau_{irr}(G)$ for non-abelian groups of order $p^n$, $3\leq n \leq 5$} \label{groupsoforderp5}
	%\Nayak{Check that if $\tau_{irr}=p^2$, all irr has degree $p^2$ for groups of order $p^5$.}
	In this section, we compute $\tau(G)$ and $\tau_{irr}(G)$ for non-abelian $p$-groups of order $p^3$, $p^4$ for all primes $p$, and for groups of order $p^5$ for primes $p\geq 5$. 
	\iffalse
	Recall that $G^*$ is a representation group of $G$ if $\exists$ a subgroup $A \subseteq {G^*}' \cap Z(G^*)$ such that $A \cong M(G)$ and $G^*/A \cong G$. With these notations, we state our next theorem. 
	
	\begin{lemma}
Let $G$ be a non-abelian $p$-group such that $Z
(G) \subseteq Z(G^*)/A$. Then, the following hold:
\begin{enumerate}
	\item $\ti$ does not exist.
	\item $\tg=\operatorname{min}\Biggl\{
	\begin{array}{cl} \mathrm{deg}(\Gamma_1 \oplus \Gamma_2)  \mid & \Gamma_i\in \irr(G^*\mid \chi) \text{ for some } \chi \in \Hom(A,\C^*), \Gamma_i |_{Z(G)}=\tau_i I, \\
		&  \text{ such that } \tau_1 \oplus \tau_2  \text{ faithful on } Z(G).
	\end{array}\Bigg\}$
	
\end{enumerate}

\end{lemma}

\begin{proof}
\begin{enumerate}
	\item This is a consequence of \cite[Theorem 3.1]{karpilovsky2}.
	\item 
\end{enumerate}    
\end{proof}
\fi
% As we will use the classification of $p$-groups of order $\leq p^{5}$, which is based on the isoclinism of two groups, we start with the following definition.
The next theorem derives the minimum number of elements in $\mathrm{Irr}(G)$ such that their direct sum yields a faithful projective representation of $G$. 
\begin{lemma} \label{direct sum of k irreducible ordinary representations}
Let $G$ be a non-abelian $p$-group such that the rank of $Z(G)$ is $r$. Suppose $\rho_i \in \mathrm{Irr}(G)$ are such that $\oplus_{i=1}^k \rho_i$ is a faithful projective representation of $G$. 
\begin{enumerate}
	\item[(i)] If $Z(G) \not \cong (\Z/2\Z)^2, (\Z/2\Z)^4$, then $k >r$.
	\item[(ii)] If $Z(G) \cong (\Z/2\Z)^2, (\Z/2\Z)^4$, then $k >r-1$.
\end{enumerate}

\end{lemma}

\begin{proof}
(i) Let $\rho = \oplus_{i=1}^k\rho_i$, where $k \leq r$. Suppose $Z(G) \not \cong (\Z/2\Z)^2, (\Z/2\Z)^4$. Let $\chi_i \in \Hom(Z(G), \mathbb C^\times)$ such that $\rho_i|_{Z(G)}= \chi_i I_{n_i\times n_i}$. Then, $\rho|_{Z(G)}= \oplus_{i=1}^k\chi_i I_{n_i\times n_i}$.
Hence, if $\rho$ is faithful on $G$, then by Lemma \ref{restriction}$(2)$, $\oplus_{i=1}^k\chi_i$ is a faithful projective representation of $Z(G)$. Since $\tau(Z(G)) = r+1$ (by Theorem \ref{elementary abelian p-group}), we must have $k > r$. 

(ii) By similar arguments, it can be shown that if $Z(G) \cong (\Z/2\Z)^2, (\Z/2\Z)^4$, then $k >r-1$. This completes the proof. 

\end{proof}
We now state the next result from \cite[Theorem 3.1, Theorem 3.3]{kaur}.
\begin{theorem} \label{embedding degree of a p group with cyclic centre}
Let $G$ be a non-abelian $p$-group with cyclic centre. If $cd(G) =\{1,p^a\}$ or $\{1,p,p^a\}$ for $a \geq 1$, then $\delta(G)= \delta_{irr}(G)= p^a$. 
%(i) $\delta(G)=p$ if and only if, $cd(G) =\{1,p\}$. \\(ii) 
% \poonam{a can also be 1}
\end{theorem}
Recall that $A.G$ is a covering of $G$ if $\exists$ a subgroup $A \subseteq {(A.G)}' \cap Z(A.G) \cap M(G)$ such that $(A.G)/A \cong G$. %If $\psi$ is an ordinary representation of $G$, let $Z(\psi)= {\psi^{-1}(\lambda I), \lambda \in \C^\times}$.

\begin{theorem} 
\label{projective embedding degrees}
Let $B_G= \{A.G \mid A \cong Z(A.G) \text{ and } \delta_{irr}(A.G) \text{ exists} \}$. Then $\tau_{irr}(G)$ exists if and only if $B_G \neq \emptyset$, in which case $\tau_{irr}(G) = \min\{\delta_{irr}(A.G) \mid A.G \in B_G\}$. 
Furthermore, $\tau(G) \geq \min\{\delta(A.G) \} $.
\end{theorem}

\begin{proof}
Let $\rho$ be a faithful projective representation of $G$ and $\rho^*$ be a lift of $\rho$ in $G^*$. Suppose $H= G^*/K_{\rho^*}$. It is easy to see that $\rho^*$ induces a faithful ordinary representation of $H$.
Suppose $\tau_{irr}(G)$ exists. By Theorem \ref{projective embedding degrees}, $B_G \neq \emptyset$. Let $H= A.G \in B_G$. Since $\delta_{irr}(H)$ exists, $Z(H)$ is cyclic due to \cite[Theorem 2.32(a)]{isaacs}. Since $A \subseteq \mathrm{H}^2(G, \C^{\times})$ and $A = Z(H)$, we have $A \cong \Z/p\Z$. Hence, $|H|= p^{k+1}$.

Conversely, let $H$ be a covering of $G$ of order $p^{k+1}$ such that $G \cong H/Z(H)$. Since $Z(H)$ is cyclic, by \cite[Theorem 2.32(b)]{isaacs}, $\delta_{irr}(H)$ exists. Hence, $H \in B_G$. By Theorem \ref{projective embedding degrees}, $\tau_{irr}(G)$ exists.
From the above argument and Theorem \ref{projective embedding degrees}, it follows that $\tau_{irr}(G) = \min\{\delta_{irr}(H) \mid |H|= p^{k+1} \text{ s.t. } Z(H) \cong \Z/p\Z \}$. 

Now, let $\eta$ be a faithful projective representation of $G$ of degree $\tau(G)$ and $\eta^*$ be a lift of $\eta$ in $G^*$. Suppose $H= G^*/K_{\eta^*}$. It is easy to see that $\eta^*$ induces a faithful ordinary representation of $H$. By \cite[Proposition 2.15]{gonzalo}, $H$ is a covering of $G$. Hence, $\tau(G) \geq \min\{\delta(A.G) \}$. 
\iffalse
Let $\rho$ be a faithful projective representation of $G$ and $\rho^*$ be a lift of $\rho$ in $G^*$. Suppose $H= G^*/K_{\rho^*}$. It is easy to see that $\rho^*$ induces a faithful ordinary representation of $H$. 
Suppose $\tau_{irr}(G)$ exists and $\rho$ is irreducible. By \cite[Chapter 4, proof of Theorem 3.1]{karpilovsky2}, $G \cong H/Z(H)$ and $H$ is a covering of $G$. Since $\rho^*$ is irreducible, $H \in B_G$. Conversely, let $B_G \neq \emptyset$. Then, by Theorem \ref{capable groups}, $\tau_{irr}(G)$ exists and $\tau_{irr}(G) \leq \min\{\delta_{irr}(A.G) \mid A.G \in B_G\}$. Also $\tau_{irr}(G) \geq \min\{\delta_{irr}(A.G) \mid A.G \in B_G\}$ by the above arguments. Thus, $\tau_{irr}(G) =\min\{\delta_{irr}(A.G) \mid A.G \in B_G\}$. 
Now, let $\rho$ be of degree $\tau(G)$. By \cite[Proposition 2.15]{gonzalo}, $H$ is a covering of $G$. Hence, $\tau(G) \geq \min\{\delta(A.G) \}$. 
\fi


%\poonam{Theorem \ref{capable groups} used twice}
%From \cite[Proposition 2.15]{gonzalo}, it follows that
%Let $A.G$ be a covering of $G$ and $\psi$ be a faithful ordinary representation of $A.G$ of degree $\delta(A.G)=m$. Suppose $A \cong \ker(\psi)$ and $L$ is a transversal for $\ker(\psi)$ in $A.G$ containing $1$ and let $\psi': (A.G)/A \to \GL_m(\C)$ be defined by $\psi'(xA)= \psi(x)$ ($x \in L)$. It follows from \cite[Page 303]{karpilovsky} that $\psi'$ is a faithful projective representation of $G$ of degree $m$. Hence, the result follows.
\end{proof}

%Using the notations from the proof of the above theorem, we state our next corollary. 
\begin{corollary} \label{Elementary abelian Schur multiplier}
(i) Let $G$ be a non-abelian group of order $p^k$ and $\mathrm{H}^2(G, \C^{\times})$ be an elementary abelian $p$-group. Then, $\tau_{irr}(G)$ exists if and only if there exists a covering $H$ of order $p^{k+1}$ such that $G \cong H/Z(H)$. Furthermore, $\tau_{irr}(G) = \min\{\delta_{irr}(H) \mid |H|=p^{k+1} \text{ s.t. } Z(H) \cong \Z/p\Z \}$.\\
(ii) Let $G$ be a non-abelian $p$-group such that $cd(G^*)= \{1, p^n\}$. If $\tau_{irr}(G)$ exists, then $\tau(G)= \tau_{irr}(G)= p^n$. 
%Then the following hold:
%\begin{enumerate}[label={(\alph*)}]
%\item \item If $\tau_{irr}(G)$ does not exist, then $\tau(G)= p^k+1$.
%\end{enumerate}
%(iii) Let $G$ be a non-abelian group such that $M(G) \cong \Z/p\Z$. If $Z(G) \neq 1$ , then $\tau_{irr}(G)= \delta_{irr}(G^*)$ if $Z(G^*)$ is cyclic and $\tau_{irr}(G)$ does not exist if $Z(G^*)$ is not cyclic.
%\begin{enumerate}[label={(\alph*)}]
%\item If $Z(G) = 1$, then $\tau_{irr}(G)= \min\{\delta_{irr}(G), \delta_{irr}(G^*)\}$. 
%\end{enumerate}    
\end{corollary}

\begin{proof}
$(i)$ Suppose $\tau_{irr}(G)$ exists. By Theorem \ref{projective embedding degrees}, $B_G \neq \emptyset$. Let $H= A.G \in B_G$. Since $\delta_{irr}(H)$ exists, $Z(H)$ is cyclic due to \cite[Theorem 2.32(a)]{isaacs}. Since $A \subseteq \mathrm{H}^2(G, \C^{\times})$ and $A = Z(H)$, we have $A \cong \Z/p\Z$. Hence, $|H|= p^{k+1}$.

Conversely, let $H$ be a covering of $G$ of order $p^{k+1}$ such that $G \cong H/Z(H)$. Since $Z(H)$ is cyclic, by \cite[Theorem 2.32(b)]{isaacs}, $\delta_{irr}(H)$ exists. Hence, $H \in B_G$. By Theorem \ref{projective embedding degrees}, $\tau_{irr}(G)$ exists and also $\tau_{irr}(G) = \min\{\delta_{irr}(H) \mid |H|= p^{k+1} \text{ s.t. } Z(H) \cong \Z/p\Z \}$ by the above arguments. %\poonam{Theorem \ref{projective embedding degrees} used twice}

$(ii)$ Let $A.G$ be any covering of $G$. If $p^m \in cd(A.G)$, let $\Gamma \in \irr(A.G)$ be of degree $p^m$. Then, there is a projective representation $\rho$ of $G$ of degree $p^m$ which lifts to $\Gamma$. Since $cd(G^*)=\{1, p^n\}$, by Theorem \ref{bijective correspondence}, $cd(A.G)=\{1, p^n\}$. If $\tau_{irr}(G)$ exists, then by Theorem \ref{projective embedding degrees}, $B_G \neq \emptyset$. Let $H \in B_G$. Since $\delta_{irr}(H)$ exists, $Z(H)$ is cyclic. By Theorem \ref{embedding degree of a p group with cyclic centre}, $\delta_{irr}(H)= p^n$. Hence, $\tau_{irr}(G) = p^n$ due to Theorem \ref{projective embedding degrees}. By Lemma \ref{one dim direct sum in non-abelian group}$(1)$, $\tau(G)= p^n$. %\poonam{Theorem \ref{projective embedding degrees} used twice}
%By the above arguments, if $\tau_{irr}(G)$ exists, then $B=\{H\}$. Hence, the result follows.
\end{proof}

%\textbf{Case (a)} Suppose $\mathrm{H}^2(G, \C^{\times})$ is not cyclic. Since $Z(G^*) \cong \mathrm{H}^2(G, \C^{\times})$ (by \cite[Theorem 3.1]{karpilovsky2}), $S$ must be non-trivial. Furthermore, by (\cite[Proposition 2.15]{gonzalo}), $S\subseteq \mathrm{H}^2(G, \C^{\times})$. As $Z(G^*)/S \subseteq Z\big(G^*/S\big)$, it follows $|H|=p^{k+1}$. By the above arguments, we obtain $\delta_{irr}(H) \leq \tau_{irr}(G)$. The result now follows from Theorem \ref{capable groups}. \\

%\textbf{Case (b)} Suppose $\mathrm{H}^2(G, \C^{\times})$ is cyclic. Then, $H=G^*$. The proof now follows similarly to the proof of Case (a).

% Since $G_1/A_1 \cong G_2/A_2$, $\tilde{\rho_i}\in \irr^{\beta\tra(\chi_i)}(G_2/A_2)$. Let $\gamma \in Z^2(G_2,\C^\times)$ be $\gamma(g,h)= \beta(gA_2,hA_2)$ $\forall$ $g$,$h$ $\in$ $G_2$ and $\phi_i \in \irr^{\gamma}(G_2)$ correspond to $\tilde{\rho_i}$. Let $\sigma= \oplus_{i=1}^l \sigma_i$. Since $\oplus_{i=1}^l \chi_i$ is faithful on $A_2$, by Theorem \ref{direct_sum}, $\sigma$ is faithful on $G_2$. Hence, $\tau(G_2) \leq \tau(G_1)$. By similar arguments, $\tau(G_1) \leq \tau(G_2)$. This completes the proof.

We suggest that readers keep all the cited papers listed below handy, as they will be used without further reference.
\begin{enumerate}
\item The classification of groups $G$ of order $p^3$, $p^4$ and $p^5$ for an odd prime $p$ is given in \cite{james}. 

\item For the capability and Schur multiplier of groups of order $p^5$, $p\geq 5$, see \cite[Table 1 and Table 2]{hatuip5}.

\item The Schur multiplier of groups of order $p^4$ are given in \cite{Ellis} and \cite{OK} for odd primes $p$, and we use GAP (\cite{GAP4}) for $p=2$.

\item For non-abelian groups $G$ of order $p^4$ and $p^5$ with elementary abelian Schur multiplier, to determine all suitable groups $H$ satisfying the hypotheses of Corollary~\ref{Elementary abelian Schur multiplier}(i), we refer to \cite[Table~4.1]{james} for odd primes $p$, and to \cite{hall2n} for groups of order $2^4$.

\item We also require information on $\delta(G)$ and $\delta_{irr}(G)$ for the above-mentioned groups $G$, which are discussed in \cite{kaur}. If $G$ is a group of order $p^5$ or $p^6$, we refer to \cite[Table~4.1]{james} to determine $\mathrm{cd}(G)$. 
\end{enumerate}
%\url{https://groupprops.subwiki.org/wiki/Groups_of_order_32} 
%Recall that if $G$ is not capable, then $\tau_{irr}(G)$ does not exist, in view of Corollary \ref{non-capable groups}. \Nayak{Can we write this in each result?} \poonam{Theorem 2.15 has Schur multiplier of groups of order $p^3$ and $p^4$ for $p \geq 5$.}
\begin{theorem} \label{Groups of order $p^3$}
For non-abelian groups $G$ of order $p^3$,  $\tau(G)$ and $\tau_{irr}(G)$ are given below.
(i) $\tau_{irr}(\Phi_2(21))$ does not exist and $\tau(\Phi_2(21)) =p+1$. \\
(ii) $\tau_{irr}(\Phi_2(1^3)) = \tau(\Phi_2(1^3)) =p$. \\
(iii) $\tau_{irr}(Q_8)$ does not exist and  $\tau(Q_8) =3$.  \\
(iv) $\tau_{irr}(D_8) = \tau(D_8) =2$. 

%We have the following table for non-abelian groups of order $p^3$.    
\end{theorem}

\begin{proof} 
$(i)$ Let $G= \Phi_2(21)$. Since $G$ is extra-special, by Lemma \ref{capability of extra special $p$-group}, $G$ is not capable. So $\tau_{irr}(G)$ does not exist due to Corollary \ref{non-capable groups}.  
By \cite[Theorem 3.1(2)]{kaur}, $\delta(G)= p$. Therefore, Lemma \ref{relation} and Remark \ref{some results on p-groups}$(iii)$ yield that $\tau(G)= p+1$. \\
$(ii)$ For $G= \Phi_2(1^3)$, $\tau(G)= \tau_{irr}(G)= p$, by Theorem \ref{Heisenberg group}. \\ %and Remark \ref{some results on p-groups}$(ii)$. \\ \poonam{Keep in mind that this is a Heisenberg group}
$(iii)$ The result follows from Example 5.4. \\
$(iv)$ The result follows from Example 5.1. \\
\end{proof}

\begin{theorem} \label{$p^4$}
For non-abelian groups $G$ of order $p^4$  ($p\geq 3)$, $\tau(G)$ and $\tau_{irr}(G)$ are given below. \\
(i) $\tau_{irr}(G)= p^2$ if $G$ is $\Phi_{2}(1^{4})$ or $\Phi_2(22)$ and $\tau_{irr}(\Phi_{3}(1^{4}))= \tau(\Phi_{3}(1^{4})) =p$. \\
(ii) $\tau_{irr}(G)$ exists if and only if $G$ is one of the groups: $\Phi_2(1^4)$, $\Phi_2(22)$ or $\Phi_3(1^4)$. \\
(iii) $\tau(G) = p+1$, if $Z(G)$ is cyclic and $G \ncong \Phi_3(1^4)$. \\
(iv) $\tau(G) = p+2$, if $Z(G) \cong \Z/p\Z \times \Z/p\Z$.

\iffalse
We have the following table for non-abelian groups of order $p^4$ $(p\geq 3)$. 

For non-abelian groups of order $p^4$$(p \geq 3$), the values of $\tau(G)$ and $\tau_{irr}(G)$ are as follows:

\begin{enumerate}
	\item $\tau_{irr}(G)= p^2$ if and only if $G$ is isomorphic to $\Phi_{2}(1^{4}),\Phi_{2}(22)$.
	\item For $G \cong \Phi_{3}(1^{4})$, $\tau(G)= \tau_{irr}(G)= p$. For the remaining groups, $\tau(G)= p+1$ if $Z(G)$ is cyclic, $p +2$ if $Z(G) \cong \Z/p\Z \times \Z/p\Z$ and $\tau_{irr}(G)$ does not exist.
	
\end{enumerate}
\fi
%The following table lists $\tau(G)$ and $\tau_{irr}(G)$ for non-abelian groups of order $p^4$, $p \geq 3$.     
\end{theorem}

\begin{proof}
$(i)$ Let $G_1 = \Phi_{2}(1^{4})$. Then, $B_{G_1}= \{H \mid |H|= p^5, H \in \Phi_{7} \}$. If $H \in B_{G_1}$, then $\delta_{irr}(H)= p^2$ in view of Theorem \ref{embedding degree of a p group with cyclic centre}. Hence, $\tau_{irr}(G_1)= p^2$ due to Corollary \ref{Elementary abelian Schur multiplier}$(i)$.

Let $G_2 = \Phi_2(22)$ and $G_3 = \Phi_{3}(1^{4})$. Then, $B_{G_2}= \{H \mid |H|= p^5, H \in \Phi_{8} \}$ and $B_{G_3}= \{H \mid |H|= p^5, H \in \Phi_{9} \text{ or } \Phi_{10}\}$. Using similar arguments, we get $\tau_{irr}(G_2)= p^2$ and $\tau_{irr}(G_3)= p$. By Remark \ref{some results on p-groups}$(ii)$, $\tau(G_3)= p$. 
%If $H \in B_{G_2}$, then $\delta_{irr}(H)= p^2$ in view of Theorem \ref{embedding degree of a p group with cyclic centre}. Hence, $\tau_{irr}(G_2)= p^2$ due to Corollary \ref{Elementary abelian Schur multiplier}$(i)$. 

%Consider $G_3 = \Phi_{3}(1^{4})$. Then, $B_{G_3}= \{H \mid |H|= p^5, H \in \Phi_{9} \text{ or } \Phi_{10}\}$. If $H \in \Phi_{9}$ with $|H|= p^5$, then $\delta_{irr}(H)= p$ in view of Theorem \ref{embedding degree of a p group with cyclic centre}. Similarly, if $H \in\Phi_{10}$ with $|H|= p^5$, then $\delta_{irr}(H)= p^2$. Hence, $\tau_{irr}(G_3)= p$ due to Corollary \ref{Elementary abelian Schur multiplier}$(i)$. 

%Consider the groups $H_1 \in \Phi_{7}$, $H_2 \in \Phi_{8}$ and $H_3 \in \Phi_{9}$ of order $p^6$. Since $G_i \cong H_i/Z(H_i)$, $i=1,2,3$, it follows from Corollary \ref{Elementary abelian Schur multiplier}$(i)$ and \cite[Theorem 3.1, Theorem 3.2]{kaur} that $\tau_{irr}(G_1)= \tau_{irr}(G_2)= p^2$ and $\tau_{irr}(G_3)= p$. \poonam{By \cite[Theorem 3.1, Theorem 3.2]{kaur}, $\delta_{irr}(H_1)= \delta_{irr}(H_2)= p^2$ and $\delta_{irr}(H_3)= p$. Since $G_i \cong H_i/Z(H_i)$, $i=1,2,3$, it follows from Corollary \ref{Elementary abelian Schur multiplier}$(i)$ that $\tau_{irr}(G_1)= \tau_{irr}(G_2) =p^2$ and $\tau_{irr}(G_3)=p$. Should I add like this everywhere?}

$(ii)$ If $G$ is one of the groups $\Phi_2(1^4)$, $\Phi_2(22)$ or $\Phi_3(1^4)$, then $\tau_{irr}(G)$ exists due to $(i)$. By \cite{hasanat}, these are the only three non-abelian groups that are capable. Hence, by Corollary \ref{non-capable groups}, $\tau_{irr}(G)$ does not exist for all other groups $G$.

%By [\cite{james}, Table 4.1], $\Phi_2(22)= \Gamma_2/Z(\Gamma_2)$, where $\Gamma_2 \in$ $\Phi_8$. Since, $\delta_{irr}(\Gamma_2)=p^2$ [\cite{kaur}, Theorem 3.1, 4,(vi)], by Remark \ref{some results on p-groups}, $\tau_{irr}(\Delta_2)\leq p^2$. \\
$(iii)$ From $(ii)$ and Remark \ref{some results on p-groups}$(iii)$, it follows that $\tau(G)>p$. Since $\delta(G)= p$ if $Z(G)$ is cyclic (by \cite[Theorem 3.1(3)]{kaur}), $\tau(G)= p+1$ due to Lemma \ref{relation}. \\
$(iv)$ By Remark \ref{some results on p-groups}$(ii)$ and Lemma \ref{direct sum of k irreducible ordinary representations}, $\tau(G)>p+1$. Since $\delta(G)= p+1$, by Lemma \ref{relation}, $\tau(G) \leq p+2$. Hence, the result follows.  

\end{proof}
In the following result, we use the notations for the groups of order $2^4$ as given in \cite{ConradGroups16}.
\begin{theorem}
For non-abelian groups $G$ of order $2^4$, $\tau(G)$ and $\tau_{irr}(G)$ are given below: \\
(i) If $G$ is one of the groups $ \mathbb{Z}/8\mathbb{Z} \rtimes_{3} \mathbb{Z}/2\mathbb{Z}$, $\mathbb{Z}/8\mathbb{Z} \rtimes_{5} \mathbb{Z}/2\mathbb{Z}$, $Q_{8} \rtimes \mathbb{Z}/2\mathbb{Z}$, $(\mathbb{Z}/4\mathbb{Z}) \rtimes (\mathbb{Z}/4\mathbb{Z})$ or $Q_{8} \times \mathbb{Z}/2\mathbb{Z}$, then $\tau_{irr}(G)$ does not exist. \\
(ii) $\tau(D_{16})= \tau_{irr}(D_{16})=2$;  $\tau(Q_{16}) =3$ and $\tau_{irr}(Q_{16})$ does not exist. \\
(iii) If $G$ is one of the groups $D_{8} \times \mathbb{Z}/2\mathbb{Z}$ or $(\mathbb{Z}/2\mathbb{Z})^{2} \rtimes \mathbb{Z}/4\mathbb{Z}$, then $\tau_{irr}(G) = \tau(G) = 4$. \\
(iv) If $G$ is one of the groups $(\mathbb{Z}/4\mathbb{Z}) \rtimes (\mathbb{Z}/4\mathbb{Z})$ or $Q_{8} \times \mathbb{Z}/2\mathbb{Z}$, then $\tau(G)= 4$. \\
(v) If $G$ is one of the groups $\mathbb{Z}/8\mathbb{Z} \rtimes_{5} \mathbb{Z}/2\mathbb{Z}$, $\mathbb{Z}/8\mathbb{Z} \rtimes_{3} \mathbb{Z}/2\mathbb{Z}$ or $Q_{8} \rtimes \mathbb{Z}/2\mathbb{Z}$, then $\tau(G)=3$. 
%If $G = D_8$, then $\tau(G)=\tau_{irr}(G)=2$, and $\tau(Q_{16})=2$. 
\end{theorem}

\begin{proof} 
$(i)$ For $G = \mathbb{Z}/8\mathbb{Z} \rtimes_{3} \mathbb{Z}/2\mathbb{Z}$ or $ \mathbb{Z}/8\mathbb{Z} \rtimes_{5} \mathbb{Z}/2\mathbb{Z}$, $M(G)= 1$. Hence, by Lemma \ref{one dim direct sum in non-abelian group}$(2)$, $\tau_{irr}(G)$ does not exist. If $G= Q_{8} \rtimes \mathbb{Z}/2\mathbb{Z}$, then $\tau_{irr}(G)$ does not exist due to \cite[Corollary 7.15]{NG}. \\ %For $G = D_8$, $\tau_{irr}(G)=2$, follows from Theorem \ref{examples}(3)(b). \\
If $G$ is one of the groups $(\mathbb{Z}/4\mathbb{Z}) \rtimes (\mathbb{Z}/4\mathbb{Z})$ or $Q_{8} \times \mathbb{Z}/2\mathbb{Z}$, then by \cite{hall2n}, there does not exist any group $H$ of order $2^5$ such that $H/Z(H) \cong G$. Hence, the result follows from Corollary \ref{Elementary abelian Schur multiplier}$(i)$. \\

%\poonam{Do I need to mention here that I have used GAP to check that their Schur multiplier is Elementary abelian?}
$(ii)$ For $G = D_{16}$ or $Q_{16}$, the result follows from Example 5.1 and Example 5.4. \\ 

$(iii)$ Let $G_1 = D_{8} \times \mathbb{Z}/2\mathbb{Z}$ and $G_2 = (\mathbb{Z}/2\mathbb{Z})^{2} \rtimes \mathbb{Z}/4\mathbb{Z}$. Consider $H_1 \in \Gamma_6$ and $H_2 \in \Gamma_7$ of order $2^5$. By \cite{hall2n}, $G_i \cong H_i/Z(H_i)$, $i=1,2$, and so $\tau_{irr}(G_i)$ exists due to Theorem \ref{capable groups}. By Example 5.1, Lemma \ref{direct sum of k irreducible ordinary representations} and Theorem \ref{one dim direct sum in non-abelian group}$(1)$, $\tau(G_i) > 3$. Hence, $\tau(G_i)= \tau_{irr}(G_i)= 4$. \\   %the result follows from Theorem \ref{Elementary abelian Schur multiplier} and \cite[Theorem 3.1, Theorem 3.2]{kaur}. By GAP \cite{GAP4}, $B_{G_1}= \{H \mid |H|= 2^5, H \in \tau_{6} \}$ and $B_{G_2}= \{H \mid |H|= 2^5, H \in \tau_{7} \}$. Hence, $\tau_{irr}(G_i)=2$, by Corollary \ref{Elementary abelian Schur multiplier}$(i)$ and \cite[Proof of Theorem 3.1,(1), (5)]{kaur}. \poonam{wrong, need to be checked}

%Since $(\Z/2\Z)^3 \subset G_1$, by Theorem \ref{elementary abelian p-group}, $\tau(G_1) \geqslant 4$. Also, as $\tau(G_1) \leq \tau_{irr}(G_1)=4$, therefore, $\tau(G_1)=4$. 
%Since $\Z/4\Z\ \times \Z/2\Z \subset G_2$, $\tau(G_2) \geqslant 3$ due to Theorem \ref{projective embedding degree of finite abelian groups}. 
%By Lemma  \ref{direct sum of k irreducible ordinary representations} and Theorem \ref{one dim direct sum in non-abelian group}$(1)$, $\tau(G_2) > 3$. Since, $\tau(G_2) \leq \tau_{irr}(G_2)=4$, it follows that $\tau(G_2) = 4$.  \\ 

$(iv)$ Suppose $G_3 = Q_{8} \times \Z/2\Z$. By Lemma \ref{direct sum of k irreducible ordinary representations} and Theorem \ref{one dim direct sum in non-abelian group}$(1)$, $\tau(G_2) > 3$.  Since $\delta(G_3)= 3$ (by \cite[Theorem 3.1(3)]{kaur}), $\tau(G_3) \leq 4$ due to Lemma \ref{relation}. Hence, $\tau(G_3)= 4$.  
Let $G_4 = D_{8} \times \mathbb{Z}/2\mathbb{Z}$. Since $(\Z/2\Z)^3 \subset G_4$, by Theorem \ref{elementary abelian p-group}, $\tau(G_4) \geqslant 4$. Also, as $\delta(G_4)= 3$, by Lemma \ref{relation}, $\tau(G_4) \leq 4$. Hence, $\tau(G_4)= 4$. \\  %\poonam{Lemma \ref{relation} used twice} 

$(v)$ %Using $(i)$ and Lemma \ref{one dim direct sum in non-abelian group}, $\tau(G)>2$. Since $\dg=2$, the result follows. 
Since $\delta(G)= 2$, the result follows from Example 5.2.

\end{proof}

In the next result, we study $\tau(G)$ and $\tau_{irr}(G)$ for non-abelian groups $G$ of order $p^5$, where $p \geq 5$.

\begin{theorem} \label{p5}
For non-abelian groups $G$ of order $p^5$, with $p \geq 5$, the following hold. 

(i) $G$ is not capable if and only if $G$ is one of the following groups: $\Phi_{2}(221)b, \Phi_{2}(311)a$, $ \Phi_{2}(311)b$, $\Phi_{2}(311)c$, $\Phi_{2}(32)a_1$, $\Phi_{2}(32)a_2$, $\Phi_{2}(41)$, $\Phi_{2}(2111)a$, $\Phi_{2}(2111)b$, $\Phi_{2}(2111)c$, \\ $\Phi_{2}(2111)d$, $\Phi_{3}(311)a, \Phi_{3}(311)b_r, \Phi_{3}(221)a, \Phi_{3}(2111)a$, $ \Phi_{3}(2111)b_r$, $\Phi_{3}(2111)c$, $\Phi_{3}(2111)d$, $\Phi_{3}(2111)e$, $\Phi_{4}(221)a, \Phi_{4}(221)c, \Phi_{4}(221)d_r \ (r \neq \tfrac{1}{2}(p-1))$, $\Phi_{4}(221)e$, $\Phi_{4}(221)f_r$, $\Phi_{4}(2111)a, \\ \Phi_{4}(2111)b, \Phi_{4}(2111)c$, $\Phi_{5}(2111)$, $\Phi_{5}(1^5)$, $\Phi_{6}(221)a$, $\Phi_{6}(221)b_r \ \big(r \neq \tfrac{1}{2}(p-1)\big)$, $\Phi_{6}(221)c_r$, $\Phi_{6}(221)d_r$, $\Phi_{6}(2111)a$, $\Phi_{6}(2111)b_r$, $\Phi_{7}(2111)a, \Phi_{7}(2111)b_r$, $\Phi_{7}(2111)c$, $\Phi_8(32)$, $\Phi_{9}(2111)a, \\ \Phi_{9}(2111)b_r$, $\Phi_{10}(2111)a_r$, $\Phi_{10}(2111)b_r$. 
For these groups, $\tau_{irr}(G)$ does not exist. \\

%\item For $G \in \Phi_{2}$, the following hold: \\

%\begin{enumerate}[label=(\alph*)]

(ii) If $G = \Phi_{2}(221)c$, then $\ti = p^2$. \\

(iii) If $G = \Phi_{2}(1^5)$, then $\ti = p^2$. \\

(iv) If $G= \Phi_{2}(221)a$ or $\Phi_{2}(221)d$, then $\ti$ does not exist.   \\ 

(v) Let $G \in \Phi_{2}$. Then one of the following holds:
\begin{enumerate}[label=(\alph*)]
	\item $\tau(G) = p+1$, if $Z(G)$ is cyclic;
	\item $\tau(G) = p+2$, if $Z(G) \cong \Z/p^{2}\Z \times \Z/p\Z$;
	\item $\tau(G) = p+3$ if $Z(G) \cong \Z/p\Z \times \Z/p\Z \times \Z/p\Z$. 
	\\
\end{enumerate}

(vi) If $G = \Phi_{3}(1^5)$ or $\Phi_{3}(221)b_{r}$, then $\ti= p^2$.  \\

(vii) For $G \in \Phi_{3} $, the following hold:
\begin{enumerate}[label=(\alph*)]
	\item $\tau(G) = p+1$, if $Z(G)$ is cyclic;
	\item $\tau(G) = p+2$, if $Z(G) \cong \Z/p\Z \times \Z/p\Z$.  \\
\end{enumerate} 

(viii) If $G = \Phi_{4}(221)b$ or $\Phi_{4}(1^5)$, then $\ti=p^2$.  \\

(ix) If $G = \Phi_{4}(221)f_0$ or $\Phi_{4}(221)d_{\tfrac{1}{2}(p-1)}$, then $\ti$ does not exist. \\ 

(x) If $G \in \Phi_{4}$ or $\Phi_{5}$, then $\tg = 2p$. \\
%\end{enumerate}

%(xi) If $G \in \Phi_{5}$, then $\tg = 2p$. \\

%\item For $G \in \Phi_{6} $, the following hold: \\

%\begin{enumerate}[label=(\alph*)]

(xi) If $G$ is one of the groups $\Phi_{6}(221)b_{\tfrac{1}{2}(p-1)}$, $\Phi_{6}(221)d_0$ or $\Phi_{6}(1^5)$, then $\ti= p^2$.  \\

(xii) If $G$ is one of the groups $\Phi_{6}(221)a$, $\Phi_{6}(221)b_r \ \big(r \neq \tfrac{1}{2}(p-1)\big)$, $\Phi_{6}(221)c_r$, $\Phi_{6}(221)d_r$, $\Phi_{6}(221)b_{\tfrac{1}{2}(p-1)}$, $\Phi_{6}(221)d_0$, $\Phi_6(2111)a$ or  $\Phi_6(2111){b_r}$, then $\tg = 2p+1$. \\

(xiii) If $G= \Phi_{6}(1^5)$, then $\tg = 2p$. \\

%\item For $G \in \Phi_{7}$, the following cases arise: \\

(xiv) If $G = \Phi_{7}(1^5)$, then $\ti= p^2$. \\

(xv) If $G \in \Phi_{7}$, then $\tg = 2p$. \\

%\end{enumerate}

(xvi) If $G = \Phi_8(32)$, then $\tau(G)=p^2+1$. \\

% For $G \in \Phi_{9}$, the following cases arise: \\

%\begin{enumerate}
(xvii) If $G= \Phi_{9}(2111)a$ or $\Phi_{9}(2111)b_r$, then $\tau(G) = p+1$. \\

(xviii) If $G = \Phi_{9}(1^5)$, then $\ti = \tg = p$. \\

(xix) If $G = \Phi_{10}(2111)a_r$ or $\Phi_{10}(2111)b_r$, then $\tg= p^2+1.$ \\ 

(xx) If $G = \Phi_{10}(1^5)$, then $\ti = \tg= p^2$.

\end{theorem}

\begin{proof}

$(i)$ By considering the list of capable groups given in \cite[Table 2]{hatuip5}, the result follows from Corollary \ref{non-capable groups}.  \\

$(ii)$ Suppose $G= \Phi_{2}(221)c =\langle \alpha, \alpha_1, \alpha_2, \gamma \mid [\alpha_1,\alpha] =\gamma^p= \alpha_2,\; \alpha^{p^2}= \alpha_1^{p} = \alpha_2^{p}=1 \rangle$. 
%\poonam{If somewhere in the presentation, it is $x^{(p)}$, I have considered it $x^p$}
Then the group 
\begin{eqnarray*}
	G^* &=& \langle \Bar{\alpha}, \Bar{\alpha_i}, \Bar{\gamma}, \gamma_i, \mid [\Bar{\alpha_1}, \Bar{\alpha}] = \Bar{\gamma}^p = \Bar{\alpha_2},\ [\Bar{\gamma}, \Bar{\alpha_1}] = \gamma_1,\ \\
	&& [\Bar{\gamma}, \Bar{\alpha}] = \gamma_2,\ \Bar{\alpha}^{p^2}= \Bar{\alpha_i}^{p} = \gamma_1^{p}= \gamma_2^{p^2}= 1 \ (i = 1,2) \rangle  
\end{eqnarray*}
is a representation group of $G$ such that ${G}^*/\langle \gamma_1, \gamma_2 \rangle \cong G$. Let $A= \langle \gamma_1, \gamma_2 \rangle$.
The proof uses the same argument as in the proof of Theorem~\ref{Heisenberg group}(1).
By \cite[Theorem 9.4]{Mackey}, it follows that any cocycle of $G$ is cohomologous to a cocycle $\nu$ of the form
$$
\nu\big(\gamma^{m_1}\alpha^{n_1}\alpha_1^{r_1}, \gamma^{m_2}\alpha^{n_2}\alpha_1^{r_2}\big) = \epsilon^{(m_2+p n_2r_1)n_1+\frac{r_1pn_2(n_2-1)}{2}}\lambda^{r_1m_2},
$$ for some $\epsilon$, $\lambda \in \C^\times$ such that $\epsilon^{p^2}= \lambda^{p}= 1$. It is easy to check that $1$ is the only $\nu$-regular element of $Z(G)$ if and only if $\nu(\alpha, \gamma)= \epsilon$ is a primitive $p^2$th root of unity. By Lemma \ref{faithful projective representation}, $G$ admits a faithful irreducible $\nu$-representation only for such $\nu$, and let $S$ be the set of all such $\nu$. Then $\tau_{irr}(G)= \operatorname{min}\{\operatorname{deg}(\operatorname{\rho}) \mid \rho \in \mathrm{Irr}^\nu(G), \nu \in S\}$. By Theorem \ref{bijective correspondence}, it follows that $$\tau_{irr}(G)=\operatorname{min}\{\operatorname{deg}(\operatorname{\Gamma}) \mid \Gamma \in \mathrm{Irr}(G^*\mid \chi), \chi\in \Hom(A, \C^\times) \text{ such that } \operatorname{tra}(\chi) \in S\}.$$ 
Let $T= \{\chi\in \Hom(A, \C^\times)\mid \chi(\gamma_2) \text{ is a primitive } p^2 \text{th root of unity}\}.$
First, we claim that $\operatorname{tra}(\chi) \in S$ if and only if $\chi\in T$. If $\chi^p= 1$, then $\operatorname{tra}(\chi)^p= 1$.  
%It is easy to observe that, if $\chi(\gamma_1)=\lambda$ and $\chi(\gamma_2) =\delta$, then \[ \operatorname{tra}(\chi) \big((\gamma^{m_1} \alpha^{n_1} \alpha_1^{p_1}), (\gamma^{m_2} \alpha^{n_2} \alpha_1^{p_2})\big) = \delta^{(m_2+p n_2p_1)n_1+ \frac{p_1pn_2(n_2-1)}{2}}\lambda^{p_1m_2},\]
Hence, if $\operatorname{tra}(\chi) \in S$, then $\chi\in T$. Conversely, let $\chi\in T$ and $\mu: G \to G^*$ be a section of $G^* \twoheadrightarrow G$ defined by $\mu(\gamma^{m} \alpha^{n} \alpha_1^{r}) = \Bar{\gamma}^{m} \Bar{\alpha}^{n} \Bar{\alpha_1}^{r}$. Then, 
$$\operatorname{tra}(\chi)(\alpha, \gamma) = \chi(\mu(\alpha)\mu(\gamma)\mu(\alpha\gamma)^{-1})=\chi(\Bar{\alpha}\Bar{\gamma}(\Bar{\gamma} \Bar{\alpha})^{-1})=\chi(\gamma_2^{-1}).$$ Thus, $\operatorname{tra}(\chi) \in S$ and the claim is proved.
Therefore, $$\tau_{irr}(G)=\operatorname{min}\{\operatorname{deg}(\operatorname{\Gamma}) \mid \Gamma \in \mathrm{Irr}(G^*\mid \chi),  \chi \in T\}.$$

Let $\chi \in T$ and $\Gamma \in \mathrm{Irr}(G^*\mid \chi)$. We have $\chi(\gamma_2)= \xi$,  where $\xi$ is a primitive $p^2$th root of unity. 
Now we 
%\poonam{use Clifford's theory to}(Clifford's theory written twice, can I remove from here?) 
show that $\operatorname{deg}(\Gamma)$ cannot be $p$.
Consider the normal abelian subgroup $N= \langle \Bar{\gamma}, \gamma_1, \gamma_2 \rangle$ of $G^*$ of order $p^5$ and $\chi_1 \in \Hom(N, \C^\times)$ such that $\chi_1|_A= \chi$. Let $\Gamma \in \mathrm{Irr}(G^*\mid \chi_1)$.
If $\Bar{\alpha_1} \notin {I_{{G}^*}(\chi_1)}$ and $\operatorname{deg}(\Gamma)=p$, then by  Clifford's theory, $\Bar{\alpha} \in {I_{{G}^*}(\chi_1)}$, which is not possible. Hence, $\Bar{\alpha_1} \in {I_{{G}^*}(\chi_1)}$. Then, $\chi_1$ extends to $\chi_2: N_1 \to \C^\times$, where $N_1= \langle N, \Bar{\alpha_1} \rangle$. Since $\chi_2^{(\Bar{\alpha}^p)}(\Bar{\gamma})\neq \chi_2(\Bar{\gamma})$, hence $\operatorname{deg}(\Gamma)$ cannot be $p$.  
Therefore, $\tau_{irr}(G)= p^2$. \\
% \Nayak{Same as Heisenberg groups.} \\

%let $[\beta]\in \mathrm{H}^2(H, \mathbb{C}^{\times}) \setminus \operatorname{ker}\operatorname{(inf)}$.

%Since $Z(G) \cong (\Z/p\Z)^3 $, from Theorem \ref{one dim direct sum in non-abelian group} and Lemma \ref{direct sum of k irreducible ordinary representations}, it follows that $\tau(G)= p+3$. \\

$(iii)$  Let $G= \Phi_{2}(1^{5})$. Then, $B_{G}=\{H \mid |H|= p^6, H \in \Phi_{22}\}$. If $H \in B_{G}$, then $\delta_{irr}(H)= p^2$ in view of Theorem \ref{embedding degree of a p group with cyclic centre}. Hence, $\tau_{irr}(G)= p^2$ due to Corollary \ref{Elementary abelian Schur multiplier}$(i)$. \\
%Let $G= \Phi_2(1^5)$ and $H \in \Phi_{22}$ be of order $p^6$. Since $G \cong H/Z(H)$, the result follows using Corollary \ref{Elementary abelian Schur multiplier}$(i)$ and \cite[Theorem 3.1, Theorem 3.2]{kaur}. \\

%\poonam{Should I give \cite[Theorem 3.2]{kaur} in preliminary?} \\

$(iv)$ If $G= \Phi_{2}(221)a$ or $\Phi_{2}(221)d$, it is easy to check that there does not exist any covering $H$ of order $p^6$ in \cite[Table 4.1]{james} such that $H/Z(H) \cong G$. Hence, the result follows from Corollary \ref{Elementary abelian Schur multiplier}$(i)$. \\ 

$(v)$ Let $G \in \Phi_2$. 
\begin{enumerate}[label=(\alph*)]
	\item If $Z(G)$ is cyclic, then $\delta(G)= p$, by \cite[Theorem 3.1(4)(i)]{kaur}. The result now follows from Remark \ref{some results on p-groups}$(iii)$. 
	\item  If $Z(G) \cong \Z/p^{2}\Z \times \Z/p\Z$, then $\delta(G)= p+1$ and the result follows from Remark \ref{some results on p-groups}$(iii)$, Lemma \ref{direct sum of k irreducible ordinary representations}. %\poonam{Remark \ref{some results on p-groups}$(iii)$ used twice}
	\item If $Z(G) \cong \Z/p\Z \times \Z/p\Z \times \Z/p\Z$, then $\delta(G)= p+2$. Using arguments similar to those in $(ii)$, the result follows. 
\end{enumerate}

$(vi)$  Let $G_1 = \Phi_{3}(1^{5})$. Then, $B_{G_1}= \{H \mid  |H|= p^6, H \in \Phi_{24} \text{ or } \Phi_{27}\}$. If $H \in B_{G_1}$, then $\delta_{irr}(H)= p^2$ in view of Theorem \ref{embedding degree of a p group with cyclic centre}. Hence, $\tau_{irr}(G_1)= p^2$ due to Corollary \ref{Elementary abelian Schur multiplier}$(i)$.
Let $G_2 = \Phi_{3}(221)b_{1}$ and $G_3 = \Phi_{3}(221)b_{\nu}$. Then $B_{G_2}= \{H \mid |H|= p^6, H \in \Phi_{25} \text{ or } \Phi_{28}\}$ and $B_{G_3}= \{H \mid |H|= p^6, H \in \Phi_{26} \text{ or } \Phi_{29} \}$.
By similar arguments, it can be shown that $\tau_{irr}(G_2)=\tau_{irr}(G_3)= p^2$. \\

% Let $G_2 = \Phi_{3}(221)b_{1}$. Then, $B_{G_2}= \{H \mid H \in \Phi_{25} \text{ or } \Phi_{28}, |H|= p^6\}$. If $H \in B_{G_2}$, then by Theorem \ref{embedding degree of a p group with cyclic centre} $\delta_{irr}(H)= p^2$. Hence, $\tau_{irr}(G_2)= p^2$, by  Corollary \ref{Elementary abelian Schur multiplier}$(i)$. 

% Consider $G_3 = \Phi_{3}(221)b_{\nu}$. Then, $B_{G_3}= \{H \mid H \in \Phi_{26} \text{ or } \Phi_{29}, |H|= p^6 \}$. If $H \in  B_{G_3}$, then $\delta_{irr}(H)= p^2$ in view of Theorem \ref{embedding degree of a p group with cyclic centre}. Hence, $\tau_{irr}(G_3)= p^2$ due to Corollary \ref{Elementary abelian Schur multiplier}$(i)$. \\

%Let $G_1 = \Phi_{3}(1^5)$, $G_2 = \Phi_{3}(221)b_{1}$ and $G_3= \Phi_{3}(221)b_{\nu}$. Consider the groups $H_1 \in \Phi_{24}$, $H_2 \in \Phi_{25}$ and $H_3 \in \Phi_{26}$ of order $p^6$. Since, $G_i \cong H_i/Z(H_i)$, $i=1,2,3$, the result follows  from Corollary \ref{Elementary abelian Schur multiplier}$(i)$ and \cite[Theorem 3.1, Theorem 3.2]{kaur}. \\

%\begin{enumerate}[label=(\arabic*)] 
$(vii)$ 
If $G \in \Phi_{3}$, then $cd(G)=\{1,p\}$ by \cite[Lemma 5.3]{Prajapati}. \\
$(a)$ If $Z(G)$ is cyclic, then by \cite[Theorem 2.32(b)]{isaacs}, $G$ has a faithful irreducible ordinary representation of degree $p$. So, $\delta(G)= p$, in view of Lemma \ref{one dim direct sum in non-abelian group}$(1)$. By Lemma \ref{relation} and Remark \ref{some results on p-groups}$(iii)$, $\tau(G) = p+1$. %\poonam{Remark \ref{some results on p-groups}$(iii)$ used twice}\\
$(b)$ If $Z(G)\cong \Z/p\Z \times \Z/p\Z$, then by Remark \ref{some results on p-groups}$(iii)$ and Lemma \ref{direct sum of k irreducible ordinary representations}, we have $\tau(G) > p+1$. 
Consider the group $$G= \Phi_{3}(221)a= \langle \alpha, \alpha_1, \alpha_2, \alpha_3
\mid [\alpha_1,\alpha]= \alpha_2,\; [\alpha_2,\alpha]= \alpha^{p}= \alpha_3,\;
\alpha_1^{p^2}= \alpha_2^{p}= \alpha_3^{p} =1 \rangle.$$ Then $Z(G)= \langle \alpha_3,\; \alpha_1^{p} \rangle$. Let $\chi \in \operatorname{Hom}(Z(G),\mathbb{C}^{\times})$ be defined on the generators by $\chi(\alpha_1^{p}) = 1$, $\chi(\alpha_3)= \xi$, where $\xi$ is a primitive $p$th root of unity. Let $\rho \in \operatorname{Irr}(G \mid \chi)$. As $\alpha_3 \in Z(G) \cap G^\prime$ and $\chi(\alpha_3)\neq 1$, we have $\operatorname{deg}(\rho)= p$. 
%$\chi$ extends to an ordinary character $\Phi:N \to \mathbb{C}^{\times}$ such that $\operatorname{deg}(\operatorname{Ind}^{G}_{N}(\Phi))=p$. Denote $\operatorname{Ind}^{G}_{N}(\Phi)$ by $\chi_1$. \\
Now, let $\chi_1 \in \operatorname{Hom}(G,\mathbb{C}^{\times})$ be defined on the generators by $\chi_1(\alpha_1)= \xi_1$, $\chi_1(\alpha)=1$, where $\xi_1$ is a primitive $p^2$th root of unity. Since $\chi \oplus \chi_1|_{Z(G)} \oplus 1_{Z(G)}$ is faithful on $Z(G)$, $\rho \oplus \chi_1 \oplus 1_G$ is faithful on $G$, follows from Lemma \ref{restriction}$(2)$. Hence, $\tau(G) = p+2$.  %\poonam{This paragraph also shows 55\% AI generated.} \\
Similarly, for all other groups $G$ whose centre is not cyclic, it can be shown that $\tg = p+2$. \\

%By Theorem \ref{projective embedding degree for direct product}, Theorem \ref{$p^4$} and Remark \ref{some results on p-groups}, $p < \tg \leq p+3$. By Lemma  \ref{direct sum of k irreducible ordinary representations}, $\tg$ cannot be $p+1$.  \\
$(viii)$ Let $G_1 = \Phi_{4}(221)b$. Then, $B_{G_1}= \{H \mid H \in \Phi_{34}, |H|= p^6 \}$. If $H \in B_{G_1}$, then $\delta_{irr}(H)= p^2$ by Theorem \ref{embedding degree of a p group with cyclic centre}. Hence, using Corollary \ref{Elementary abelian Schur multiplier}$(i)$, $\tau_{irr}(G_1)= p^2$.
Let $G_2 = \Phi_{4}(1^5)$. Then, $B_{G_2} = \{H \mid H \in \Phi_{31}, \Phi_{32} \text{ or } \Phi_{33}, |H|= p^6\}$. A similar argument shows that $\tau_{irr}(G_2)= p^2$. \\
%Then, $B_{G_2} = \{H \mid H \in \Phi_{31}, \Phi_{32} \text{ or } \Phi_{33}, |H|= p^6\}$. If $H \in B_{G_2}$, then $\delta_{irr}(H)= p^2$ in view of Theorem \ref{embedding degree of a p group with cyclic centre}. Hence, $\tau_{irr}(G_2)= p^2$ due to Corollary \ref{Elementary abelian Schur multiplier}$(i)$. 
%Let $G_1 = \Phi_{4}(221)b$ and $G_2 = \Phi_{4}(1^5)$. Consider $H_1 \in \Phi_{34}$ and $H_2 \in \Phi_{33}$ of order $p^6$. Since $G_i \cong H_i/Z(H_i)$, $i=1,2$, the result follows  from Corollary \ref{Elementary abelian Schur multiplier}$(i)$ and \cite[Theorem 3.1, Theorem 3.2]{kaur}. \\

$(ix)$
Consider the group $$G_1= \Phi_4(221){f_0}
= \langle \alpha, \alpha_i,  \beta_i \mid [\alpha_i, \alpha] = \beta_1,\ 
\alpha_1^p = \beta_2,\ \alpha_2^p = \beta_1^{\,\nu}, \ \alpha^p = \beta_1^p = 1 \ (i = 1,2) \rangle.$$ Then $G_1^* = \langle \alpha, \alpha_i,  \beta_i, \gamma \mid [\alpha_i, \alpha] = \beta_i,\ [\alpha_1, \alpha_2] = \gamma,\ \alpha_1^p = \beta_2,\ \alpha_2^p = \beta_1^{\,\nu}, \ \alpha^p = \beta_i^p = \gamma^{p^2}= 1 \ (i = 1,2) \rangle$ is a representation group of $G_1$ and $Z(G_1^*) = \langle \gamma, \beta_1, \beta_2 \rangle$. Since $M(G_1) \cong \Z/p^2\Z$, $\tau_{irr}(G_1)$ does not exist due to \cite[Chapter 4, Theorem 3.1]{karpilovsky2}. 
Next, consider the group $$G_2 = \Phi_4(221)d_{\tfrac{1}{2}(p-1)} = \langle \alpha, \alpha_i,  \beta_i \mid [\alpha_i, \alpha] = \beta_i,\ \alpha_1^p = \beta_1^{-1},\ \alpha_2^p = \beta_2,\ \alpha^p = \beta_i^p = 1 \ (i = 1,2) \rangle.$$ Then, $G_2^* = \langle \alpha, \alpha_i,  \beta_i, \gamma \mid [\alpha_i, \alpha] = \beta_i,\ [\alpha_1, \alpha_2] = \gamma,\ \alpha_1^p = \beta_1^{-1},\ \alpha_2^p = \beta_2,\ \alpha^p = \beta_i^p = \gamma^{p^2}= 1 \ (i = 1,2) \rangle$ is a representation group of $G_2$ and $Z(G_2^*) = \langle \gamma, \beta_1, \beta_2 \rangle$. Using the same argument as above, $\tau_{irr}(G_2)$ does not exist. \\

%Let $\beta$ be the corresponding cocycle in $\mathrm{Z}^2(G/Z(G), \C^{\times})$, where the correspondence is given by Remark \ref{remark for bijective correspondence via inflation}. Let $\tilde{\rho_1}$ and $\tilde{\rho_2}$ irreducible $\beta\mathrm{tra}(\chi_1)$ and $\beta\mathrm{tra}(\chi_2)$-representations of $G/Z(G)$, where $\chi_1$, $\chi_2 \in \operatorname{Hom}(Z(G),\mathbb{C}^{\times})$ such that $\chi_1 \oplus \chi_2$ is a faithful projective representation of $Z(G)$. Consequently, $G$ has a faithful projective representation of degree $2p$, in view of Theorem \ref{direct_sum}.

$(x)$ \label{Phi4} If $G \in \Phi_{4}$, then $Z(G) \cong \Z/p\Z \times \Z/p\Z$ and $G/Z(G) \cong (\Z/p\Z)^3$. By Lemma \ref{direct sum of k irreducible ordinary representations}, the direct sum of $p$-dimensional and one-dimensional ordinary representations of $G$ can not be faithful as projective. Hence, $\tau(G) >p+1$, in view of Remark \ref{some results on p-groups}$(iii)$. Since $Z(G) \subseteq G^\prime$,  $\psi|_{Z(G)}= 1$ for any $\psi \in \Hom(Z(G), \C^\times)$. This forces $\tau(G) >2p-1$. 
Now, Lemma \ref{exact sequence} yields the following exact sequence: \\

$1 \to \mathrm{Hom}(Z(G), \mathbb{C}^{\times}) 
\xrightarrow{\mathrm{tra}} 
\mathrm{H}^2(G/Z(G), \mathbb{C}^{\times}) 
\xrightarrow{\mathrm{inf}} \mathrm{H}^2(G, \mathbb{C}^{\times})$. \\

Let $\alpha$ be a non-trivial cocycle in $\operatorname{Im} \operatorname{(inf)}$, and let $\chi_i \in\mathrm{Hom}(Z(G), \mathbb{C}^{\times})$, $i=1,2$, be such that their direct sum $\chi_1 \oplus \chi_2$ is faithful on $Z(G)$. Let $\tilde{\rho_i} \in \mathrm{Irr}^{\beta \mathrm{tra}(\chi_i)}\big(G/Z(G)\big)$, $i=1,2$, and $\rho_i \in \mathrm{Irr}^{\alpha}(G)$ correspond to $\tilde{\rho_i}$ as described in Remark \ref{remark for bijective correspondence via inflation}. As $|G/Z(G)|= p^3$, so $\operatorname{deg}(\tilde{\rho_i})= p$. Now, by Theorem \ref{direct_sum}, $\rho_1 \oplus \rho_2$ is faithful of degree $2p$. Hence, $\tau(G)= 2p$.
If $G \in \Phi_{5}$, the result follows directly from Theorem \ref{extraspecial $p$-group}. \\
%Since $Z(G) \cong \Z/p\Z\times \Z/p\Z$, the result follows directly from [\cite{kaur}, Theorem 3.1], Lemma \ref{relation}, Theorem \ref{one dim direct sum in non-abelian group} and Lemma \ref{direct sum of k irreducible ordinary representations}.  \\  \poonam{Remark \ref{remark for bijective correspondence via inflation} used twice}  

$(xi)$ Let $G_1 = \Phi_{6}(221)b_{\tfrac{1}{2}(p-1)}$. Then, $B_{G_1}= \{H \mid |H|= p^6, H \in \Phi_{42}\}$. If $H \in B_{G_1}$, then $\delta_{irr}(H)=p^2$ in view of Theorem \ref{embedding degree of a p group with cyclic centre}. Hence, $\tau_{irr}(G_1)= p^2$ due to Corollary \ref{Elementary abelian Schur multiplier}$(i)$.
If $G = \Phi_{6}(221)d_0$ or $\Phi_{6}(1^5)$, then by similar arguments, $\tau_{irr}(G)= p^2$. \\

%Then, $B_{G_2}= \{H \mid |H|= p^6, H \in \Phi_{43}\}$. If $H \in B_{G_2}$, then $\delta_{irr}(H)= p^2$ in view of Theorem \ref{embedding degree of a p group with cyclic centre}. Hence, $\tau_{irr}(G_2)= p^2$ due to Corollary \ref{Elementary abelian Schur multiplier}$(i)$. 

%Consider $G_3 = \Phi_{6}(1^5)$. Then, $B_{G_3}= \{H \mid |H|= p^6, H \in \Phi_{40} \text{ or } \Phi_{41}\}$. If $H \in B_{G_3}$, then $\delta_{irr}(H)= p^2$ in view of Theorem \ref{embedding degree of a p group with cyclic centre}. Hence, $\tau_{irr}(G_3)= p^2$ due to Corollary \ref{Elementary abelian Schur multiplier}$(i)$. \\

%Let $G_1 = \Phi_{6}(221)b_{\tfrac{1}{2}(p-1)}$, $G_2 = \Phi_{6}(221)d_0$ and $G_3= \Phi_{6}(1^5)$. Consider $H_1 \in \Phi_{42}$, $H_2 \in \Phi_{43}$ and $H_3 \in \Phi_{41}$ of order $p^6$. Since $G_i \cong H_i/Z(H_i)$, $i=1,2,3$, the result follows from Corollary \ref{Elementary abelian Schur multiplier}$(i)$ and \cite[Theorem 3.1, Theorem 3.2]{kaur}. \\

$(xii)$ First, we show that if  $G\in \Phi_6$, then  $2p-1 < \tau(G) \leq 2p+1$. Proceeding along the same lines as the proof of $(x)$, we get $\tau(G) > 2p-1$. Now, let $\chi_1$ and $\chi_2$ be two non-trivial, distinct elements of $\Hom(Z(G), \C^\times)$ such that $\chi_1 \oplus \chi_2$ is a faithful ordinary representation of $Z(G)$. Let $\rho_i \in \operatorname{Irr}({G} \mid \chi_i)$, $i=1,2$. Since $cd(G)=\{1,p\}$ by  \cite[Table 4.1]{james}, $\operatorname{deg}(\rho_i)=p$. 
From Lemma \ref{restriction}(2), it follows that $(\oplus_{i=1}^2 \rho_i)\oplus 1_G$ is faithful on $G$. Hence, $\tau(G) \leq 2p+1$. 
By Lemma \ref{direct sum of k irreducible ordinary representations}, $G$ cannot possess an ordinary representation of degree $2p$ which is faithful. 

If $G$ is one of the groups $\Phi_{6}(221)a$, $\Phi_{6}(221)b_r \ \big(r \neq \tfrac{1}{2}(p-1)\big)$, $\Phi_{6}(221)c_r$ or $\Phi_{6}(221)d_r$, then $M(G)= 1$. Hence, $\tau(G) = 2p+1$.

Now, for other groups, it is enough to show that $G$ cannot possess a faithful $\alpha$-representation of degree $2p$ for any non-trivial cocycle $\alpha$.
\\
Let $G= \Phi_{6}(221)b_{\tfrac{1}{2}(p-1)}$ or $\Phi_{6}(221)d_0$ and $\rho \in \mathrm{Irr}^\alpha(G)$. By the proof of \cite[Theorem 4.3, $(xvii)$]{sabnam}, it follows that $\operatorname{deg}(\rho)= p^2$. Hence, the result holds. \\
Suppose $G_1 = \Phi_{6}(2111)a= \langle \alpha_1, \alpha_2, \beta, \beta_1, \beta_2
\mid [\alpha_1,\alpha_2]= \beta,\ [\beta,\alpha_i]= \beta_i,\ \alpha_1^p= \beta_1,\ \alpha_2^p= \beta^p= \beta_i^p= 1\ (i=1,2) \rangle$. From the proof of \cite[Theorem 4.3, (xvi)]{sabnam}, it follows that
\begin{eqnarray*}
	{G}_1^* & = \langle \alpha_1, \alpha_2, \beta, \beta_i,\ 1 \le i \le 3 \mid [\alpha_1,\alpha_2]  =\beta,\ [\beta,\alpha_j]= \beta_j,\ \\
	& [\beta_2,\alpha_2]= \beta_3,\ \alpha_1^p= \beta_1,\ \alpha_2^p = \beta^p = \beta_i^p =1\ (j=1,2) \rangle    
\end{eqnarray*}
is a representation group of $G_1$ of order $p^6$, and ${G}_1^*/\langle \beta_3 \rangle \cong G_1$. 
Now consider the abelian normal subgroup
$N_1= \langle \beta, \beta_i,\ 1 \le i \le 3 \rangle$ of ${G}_1^*$ of order $p^4$. 
Taking $A_1=\langle \beta_3 \rangle$, we have an isomorphism $\tra: \Hom(A_1,\C^\times) \to \mathrm{H}^2(G_1, \C^\times)$. Let $\chi' \in \mathrm{Hom}(\langle \beta_3 \rangle, \C^\times)$   such that $\chi'\neq 1$ and $\tra(\chi')=[\alpha]$. Suppose $\rho_i \in \mathrm{Irr}^\alpha(G_1)$ correspond to $\Gamma_i \in \operatorname{Irr}({G}_1^* \mid \chi')$ via the correspondence given in Theorem \ref{bijective correspondence}. Now, we show that if deg$(\rho_i) =p$, then $\rho_1 \oplus \rho_2$ is not faithful on $G_1$.
Suppose $\chi \in \mathrm{Irr}(N_1)$ such that $\chi|_{\langle \beta_3 \rangle}= \chi'$.
Let $\chi(\beta)=\xi^{i_1},\ \chi(\beta_1)=\xi^{i_2},\ 
\chi(\beta_2)=\xi^{i_3},\ \chi(\beta_3)=\xi^{i_4}$,
where $\xi$ is a primitive $p$th root of unity and
$0 \le i_1,i_2,i_3,i_4 \le (p-1)$. 
Since $i_4 \ne 0$ and deg$(\Gamma_i)= p$, from the proof of \cite[Theorem 4.3, $(xvi)$]{sabnam}, it follows that $i_2= 0$. In this case, $H_1 = I_{{G}_1^*}(\chi)= \langle N_1,\alpha_1\rangle$ is of order $p^5$. Let $\chi_1$, $\chi_2 \in \operatorname{Irr}(N_1)$ be such that $\chi_i|_{\langle \beta_3 \rangle}= \chi'$ and $\chi_i(\beta_1)=1$ for $i=1,2$. Then, $\Gamma_i = \operatorname{Ind}^{G_1^*}_{I_{G_1^*}(\chi_i)}(\tilde{\chi}_i)$, where $\tilde{\chi}_i:I_{G_1^*}(\chi_i)\to \C^\times$ is an extension of $\chi_i$. Choosing $t_i= \alpha_2^i$ for $i=1,2,..,p$, as left coset representatives of $H_1$ in $G_1^*$, we get $\Gamma_i(\beta_1)= I_{p \times p}$. Hence, $(\Gamma_1 \oplus \Gamma_2)(\beta_1)= I_{2p \times 2p}$ and so $\rho_1 \oplus \rho_2|_{Z(G_1)}$ is not faithful. Consequently, Lemma \ref{restriction}$(2)$ yields that $\tau(G_1)= 2p+1$. \\
Suppose $G_2= \Phi_6(2111){b_r}$. It follows from  \cite[Proof of Theorem 4.1(xvi)]{sabnam} that $(G_1,G_2)$ satisfies the hypotheses of Theorem \ref{tau(G_1) = tau(G_2)}. Hence, $\tau(G_2)= \tau(G_1)= 2p+1$. \\

%Since $i_4 \ne 0$. From the proof of \cite[Theorem 4.3, (xvi)]{sabnam}, it follows that if $\Gamma \in \operatorname{Irr}({G}_1^* \mid \chi)$ and $\operatorname{deg}(\Gamma)= p$, then $i_2= 0$ and in that case $H_1 = I_{{G}_1^*}(\chi)= \langle N_1,\alpha_1\rangle$ is of order $p^5$. Fix such $\chi$. Choosing $t_i= \alpha_2^i$, $i=1,2,..,p$, as left coset representatives of $H_1$ in $G^*$, we get $\Gamma(\beta_1)=I_{p \times p}$. \\ Let $\Gamma_i \in \operatorname{Irr}({G}_1^* \mid \chi_i)$, $i=1,2$, where $\chi_i \in \mathrm{Irr}(N_1)$ such that $\chi_i|_{\langle \beta_3 \rangle}=\chi'$. Let $\alpha= \operatorname{tra}(\chi')$ and $\rho_i \in \mathrm{Irr}^\alpha(G_1)$ correspond to $\Gamma_i$, where the correspondence is given by Theorem \ref{bijective correspondence}. Set $\rho= \rho_1 \oplus \rho_2$. Since $\rho|_{Z(G)}$ is not faithful, the result follows from Lemma \ref{restriction}(2). \\
%Consequently, if $\alpha=\operatorname{tra}(\chi')$ and $\rho_i \in \mathrm{Irr}^\alpha(G_1)$, $i=1,2$, then it follows from Corollary \ref{bijective correspondence} that $\rho_1 \oplus \rho_2$ cannot be faithful. Hence, the result follows.  \\

%We first show that $G_1$ cannot have a faithful $\alpha$-representation of degree $2p$, where $\alpha$ is a non-trivial cocycle. \\
\iffalse
Suppose $G_2= \Phi_6(2111){b_1} = \langle \alpha_1, \alpha_2, \beta, \beta_i \mid [\alpha_1,\alpha_2] = \beta,\ [\beta,\alpha_i] = \beta_i,\ \alpha_2^p = \beta_1,\ \alpha_1^p = \beta^p = \beta_i^p = 1\ (i = 1,2) \rangle$. Observe that the group 
\begin{eqnarray*}
	G_2^*= G_{23,4} &=& \langle \alpha_1, \alpha_2, \alpha_3, \beta, \beta_1, \beta_2 \mid [\alpha_1,\alpha_2] = \beta,\ [\beta,\alpha_i] = \beta_i,\ \\ \poonam{r=1}
	&& [\beta_2,\alpha_2] = \alpha_3,\ 
	\alpha_2^p = \beta_1 \alpha_3,\ \alpha_1^p = \alpha_3^p = \beta^p = \beta_i^p = 1\ (i = 1,2) \rangle    
\end{eqnarray*}
given in \cite{presentation} is a representation group of $G_2$ of order $p^6$ and ${G}_2^*/ \langle \alpha_3 \rangle \cong G_2$. 
Now consider the abelian normal subgroup
$N_2=\langle \alpha_3, \beta, \beta_i,\ 1 \le i \le 2 \rangle$ of ${G}_2^*$ of order $p^4$. 
Take $A_2=\langle \alpha_3 \rangle$.
Let $\chi' \in \mathrm{Irr}(A_2)$ such that $\chi'(\alpha_3) = \xi^{r}$,
where $\xi$ is a $p$th root of unity and $0 < r \le (p-1)$. Let $\chi \in \mathrm{Irr}(N_2)$ such that $\chi|_{A_2}=\chi'$. 
Let $\chi(\beta_1) = \xi^{s}$ and $\chi(\beta_2) = \xi^{t}$, where  $0 \le s, t \le (p-1)$. 

Our first claim is $\Gamma\in \mathrm{Irr}(G_2^*\mid \chi)$ is $p$-dimensional if and only if $s=0$.
We use Clifford's theory here. Now $g \in G_2^*$ can be written as $g =  \alpha_1^j \alpha_2^k n'$
for some $n' \in N_2$ and any element of $N_2$ is of the form $n_1 = \beta_2^l \beta^m z$
for some $z \in Z(G_2^*)$.
Therefore,
\[
\chi^{\alpha_1^j \alpha_2^k}(\beta_2^l \beta^m z)
= \chi(\beta_2^l \beta^m z)
\iff
\chi(\alpha_1^j \alpha_2^k \beta_2^l \beta^m  \alpha_2^{-k} \alpha_1^{-j}) = \chi(\beta_2^l \beta^m).\]
Upon simplification, we obtain that
$\alpha_1^j \alpha_2^k n' \in I_{G_2^*}(\chi)$
if and only if
\[
\chi\big(\beta_1^{m(p-j)} \beta_2^{p^2-km} \alpha_3^{p-kl+m\binom{p^2-k}{2}}\big) = 1 \text{ i.e., } \xi^{\,-smj-tkm+r(-kl+m\binom{p^2-k}{2})} = 1.
\] 
Suppose $s \neq 0$. Taking $m=0$, $l=1$, we have $\xi^{\,-rk} = 1$ which implies $k= 0$. Taking $m=1$, we have $\xi^{\,-sj} = 1$ which implies $j=0$. So $I_{{G}_2^*}(\chi)= N_2$. 
If $s= 0$, then $I_{{G}_2^*}(\chi)= \langle N_2, \alpha_1 \rangle$. 
Hence, our first claim follows.
Note that $\Gamma= \mathrm{Ind}_{I_{{G}_2^*}(\chi)}^{G_2^*}(\eta)$, where $\eta$ is a extension of $\chi$ from $N_2$ to $I_{{G}_2^*}(\chi)$.
Now, choosing $t_i= \alpha_2^i$, $i=1,2,..,p$, as left coset representatives of $I_{{G}_2^*}(\chi)$ in ${G}_2^*$, we get $\Gamma(\beta_1)= I_{p \times p}$. 
Proceeding along the same lines as the proof of $\tau(G_1)$, we get $\tau(G_2)= 2p+1$. \\

%Consequently, if $\alpha=\operatorname{tra}(\chi)$ and $\rho_i \in \mathrm{Irr}^\alpha(G_2)$, $i=1,2$, then it follows from Corollary \ref{bijective correspondence} that $\rho_1 \oplus \rho_2$ cannot be faithful. The conclusion now follows from (b) and (c). \\

$(xiii)$ 
Suppose $G_3= \Phi_6(2111){b_\nu}
= \langle \alpha_1, \alpha_2, \beta, \beta_i \mid [\alpha_1,\alpha_2] = \beta,\ 
[\beta,\alpha_i] = \beta_i,\ \alpha_2^p = \beta_1^\nu,\ \alpha_1^p = \beta^p = \beta_i^p = 1\ (i = 1,2) \rangle$. Observe that the group 
\begin{eqnarray*}
	G_3^*=G_{23,8\nu}&=& \langle \alpha_1, \alpha_2, \alpha_3, \beta, \beta_1, \beta_2 \mid [\alpha_1,\alpha_2] = \beta,\ [\beta,\alpha_i] = \beta_i,\ \\
	&& [\beta_2,\alpha_2] = \alpha_3,\ \alpha_1^p = \alpha_3,\ \alpha_2^p =\beta_1^\nu, \alpha_3^p = \beta^p = \beta_i^p = 1\ (i = 1,2) \rangle
\end{eqnarray*} given in \cite{presentation}, is a representation group of $G_3$ of order $p^6$ and ${G}_3^*/\langle \alpha_3 \rangle \cong G_3$. 
Consider the abelian normal subgroup
$N_3=\langle \beta, \beta_i, \alpha_3, \ 1 \le i \le 2 \rangle$ of ${G}_3^*$ of order $p^4$. 
Take $A_3=\langle \alpha_3 \rangle$. 
% Since $\tra: \Hom(A_3,\C^\times)\to \mathrm{H}^2(G_3,\C^\times)$ is an isomorphism, for each non-trivial cocycle $\alpha$ of $G_3$, there is a non-trivial $\chi' \in \mathrm{Hom}(A_3, \C^\times)$ such that $\tra(\chi')=[\alpha]$. 
Let $\chi' \in \mathrm{Irr}(\langle \alpha_3 \rangle)$ such that $\chi'(\alpha_3)=\xi$
where $\xi$ is a primitive $p$th root of unity.  
Our aim is to find   $\Gamma_i \in \operatorname{Irr}(G_3^* \mid \chi')$ with deg$(\Gamma_i)=p$ such that $\Gamma_1 \oplus \Gamma_2$  is  faithful on $\langle \beta_1, \beta_2\rangle \subseteq G_3^*$.
We use Clifford's theory here. Let $\chi \in \operatorname{Irr}(N_3)$ such that
$\chi|_{A_3} = \chi'$ and $\chi(\beta_2) = 1$, $\chi(\beta) = \xi$ and $\chi(\beta_1) = \xi^j$ where  $0 < j \le (p-1)$. Now $g \in G_3^*$ can be written as $g = \alpha_1^k \alpha_2^l n'$
for some $n' \in N_3$ and any element of $N_3$ is of the form $n_1 = \beta_2^m \beta^n z$ for some $z \in Z(G_3^*)$.
Therefore,
$$\chi^{(\alpha_1^k \alpha_2^l)}(\beta_2^m \beta^n z) = \chi(\beta_2^m \beta^n z)
\iff \chi(\alpha_1^k \alpha_2^l \beta_2^m \beta^n  \alpha_2^{-l} \alpha_1^{-k}) = \chi(\beta_2^m \beta^n).$$
Upon simplification,  we obtain that
$\alpha_1^k \alpha_2^l n' \in I_{G_3^*}(\chi)$
if and only if
\[
\chi(\alpha_3^{n\binom{p^2-l}{2}-lm} \beta_2^{m-nl} \beta_1^{-nk}\beta^{n}) = 1 \text{ i.e., } \xi^{(n\binom{p^2-l}{2}-lm)-nkj+n} = 1.
\] 
Taking $n=0$, $m=1$, we have $\xi^{-l} = 1$ implies $l=0$.
Hence, $I_{G_3^*}(\chi) = \langle N_3, \alpha_1 \rangle$ is of order $p^5$. 
Let $\chi_1$, $\chi_2 \in \operatorname{Irr}(N_3)$ such that $\chi_i(\beta_1)=\xi^i$ and $\chi_i|_{\langle \beta,\beta_2, \alpha_3\rangle} = \chi |_{\langle \beta,\beta_2, \alpha_3\rangle}$. Let $\Gamma_i = \operatorname{Ind}^{G_3^*}_{I_{G_3^*}(\chi_i)}(\tilde{\chi}_i)$, where $\tilde{\chi}_i:I_{G_3^*}(\chi_i)\to \C^\times$ is an extension of $\chi_i$. Let $\{l_i=\alpha_2^i, i=1,2,..,p\}$ be a set of left coset representatives of $I_{G_3^*}(\chi_i)$ in $G_3^*$. Then $\Gamma_1(\beta_1)= \xi I_{p \times p}$, $\Gamma_1(\beta_2)= {D}_{p \times p}$ and $\Gamma_2(\beta_1)= \xi^2 I_{p \times p}$, $\Gamma_2(\beta_2)=  {D}_{p \times p}$, where $D_{p \times p}$ is a diagonal matrix such that $D_{ii}=\xi^{i}$. 
Consider the $\tra: \Hom(A_3,\C^\times)\to \mathrm{H}^2(G_3,\C^\times)$ and suppose $\alpha= \operatorname{tra}(\chi')$. For $i=1,2$, let  $\rho_i \in \mathrm{Irr}^\alpha(G_3)$ correspond to $\Gamma_i$ via the correspondence given in Theorem \ref{bijective correspondence}. Let $\rho= \rho_1 \oplus \rho_2$. 
%Using Clifford's theory, we can check that if $\Gamma \in \operatorname{Irr}({G}_1^* \mid \chi')$ is of degree $p$, then $|{G}_1^*/I_{{G}_1^*}(\chi)|=p$. 
Then, $$\rho(\beta_1^m \beta_2^n)= \left[\begin{array}{cccc|cccc}
	\xi^{m}\xi^{n} & & & & & & &\\
	& \xi^{m}\xi^{2n} & & & & & & \\
	& & \ddots & & & & & \\
	& & & \xi^{m} & & & & \\
	\hline
	& & & & \xi^{2m}\xi^{n} & & & \\
	& & & & & \xi^{2m}\xi^{2n} & &  \\
	& & & & & & \ddots & \\
	& & & & & & & \xi^{2m}  \\
\end{array}\right]_{2p \times 2p}, 0\leq m, n \leq p-1$$
is a diagonal matrix.
Since $\rho|_{Z(G_3)}$ is faithful, it follows from Theorem \ref{restriction}$(2)$ that $\rho$ is faithful. Hence, $\tau(G_3) = 2p$ as $2p-1 < \tau(G) \leq 2p+1$ for $G\in \Phi_6$ (follows from the proof of $(xii)$). \\
\fi  %Proceeding along the same lines as the proof of (xv), we get $\tau(G)= 2p$.  \\ 
%If for $r_1 \neq r_2$,  $\operatorname{deg}(\rho_i)=p$, then $\rho$ is faithful. 

%The proof proceeds analogously to that of part (a) in (3). 

(xiii) Let $G = \Phi_{6}(1^5) = \langle \alpha_1, \alpha_2, \beta, \beta_i \mid 
[\alpha_1,\alpha_2] = \beta,\,
[\beta,\alpha_i] = \beta_i,\,
\alpha_i^p = \beta^p = \beta_i^p = 1 \ (i=1,2) \rangle$. Now consider the following group 
\begin{eqnarray*}
	G^*&=&\langle \alpha_1, \alpha_2, \gamma ,\beta, \beta_j, 1\leq j \leq 4  \mid 
	[\alpha_1,\alpha_2] = \beta,\, 
	[\beta,\alpha_i] = \beta_i,\, [\beta_2,\alpha_2] = \beta_3,\, \\
	&& [\beta_1,\alpha_1] = \beta_4,\,
	[\beta_1,\alpha_2] = [\beta_2,\alpha_1] = \gamma,\,
	\alpha_i^p = \beta^p = \beta_j^p = \gamma^p = 1 \ (i=1,2) \rangle \\
	&=& \big(\langle  \alpha_2, \gamma ,\beta, \beta_j, 1\leq j \leq 3  \mid 
	[\beta,\alpha_2] = \beta_2,\, [\beta_2,\alpha_2] = \beta_3,\, [\beta_1,\alpha_2] =  \gamma,\, \\
	&& \alpha_2^p = \beta^p = \beta_j^p = \gamma^p = 1 \ \rangle \times \langle \beta_4 \rangle\big) \rtimes \langle \alpha_1 \rangle \\
	&=&\big(\Phi_{16}(1^6) \times  \Z/p\Z \big) \rtimes  \Z/p\Z.
\end{eqnarray*}
Then $G^*$ is a representation group of $G$ of order $p^8$ such that $G^*/Z(G^*)= G$, where $Z(G^*)= \langle \beta_3, \beta_4, \gamma \rangle$. 
Consider the abelian normal subgroup $N =\langle \beta, \gamma, \beta_i, 1 \leq i \leq 4 \rangle$ of ${G}^*$ of order $p^6$. Take $A= Z(G^*)$. Let $\chi' \in \mathrm{Irr}(A)$ such that $\chi'(\beta_3) = \xi$, $\chi'(\beta_4) = \xi$ and $\chi'(\gamma) = \xi$ where $\xi$ is a primitive $p$th root of unity. First, our aim is to find $\Gamma_i \in\operatorname{Irr}(G^* \mid \chi')$, $i=1,2$, of degree $p$. %\poonam{We are trying to find two irreducible projective representations of degree p whose direct sum is faithful.} 
We use Clifford's theory here. Let $\chi \in \operatorname{Irr}(N)$ such that
$\chi|_{A} = \chi'$. Now $g \in {G}^*$ can be written as $g = \alpha_1^m \alpha_2^n n'$
for some $n' \in N$.
%and any element of $N$ is of the form $n_1 = \beta^j \beta_1^k \beta_2^l z$ for some $z \in Z({G}^*)$.
Now we have
$$\chi^{(\alpha_1^m \alpha_2^n)}(\beta^j \beta_1^k \beta_2^l z) = \chi(\beta^j \beta_1^k \beta_2^l z)
\iff \chi(\alpha_1^m \alpha_2^n \beta^j \beta_1^k \beta_2^l \alpha_2^{-n} \alpha_1^{-m}) = \chi(\beta^j \beta_1^k \beta_2^l).$$
Upon simplification, we obtain that
$\alpha_1^m \alpha_2^nn' \in I_{{G}^*}(\chi)$
if and only if
\[
\chi(\beta_1^{-jm} \beta_2^{-jn} \beta_3^{j\binom{p-n}{2}-ln} \beta_4^{-km+j\binom{p-m}{2}} \gamma^{-jnm-lm-kn}) = 1.
\]
Let $\chi(\beta_1) = 1$ and $\chi(\beta_2) = \xi^r$ where $0 < r \le (p-1)$. Then,
\[\xi^{-rjn+(j\binom{p-n}{2}-ln)+(-km+j\binom{p-m}{2})-jnm-lm-kn} = 1.\] 
Taking $j=0$, $k=0$, $l=1$, we have $\xi^{-n-m} = 1$, which implies $n= -m$.
Hence, $N_1= I_{{G}^*}(\chi) = \langle N, \alpha_1^{-2^{-1}r}\alpha_2^{2^{-1}r} \rangle= \langle N, \alpha_1^{-1} \alpha_2 \rangle$ is of order $p^7$. Let $\chi_1$ and $\chi_2 \in \operatorname{Irr}(N)$ be such that $\chi_i(\beta_1)= 1$ and $\chi_1(\beta_2)= \xi$, $\chi_2(\beta_2)= \xi^2$, $\chi_i|_{A} = \chi'$ for $i=1,2$. 
%Then, $I_{G^*}(\chi_1)= \langle N, \alpha_1^{-2^{-1}} \alpha_2^{2^{-1}} \rangle $ and $I_{G^*}(\chi_2)= \langle N, \alpha_1^{-1} \alpha_2 \rangle$. 
Since $I_{G^*}(\chi_i)/N$, $i=1,2$ is cyclic, $\chi_i$ extends to $\tilde{\chi}_i: I_{G^*}(\chi_i)\to \C^\times$. Consider $\Gamma_i = \operatorname{Ind}^{G^*}_{I_{G^*}(\chi_i)}(\tilde{\chi}_i)$. Let $\{l_t= \alpha_2^t, t=1,2,..,p\}$ be a set of left coset representatives of $I_{G^*}(\chi_i)$ in $G^*$. Then $\Gamma_i(\beta_1)= D_{p \times p}$, $\Gamma_1(\beta_2)= \xi{D}_{p \times p}$ and $\Gamma_2(\beta_2)= \xi^2 D_{p \times p}$, where $D_{p \times p}$ is a diagonal matrix such that $D_{ii}= \xi^{i}$. Let $\alpha= \operatorname{tra}(\chi')$ and $\rho_i \in \mathrm{Irr}^\alpha(G)$ correspond to $\Gamma_i$, where the correspondence is given by Theorem \ref{bijective correspondence}. Let $\rho= \rho_1 \oplus \rho_2$. 
Then, for $0\leq m, n \leq p-1$,
$$\rho(\beta_1^m \beta_2^n)= \left[\begin{array}{cccc|cccc}
	\xi^{m}\xi^{n}\xi^{n} & & & & & & &\\
	& \xi^{2m}\xi^{n}\xi^{2n} & & & & & & \\
	& & \ddots & & & & & \\
	& & & \xi^{n} & & & & \\
	\hline
	& & & & \xi^{m}\xi^{2n}\xi^{n} & & & \\
	& & & & & \xi^{m}\xi^{2n}\xi^{2n} & &  \\
	& & & & & & \ddots & \\
	& & & & & & & \xi^{2n}  \\
\end{array}\right]_{2p \times 2p}$$
is a diagonal matrix.
Since $\rho|_{Z(G)}$ is faithful, it follows from Theorem \ref{restriction} that $\rho$ is faithful. Hence, $\tau(G) = 2p$ as $2p-1 < \tau(G) \leq 2p+1$ for $G\in \Phi_6$ (as shown in the proof of $(xii)$). \\

$(xiv)$ Let $G= \Phi_{7}(1^{5})$. Then, $B_{G}=\{H \mid |H|= p^6, H \in \Phi_{30}\}$. If $H \in B_{G}$, then $\delta_{irr}(H)= p^2$ in view of Theorem \ref{embedding degree of a p group with cyclic centre}. Hence, $\tau_{irr}(G)= p^2$ due to Corollary \ref{Elementary abelian Schur multiplier}$(i)$.  \\

$(xv)$ \label{Phi7}
%Let $G$ be one of the groups $\Phi_{7}(2111)a, \Phi_{7}(2111)b_r$ or $\Phi_{7}(2111)c$.  
%Additionally, according to Theorem \ref{cohomological characterization}, the map $\mathrm{inf}: \mathrm{H}^2(G/Z(G), \C^{\times}) \to \mathrm{H}^2(G, \C^{\times})$ is surjective. \\ \poonam{(surjective can also be shown by Theorem \ref{epicenter})}
If $G \in \Phi_{7}$, then $Z(G)\cong \Z/p\Z$ and $G/Z(G) \cong \Phi_2(1^4)$. Let $G_1= G/Z(G)$. Then 
\[
G_1 \cong \langle \alpha, \alpha_i, 1\leq i \leq 3 \mid [\alpha_1,\alpha] = \alpha_2,\; \alpha^{p}= \alpha_i^{p} =1 \rangle.
\]
Using \cite[Theorem 9.4]{Mackey}, it is easy to check that every cocycle of $G_1$ is cohomologous to a cocycle $\gamma$ of the form \[
\gamma\big(\alpha_3^{\omega_1} \alpha_2^{x_1} \alpha^{y_1} \alpha_1^{z_1},\alpha_3^{\omega_2} \alpha_2^{x_2} \alpha^{y_2} \alpha_1^{z_2}\big) = \lambda^{x_2 z_1+\frac{y_2 z_1\left(z_1-1\right)}{2}}\mu^{ y_1 x_2+\frac{z_1 y_2(y-1)}{2}+z_1 y_1 y_2}\tau_y^{y_1 \omega_2} \tau_z^{z_1 \omega_2},\] where $\lambda,\mu,\tau_y,\tau_z \in \C^\times$ such that $\lambda^p= \mu^p = \tau_y^p = \tau_z^p= 1.$ 
Let $\lambda= e^{\frac{2 \pi i j}{p}}, \mu= e^{\frac{2 \pi i k}{p}}, \tau_y= e^{\frac{2 \pi i l}{p}}, \tau_z= e^{\frac{2 \pi i m}{p}}$. First, we show that if $\left[\begin{array}{l}j \\ k\end{array}\right]$ and $\left[\begin{array}{l}m \\  l \end{array}\right]$ are linearly dependent over $\Z/p\Z$, then $Z(G_1)$ has a non-trivial $\gamma$-regular element. If $(j,k)=n(m,l)$ for some $n\in \Z/p\Z$, then it is easy to check that $\alpha_3^{-n} \alpha_2$ is a $\gamma$-regular element of $Z(G_1)$. If $(m,l)=(0,0)$, then $\alpha_3$ is a $\gamma$-regular element of $Z(G_1)$. Hence, for any such $\gamma$, all irreducible $\gamma$-representations of $G_1$ are of degree $p$, in view of \cite[Chapter 2, Lemma 5.1]{karpilovsky3}. This implies  
that, if $X= \{[\nu] \in \mathrm{H}^2(G_1,\C^\times) \mid  \irr^\nu(G_1) \text{ contains an element of degree } p^2\}$, then $|X| \leq (p^2-1)(p^2-p)$.
Now, using the exact sequence from Lemma \ref{exact sequence}
$$1 \to \mathrm{Hom}(Z(G), \mathbb{C}^{\times}) 
\xrightarrow{\mathrm{tra}} 
\mathrm{H}^2(G_1, \mathbb{C}^{\times}) 
\xrightarrow{\mathrm{inf}} \mathrm{H}^2(G, \mathbb{C}^{\times}),$$ 
we observe that there are $(p^3-1)$ non-trivial $[\alpha]$ in $\operatorname{Im}\operatorname{(inf)}$.
Therefore, for each such $\alpha$, there is a $\beta \in Z^2(G_1, \mathbb{C}^{\times})$ such that $$\mathrm{Irr}^{\alpha}(G) \longleftrightarrow \bigcup_{\chi \in \mathrm{Hom}(Z(G),\C^{\times})} \mathrm{Irr}^{\beta\mathrm{tra}(\chi)}(G_1),$$ by Remark \ref{remark for bijective correspondence via inflation}(1).
If, for every non-trivial $[\alpha] \in \operatorname{Im} \operatorname{(inf)}$, there exist at least $p-1$ distinct $\chi$ such that each $[\beta \mathrm{tra}(\chi)]\in X$, then $(p^3-1)(p-1)\leq |X|$, contradicting to the fact $|X| \leq (p^2-1)(p^2-p)$. 
%cocycles give $p^2$-dimensional representations, 
Hence, there is an $\alpha \in \mathrm{Z}^2(G, \C^{\times})$ and $\chi_1, \chi_2 \in \mathrm{Hom}(Z(G), \mathbb{C}^{\times})$ with $\chi_1 \neq \chi_2$ such that
%such that for the corresponding $\beta \in \mathrm{Z}^2(G_1, \C^{\times})$, 
the elements in $\mathrm{Irr}^{\beta \mathrm{tra}(\chi_i)}\big(G_1\big)$ are of degree $p$ for  $i=1,2$.
%irreducible $\beta\mathrm{tra}(\chi_1)$ and $\beta\mathrm{tra}(\chi_2)$-representations of $G_1$ are of degree $p$.
Now, let $\tilde{\rho_i} \in \mathrm{Irr}^{\beta \mathrm{tra}(\chi_i)}\big(G_1\big)$ and $\rho_i\in \mathrm{Irr}^{\alpha}(G)$ correspond to $\tilde{\rho_i}$, via the correspondence in Remark \ref{remark for bijective correspondence via inflation}(2). Set $\rho= \rho_1 \oplus \rho_2$. By Theorem \ref{direct_sum}, $\rho$ is faithful. Hence, $\tau(G) \leq 2p$.  
Let $\psi$ be a faithful projective representation of $G$ of degree $\tau(G)$. Since $\delta(G) =p^2$, $\psi$ can not be an ordinary representation. Therefore, $\psi$ has to be a true projective representation and $\tau(G)= p$ or $2p$. By Remark \ref{some results on p-groups}$(iii)$, $\tau(G)> p$. Hence, the result follows. \\ 

$(xvi)$ Let $G = \Phi_8(32)$. By \cite[Theorem 3.1(4)(vi)]{kaur}, $\delta(G)= p^2$. Hence, the result follows from Lemma \ref{group with cyclic centre of order p}. \\

$(xvii)$ Let $G$ be one of the groups $\Phi_{9}(2111)a$ or $\Phi_{9}(2111)b_r$. Since $\delta(G)= p$ (by \cite[Theorem 3.1(4)(vii)]{kaur}), $\tau(G) \leq p+1$ due to Lemma \ref{relation}. The result is a consequence of Remark \ref{some results on p-groups}$(iii)$. \\

$(xviii)$ Let $G= \Phi_{9}(1^{5})$. Then, $B_{G}=\{H \mid |H|= p^6, H \in \Phi_{35}, \Phi_{36} \text{ or } \Phi_{37}\}$. If $H \in \Phi_{35}$ with $|H|= p^6$, then $\delta_{irr}(H)= p$ in view of Theorem \ref{embedding degree of a p group with cyclic centre}. Similarly, if $H \in \Phi_{36}$ or $\Phi_{37}$ with $|H|= p^6$, then $\delta_{irr}(H)= p^2$. Hence, $\tau_{irr}(G)= p$ due to Corollary \ref{Elementary abelian Schur multiplier}$(i)$. By Remark \ref{some results on p-groups}$(ii)$, $\tau(G)= p$. \\

%By \cite[Table 4.1]{james}, $G= \Gamma/Z(\Gamma)$, where $\Gamma \in$ $\Phi_{35}$. By \cite[Table 4.1]{james} and [\cite{kaur}, Theorem 3.3], $\delta_{irr}(\Gamma)=p$. Hence, it follows from Remark \ref{irreducible projective embedding degree for capable groups} and \ref{some results on p-groups} that $\tg = \ti = p$. \\
%\poonam{rewritten proof using Theorem \ref{projective embedding degrees}}
$(xix)$ Suppose $G = \Phi_{10}(2111)a_r$ or $\Phi_{10}(2111)b_r$. %Since $Z^*(G)=Z(G)$, by Theorem \ref{epicenter}, inf : $\mathrm{H}^2(G/Z(G), \mathbb{C}^{\times}) \rightarrow\mathrm{H}^2(G,\mathbb{C}^{\times})$ is surjective. Let $\tau(G)=m$ and $\rho$ be a faithful $\alpha$-representation of $G$ of degree $m$. By \cite[Proposition 2.15]{gonzalo}, there is a covering $A.G$ of $G$ such that $\rho$ lifts to a faithful ordinary representation of $A.G$, where $A \subseteq M(G)$. 
As $M(G) \cong \Z/p\Z$, we have $A.G =G$ or $G^*$. Since $cd(G)=\{1,p,p^2\}$ (by \cite[Lemma 5.8]{Prajapati}), any irreducible projective representation of $G$ is of degree $1, p$ or $p^2$. Then Theorem \ref{embedding degree of a p group with cyclic centre} and Theorem \ref{bijective correspondence} yield that $\delta(G^*)= p^2$. Also, as $\delta(G)= p^2$, $\tau(G) \geq p^2$, due to Theorem \ref{projective embedding degrees}.  

Let $\rho = \oplus_{i=1}^l\rho_i$, where $\rho_i \in \mathrm{Irr}^\alpha(G)$. Let $\tilde\rho_i$ be $\beta\mathrm{tra}(\chi_i)$ representations of $G/Z(G)$ corresponding to $\rho_i$ via the correspondence given in Remark \ref{remark for bijective correspondence via inflation}. As $\rho$ is faithful, by Theorem \ref{direct_sum}, $\chi_i \neq \chi_j$ for $i \neq j$ and $\oplus_{i=1}^l\chi_i$ is faithful on $Z(G)$. Since $\chi_i \oplus \chi_j$ is faithful on $Z(G)$, we have $l=2$. If $\operatorname{deg}(\rho_i)<p^2$ for $i=1,2$, then $\operatorname{deg}(\rho) \leq 2p$ which is not possible.
Hence, $\operatorname{deg}(\rho) \geq p^2+p$ if $\alpha$ is a non-trivial cocycle and $\operatorname{deg}(\rho) \geq p^2+1$ if $\alpha=1$. Since $\delta(G)= p^2$, the result follows from Lemma \ref{relation}.  \\

%If $\rho$ is an ordinary representation, let $\rho_i \in $ \\

$(xx)$ Let $G= \Phi_{10}(1^5)$. Then, $B_{G}=\{H \mid |H|= p^6, H \in \Phi_{38} \text{ or } \Phi_{39}\}$. By arguments similar to those in the proof of (xviii), we get $\tau_{irr}(G)=p^2$. 
%Thus, $\tau_{irr}(G)=p^2$, using Corollary \ref{Elementary abelian Schur multiplier}$(i)$ and \cite[Theorem 3.1, Theorem 3.2]{kaur}. 
Now, from the proof of \cite[Lemma 4.2]{hatuip5}, it follows that the group
\begin{eqnarray*}
	G^* &=&\langle \alpha, \alpha_j, \beta_i
	\mid [\alpha_i,\alpha] = \alpha_{i+1},\ 
	[\alpha_1,\alpha_2] = \alpha_4 \beta_3,\ 
	[\alpha_4,\alpha] = \beta_1,\ [\alpha_1,\alpha_4] = \\ 
	&& [\alpha_3,\alpha_2] = \beta_2,\  [\alpha_1,\alpha_3] = \beta_1 \beta_2,\ \alpha^p = \alpha_{j}^p = \beta_i^p = 1,\ 1 \le i \le 3,\ 1 \le j \le 4 \rangle
\end{eqnarray*} 
is a representation group of $G$ of order $p^8$ and $G^*/Z({G}^*) = G$, where $Z({G}^*) = \langle \beta_1, \beta_2, \beta_3 \rangle$.  \\
Consider the abelian normal subgroup
$N = \langle \beta_1, \beta_2, \beta_3, \alpha_3, \alpha_4 \rangle$ of ${G}^*$ of order $p^5$. 
Let $\chi' \in \operatorname{Irr}(Z({G}^*))$ such that $\chi'(\beta_1) = \xi^{r}$, $\chi'(\beta_2) = \xi^{s}$, $\chi'(\beta_3)= \xi^q$, where $\xi$ is a primitive $p$th root of unity and $0 \le r, s, q \le (p-1)$ such that $r$, $s$, $q$ are  not all $0$.
We use Clifford's theory here to study $\operatorname{Irr}(G^* \mid \chi')$. 
Let $\chi \in \operatorname{Irr}(N)$ such that
$\chi|_{Z({G}^*)} = \chi'$ and $\chi(\alpha_4) = \xi^t$ where $0 < t \le (p-1)$. Now $g \in {G}^*$ can be written as $g = \alpha^j \alpha_1^k \alpha_2^l n'$ for some $n' \in N$ and any element of $N$ is of the form $n_1 = \alpha_3^m \alpha_4^n z$ for some $z \in Z({G}^*)$.
Therefore,
\[
\chi^{\alpha^j \alpha_1^k \alpha_2^l}(\alpha_3^m \alpha_4^n  z)
= \chi(\alpha_3^m \alpha_4^n z)
\iff
\chi(\alpha^j \alpha_1^k \alpha_2^l \alpha_3^m \alpha_4^n \alpha_2^{-l} \alpha_1^{-k} \alpha^{-j})
= \chi(\alpha_3^m \alpha_4^n).\]
Upon simplification,  we obtain that
$\alpha^j \alpha_1^k \alpha_2^l n' \in I_{{G}^*}(\chi)$
if and only if
\[
\chi(\beta_1^{km-jn+m\binom{p-j}{2}} \beta_2^{-lm+kn+km} \alpha_4^{-jm}) = 1
\text{ i.e., } 
\xi^{\,r(km-jn+m\binom{p-j}{2})+s(-lm+kn+km)-tjm} = 1.
\] Let $\Gamma \in \operatorname{Irr}(G^* \mid \chi)$. Now the following cases occur.

\textbf{Case (a)} Suppose $s\neq 0$. Then $I_{{G}^*}(\chi) = \langle N, \alpha\alpha_1^{s^{-1}r}\alpha_2^{-s^{-1}t+(s^{-1}r)^2+2s^{-1}r} \rangle$ and $\deg(\Gamma)=p^2$.

\textbf{Case (b)} Suppose $r\neq 0$ and $s =0$. Then $I_{{G}^*}(\chi) = \langle N, \alpha_2 \rangle$ and $\deg(\Gamma)=p^2$. 

\textbf{Case (c)} Suppose $r = 0$, $s= 0$ and $t \neq 0$. Then $I_{{G}^*}(\chi) = \langle N, \alpha_1, \alpha_2 \rangle$ is of order $p^7$. Since $\chi$ extends to $N_1= \langle N, \alpha_1  \rangle$, let $\tilde{\chi}: N_1 \to \C^\times$ be such that $\tilde{\chi}|_N = \chi$. Now,  $\tilde{\chi}^{(\alpha_2)}(\alpha_1)= \tilde{\chi}(\alpha_1)\tilde{\chi}(\alpha_4 \beta_3)^{-1}$. Therefore, $\chi$ extends to $I_{{G}^*}(\chi)$ if and only if $\chi(\alpha_4)= \chi(\beta_3)^{-1}$. 
Hence, $\deg(\Gamma) =p$ or $p^2$, and $\deg(\Gamma) =p$ if and only if $\chi$ extends to $I_{{G}^*}(\chi)$.
%if and only if $\deg(\Gamma)=p^2$ if $\chi$ does not extend to $I_{{G}^*}(\chi)$. 
If $\deg(\Gamma)= p$, choosing $l_i= \alpha_3^i$, $i=1,2,..,p$, as left coset representatives of $I_{{G}^*}(\chi)$ in $G^*$, we get $\Gamma(\alpha_4)= \xi^{-q} I_{p \times p}$.

Now, suppose $\rho$ is a faithful $\alpha$-representation.
If $\rho$ is an ordinary representation, then  $\deg(\rho) \geq \delta(G) =p^2$. Otherwise, we have $\alpha= \operatorname{tra}(\chi')$ with $\chi'\neq 1$.
Let $\rho =\oplus_{i=1}^l\rho_i$, where $\rho_i \in \irr^\alpha(G)$ such that $\deg(\rho_i)=p$. Then, by Theorem \ref{restriction}, $\rho|_{Z(G)}$ is faithful. Let $\Gamma_i$ correspond to $\rho_i$, where the correspondence is given by Theorem \ref{bijective correspondence}. If $\deg(\Gamma_i) = p$, then 
by the above discussion,
$\Gamma_i(\alpha_4)=\xi^{-q} I_{p \times p}$ and in this case, $\rho|_{Z(G)}$ cannot be faithful, and so $\rho$ is not faithful. It says that $\rho$ must be an irreducible $\alpha$-representation.
Thus, $\tg= \ti =p^2$.


%Let $\psi$ be a faithful $\alpha$-representation of $G$ of degree $\tau(G)$. If $\alpha$ is a non-trivial cocycle, then $\alpha=\tra(\chi')$, where $\chi'$ corresponds to one of the cases (a)-(c). From Theorem \ref{bijective correspondence} and Case (c), it follows that $\deg(\psi) \geq p^2$. If $\alpha=1$, then $\psi$ is an ordinary representation and $\deg(\psi) \geq \delta(G) =p^2$. Hence, the result follows. 

\iffalse
Since $\tra: \Hom(Z({G}^*),\C^\times)\to H^2(G,\C^\times)$ is an isomorphism, for any non-trivial cocycle $\alpha$ of $G$, there is a non-trivial $\chi' \in \mathrm{Hom}(A_3, \C^\times)$ such that $\tra(\chi')=[\alpha]$. Let $\chi' \in \operatorname{Irr}(Z({G}^*))$ such that $\chi'(\beta_1) = 1$, $\chi'(\beta_2) = 1$, $\chi'(\beta_3)= \xi$ where $\xi$ is a primitive $p$th root of unity. Suppose $\rho_i \in \mathrm{Irr}^\alpha(G)$ correspond to $\Gamma_i \in \operatorname{Irr}({G}^* \mid \chi')$, where the correspondence is given by Theorem \ref{bijective correspondence}. 

Our aim is to find $\rho_i$ with deg$(\rho_i)=p$ such that $\rho_1 \oplus \rho_2$ is faithful on $G$ which is equivalent to find proper $\Gamma_i \in \operatorname{Irr}(G^* \mid \chi')$ which gives such $\rho_i$, $i=1,2$. We use Clifford's theory here. Let $\chi \in \operatorname{Irr}(N)$ such that
$\chi|_{Z({G}^*)} = \chi'$ and $\chi(\alpha_4) = \xi^t$ where  $0 < t \le (p-1)$. Now $g \in {G}^*$ can be written as
$g = \alpha^j \alpha_1^k \alpha_2^l n'$
for some $n' \in N$ and any element of $N$ is of the form $n_1 = \alpha_3^m \alpha_4^n z$
for some $z \in Z({G}^*)$.
Therefore,
$$\chi^{(\alpha^j \alpha_1^k \alpha_2^l)}(\alpha_3^m \alpha_4^n  z)
= \chi(\alpha_3^m \alpha_4^n z)
\iff
\chi(\alpha^j \alpha_1^k \alpha_2^l \alpha_3^m \alpha_4^n \alpha_2^{-l} \alpha_1^{-k} \alpha^{-j})
= \chi(\alpha_3^m \alpha_4^n).$$

Upon simplification, we obtain that
$\alpha^j \alpha_1^k \alpha_2^l n' \in I_{{G}^*}(\chi)$
if and only if
\[
\chi(\beta_1^{km-jn-m\binom{-j}{2}} \beta_2^{lm+kn+km} \alpha_4^{jm}) = 1 \text{ i.e., }\xi^{tjm} = 1.
\]

Taking $m=1$, we have $\xi^{tj} = 1$ implies $j=0$. Hence, $N_1= I_{{G}^*}(\chi) = \langle N, \alpha_1, \alpha_2 \rangle$ is of order $p^7$. Let $\chi_1$ and $\chi_2 \in \operatorname{Irr}(N)$ such that $\chi_1(\alpha_4)=\xi$, $\chi_2(\alpha_4)=\xi^2$ and $\chi_i|_{Z({G}^*)} = \chi'$. Let $\Gamma_i = \operatorname{Ind}^{G^*}_{I_{G^*}(\chi_i)}(\tilde{\chi}_i)$, where $\tilde{\chi}_i:I_{G^*}(\chi_i)\to \C^\times$ is an extension of $\chi_i$. Choosing $l_i=\alpha_3^i$, $i=1,2,..,p$, as left coset representatives of $N_1$ in $G^*$, we get $\Gamma_1(\alpha_4)= \xi I_{p \times p}$ and $\Gamma_2(\alpha_4)= \xi^2 I_{p \times p}$. Let $\alpha= \operatorname{tra}(\chi')$ and $\rho_i \in \mathrm{Irr}^\alpha(G)$ correspond to $\Gamma_i$, where the correspondence is given by Theorem \ref{bijective correspondence}. Let $\rho= \rho_1 \oplus \rho_2$. 

Since $\rho|_{Z(G)}$ is faithful, it follows from Theorem \ref{restriction} that $\rho$ is faithful. Hence, $\tau(G) \leq 2p$. \\
Proceeding along the same lines as the proof of $(xv)$, we get $\tau(G)= 2p$.  \\ 
\fi


%Using Clifford's theory, we can check that if $\Gamma \in \operatorname{Irr}({G}_1^* \mid \chi')$ is of degree $p$, then $|{G}_1^*/I_{{G}_1^*}(\chi)|=p$. 

%Let $\tau(G)=m$ and $\phi$ be a faithful projective representation of $G$ of degree $m$. Since $\delta(G) =p^2$, if $\phi$ is an ordinary representation, then $\operatorname{deg}(\phi)\geq p^2$. Therefore, $\phi$ has to be a true projective representation and $\tau(G)=p$ or $2p$. Since $\tau_{irr}(G)>p$, the result follows by Remark \ref{some results on p-groups}(iii). 

%We propose the following open question. \textbf{Open question:} If $G$ and $H$ are $p$-groups and $G \sim_\mathbb C H$ (See \cite[Definition 1.1]{sabnam}, then $G$ is capable if and only if $H$ is capable.

% By \cite[Table 4.1]{james}, $G= \Gamma/Z(\Gamma)$, where $\Gamma \in$ $\Phi_{38}$. By \cite[Table 4.1]{james} and [\cite{kaur}, Theorem 3.3], $\delta_{irr}(\Gamma)=p^2$. Hence, by Remark \ref{irreducible projective embedding degree for capable groups}, $\tau_{irr}(G)$ exists and satisfies $\tau_{irr}(G)\leq p^2$. Since $\Phi_{2}(1^4) \subseteq G$, it follows from Theorem \ref{$p^4$} that $\tau(G)\geq p+2$. This implies that $\tau_{irr}(G)$ must be equal to $p^2$. \\

\end{proof}

We summarize our results on $\tau(G)$ and $\tau_{irr}(G)$, given in this section, in the following tables. If $\tau_{irr}(G)$ does not exist for a group $G$, then we denote this by writing $\tau_{irr}(G)=$ -.

\begin{table}[h!]
\centering
\small
\setlength{\extrarowheight}{2pt}
\begin{tabular}{|c|c|c|}
	\hline
	$G$ & $\tau(G)$ & $\tau_{irr}(G)$ \\ 
	\hline
	$\Phi_2(21)$ & $p+1$ & - \\ 
	$\Phi_2(1^3)$ & $p$ & $p$ \\ \
	$Q_8$ & 3 & - \\ 
	$D_8$ & 2 & 2 \\ 
	\hline
\end{tabular}
\caption{$\tau(G)$ and $\tau_{irr}(G)$ for non-abelian groups of order $p^3$}
\end{table}

\begin{table}[h!]
\centering
\small
\begin{tabular}{|c|c|c|c|}
	\hline
	Group ID & Group & $\tau(G)$ & $\tau_{irr}(G)$ \\
	\hline
	3  & $(\mathbb{Z}/2\mathbb{Z})^{2} \rtimes \mathbb{Z}/4\mathbb{Z}$ & 4 & 4 \\
	4  & $(\mathbb{Z}/4\mathbb{Z}) \rtimes (\mathbb{Z}/4\mathbb{Z})$ & 4 & - \\
	6  & $\mathbb{Z}/8\mathbb{Z} \rtimes_{5} \mathbb{Z}/2\mathbb{Z}$ & 3 & - \\
	7  & $D_{16}$ & 2 & 2 \\
	8  & $\mathbb{Z}/8\mathbb{Z} \rtimes_{3} \mathbb{Z}/2\mathbb{Z}$ & 3 & - \\
	9  & $Q_{16}$ & 2 & - \\
	11 & $D_{8} \times \mathbb{Z}/2\mathbb{Z}$ & 4 & 4 \\
	12 & $Q_{8} \times \mathbb{Z}/2\mathbb{Z}$ & 4 & - \\
	13 & $Q_{8} \rtimes \mathbb{Z}/2\mathbb{Z}$ & 3 & - \\
	\hline
\end{tabular}
\caption{$\tau(G)$ and $\tau_{irr}(G)$ for non-abelian groups of order $2^4$}
\end{table}

\begin{table}[!hbt]
\centering
\small
\begin{tabular}{|c|c|c|}
	\hline
	Group & $\tau(G)$ & $\tau_{irr}(G)$ \\
	\hline
	$\Phi_{2}(211)a$ & $p+2$ & - \\
	$\Phi_{2}(1^{4})$ & $p+2$ & $p^2$ \\
	$\Phi_{2}(31)$ & $p+1$ & - \\
	$\Phi_{2}(22)$ & $p+2$ & $p^2$ \\
	$\Phi_{2}(211)b$ & $p+1$ & -\\
	$\Phi_{2}(211)c$ & $p+2$ & - \\
	$\Phi_{3}(211)a$ & $p+1$ & -\\
	$\Phi_{3}(211)b_{r}$ & $p+1$ & -\\
	$\Phi_{3}(1^{4})$ & $p$ & $p$ \\
	\hline
\end{tabular}
\caption{$\tau(G)$ and $\tau_{irr}(G)$ for non-abelian groups of order $p^4$, $p \geq 3$} %for which $\tau_{irr}(G)$ does not exist}
\end{table}

\begin{table}[H]
\centering
\small
\begin{tabular}{|c|c|c||c|c|c|}
\hline
$G$ & $\tau(G)$ & $\tau_{irr}(G)$ & $G$ & $\tau(G)$ & $\tau_{irr}(G)$  \\ 
\hline
$\Phi_2(311)a$ & $p+2$ & - & $\Phi_4(221)e$ & $2p$ & - \\ 
$\Phi_2(221)a$ & $p+3$ & - & $\Phi_4(221)f_0$ & $2p$ & - \\ 
$\Phi_2(221)b$ & $p+2$ & - & $\Phi_4(221)f_r$ & $2p$ & - \\ 
$\Phi_2(2111)a$ & $p+3$ & - & $\Phi_4(2111)a$ & $2p$ & - \\ 
$\Phi_2(2111)b$ & $p+2$ & - & $\Phi_4(2111)b$ & $2p$ & - \\ 
$\Phi_2(2111)c$ & $p+3$ & - & $\Phi_4(2111)c$ & $2p$ & - \\ 
$\Phi_2(2111)d$ & $p+2$ & - & $\Phi_4(1^5)$ & $2p$ & $p^2$ \\ 
$\Phi_2(1^5)$ & $p+3$ & $p^2$ & $\Phi_5(2111)$ & $2p$ & - \\ 
$\Phi_2(41)$ & $p+1$ & - & $\Phi_5(1^5)$ & $2p$ & - \\ 
$\Phi_2(32)a_1$ & $p+2$ & - & $\Phi_6(221)a$ & $2p+1$ & - \\ 
$\Phi_2(32)a_2$ & $p+2$ & - & $\Phi_6(221)b_r, \, r \neq \tfrac{1}{2}(p-1)$ & $2p+1$ & - \\ 
$\Phi_2(311)b$ & $p+1$ & - & $\Phi_6(221)b_{\tfrac{1}{2}(p-1)}$ & $2p+1$ & $p^2$ \\ 
$\Phi_2(311)c$ & $p+2$ & - & $\Phi_6(221)c_r$ & $2p+1$ & - \\ 
$\Phi_2(221)c$ & $p+2$ & $p^2$ & $\Phi_6(221)d_0$ & $2p+1$ & $p^2$ \\ 
$\Phi_2(221)d$ & $p+3$ & - & $\Phi_6(221)d_r$ & $2p+1$ & - \\ 
$\Phi_3(2111)a$ & $p+2$ & - & $\Phi_6(2111)a$ & $2p+1$ & - \\ 
$\Phi_3(2111)b_r$ & $p+2$ & - & $\Phi_6(2111)b_1$ & $2p+1$ & - \\ 
$\Phi_3(1^5)$ & $p+2$ & $p^2$ & $\Phi_6(2111)b_\nu$ & $2p+1$ & - \\ 
$\Phi_3(311)a$ & $p+1$ & - &  $\Phi_6(1^5)$ & $2p$  & $p^2$ \\ 
$\Phi_3(311)b_r$ & $p+1$ & - & $\Phi_7(2111)a$ & $2p$ & - \\ 
$\Phi_3(221)a$ & $p+2$ & - & $\Phi_7(2111)b_r$ & $2p$ & - \\ 
$\Phi_3(221)b_r$ & $p+2$ & $p^2$ & $\Phi_7(2111)c$ & $2p$ & - \\ 
$\Phi_3(2111)c$ & $p+1$ & - & $\Phi_7(1^5)$ & $2p$ & $p^2$ \\ 
$\Phi_3(2111)d$ & $p+2$ & - & $\Phi_8(32)$ & $p^2+1$ & - \\ 
$\Phi_3(2111)e$ & $p+2$ & - & $\Phi_9(2111)a$ & $p+1$ & - \\ 
$\Phi_4(221)a$ & $2p$ & - & $\Phi_9(2111)b_r$ & $p+1$ & - \\ 
$\Phi_4(221)b$ & $2p$ & $p^2$ & $\Phi_9(1^5)$ & $p$ & $p$ \\ 
$\Phi_4(221)c$ & $2p$ & - & $\Phi_{10}(2111)a_r$ & $p^2+1$ & - \\ 
$\Phi_4(221)d_r, \, r \neq \tfrac{1}{2}(p-1)$ & $2p$ & - & $\Phi_{10}(2111)b_r$ & $p^2+1$ & -  \\ 
$\Phi_4(221)d_{\tfrac{1}{2}(p-1)}$ & $2p$ & - &  $\Phi_{10}(1^5)$ & $p^2$ & $p^2$ \\ 
\hline
\end{tabular}
\caption{$\tau(G)$ and $\tau_{irr}(G)$ for non-abelian groups of order $p^5$, $p \geq 5$.}
\end{table}

%\poonam{Should I change $\Phi_6(2111)b_1$ and $\Phi_6(2111)b_\nu$ to $\Phi_6(2111)b_r$?}
\iffalse
\poonam{The title page should contain: If available, the 16-digit ORCID of the author(s), The affiliation(s) of the author(s), i.e. institution, (department), city, (state), country. A clear indication and an active e-mail address of the corresponding author}

\poonam{The list of references should only include works that are cited in the text and that have been published or accepted for publication. Personal communications and unpublished works should only be mentioned in the text.}

\poonam{Authors wishing to include figures, tables, or text passages that have already been published elsewhere are required to obtain permission from the copyright owner(s) for both the print and online format and to include evidence that such permission has been granted when submitting their papers. Any material received without such evidence will be assumed to originate from the authors.}

\poonam{Please ensure you provide all relevant editable source files. Failing to submit these source files might cause unnecessary delays in the review and production process.}

\poonam{If available, please always include DOIs as full DOI links in your reference list (e.g. “https://doi.org/abc”).}

\poonam{All tables are to be numbered using Arabic numerals.}

\poonam{Acknowledgements of people, grants, funds, etc. should be placed in a separate section on the title page.}

\poonam{The above should be summarized in a statement and placed in a ‘Declarations’ section before the reference list under a heading of ‘Funding’ and/or ‘Competing interests’.}

\poonam{The authors have no relevant financial or non-financial interests to disclose. No funding was received for conducting this study. }

\poonam{The following statements should be included under the heading "Statements and Declarations" for inclusion in the published paper.}

\poonam{Each figure should have a concise caption describing accurately what the figure depicts. Include the captions in the text file of the manuscript, not in the figure file. Figure captions begin with the term Fig. in bold type, followed by the figure number, also in bold type. No punctuation is to be included after the number, nor is any punctuation to be placed at the end of the caption.}

\poonam{Identify previously published material by giving the original source in the form of a reference citation at the end of the figure caption.}

\poonam{If you include figures that have already been published elsewhere, you must obtain permission from the copyright owner(s) for both the print and online format.}

\poonam{Check if any Theorem is referred to twice in a result.}

\poonam{Always use the standard abbreviation of a journal’s name according to the ISSN List of Title Word Abbreviations, see

ISSN.org LTWA If you are unsure, please use the full journal title.}
\Nayak{How to include DOI in reference?}
\Nayak{How to cite Conrad, GAP, Dolgachev?}
\Nayak{In some references, journals have full names while in others, they are abbreviated}
\poonam{Representation groups written twice in keywords}
\iffalse
\begin{table}[h!]
\centering
\[
\begin{array}{|c|}
\hline
p \\
\hline
\begin{array}{|c|c|c|}
	\hline
	\text{Ring} & xy & yx \\
	\hline
	&  &  \\
	&  &  \\
	&  &  \\
	\hline
\end{array}
\\
\hline
ab \\
\hline
\begin{array}{|c|c|c|}
	\hline
	\text{Ring} & xy & yx \\
	\hline
	&  &  \\
	&  &  \\
	&  &  \\
	\hline
\end{array}
\\
\hline
\end{array}
\]
\caption{r}
\end{table}
\fi
\fi


\section*{Acknowledgements}
%The authors are thankful to the editor for providing valuable suggestions which has greatly improved the presentation of the paper. We thank Prof. Eamonn O’Brien for pointing out a representation group of the group $\Phi_2(3^2)a_2$. We are also grateful to the anonymous referee for carefully and critically reading the manuscript and providing many valuable comments. 
The authors acknowledge the support of the National Institute of Science Education and Research (NISER), Bhubaneswar and Homi Bhabha National Institute (HBNI), Mumbai.

%\poonam{The paper by Gurleen mam that we are using is not available online; only the old one is available, which I have cited.}
%\section*{Data Availability}
%No data was used for the research described in the article.

%\section*{Competing interests}
%The authors declare that they have no financial, professional, or personal conflicts of interest that could have affected the research or its interpretation.

\bibliographystyle{plain}
\bibliography{ref}

\end{document}
